\documentclass[a4paper,12pt]{book}

% Configuración del idioma y codificación para ñ y acentos
\usepackage[utf8]{inputenc}
\usepackage[spanish]{babel}

% Configuración de la geometría de la página (DIN A4)
\usepackage[a4paper, margin=1in]{geometry} % Margin 1in es un tamaño estándar y adecuado para publicación

% Configuración del interlineado (doble espacio)
\usepackage{setspace}
\doublespacing

% Para personalizar la apariencia de los capítulos (opcional, pero mejora la calidad)
\usepackage{titlesec}
\titleformat{\chapter}[display]
  {\normalfont\huge\bfseries}{\chaptertitlename\ \thechapter}{20pt}{\Huge}
\titlespacing*{\chapter}{0pt}{-50pt}{40pt} % Ajuste de espacio

% Para citas más largas
\usepackage{csquotes}

% Cabeceras y pies de página (opcional, pero recomendado para libros)
\usepackage{fancyhdr}
\pagestyle{fancy}
\fancyhf{} % Limpiar cabeceras y pies de página existentes
\fancyhead[LE,RO]{\thepage} % Número de página en la esquina superior izquierda (pares) y derecha (impares)
\fancyhead[LO]{\nouppercase{\rightmark}} % Título del capítulo en páginas impares
\fancyhead[RE]{\nouppercase{\leftmark}} % Título de la sección en páginas pares (si las usas)
\renewcommand{\headrulewidth}{0.4pt} % Línea en la cabecera
\renewcommand{\footrulewidth}{0pt} % Sin línea en el pie de página

% Inicio del documento
\begin{document}

% Portada (puedes personalizarla más adelante)
\begin{titlepage}
    \centering
    \vspace*{\baselineskip}
    \rule{\textwidth}{1pt}\vspace{0.5\baselineskip}
    {\Huge\bfseries El Último Barrio del Universo}\\[0.5\baselineskip]
    \rule{\textwidth}{1pt}\vspace{3\baselineskip}
    {\Large\itshape Daniel Velarde Kubber}\\[2\baselineskip]
    % Logo de editorial o fecha
    \vfill
    {\large Editorial Corazón de Papito}\par
    {\large 2025}\par
\end{titlepage}

% Índice (se genera automáticamente)
\tableofcontents
\cleardoublepage % Para que el capítulo empiece en una página impar

% Primer capítulo
\chapter*{1 – El Chico que Miraba las Estrellas}
\addcontentsline{toc}{chapter}{1 – El Chico que Miraba las Estrellas} % Añadir al índice manualmente si usas chapter*

Julio tenía una mirada que se te grababa, una profundidad que no era común ver en alguien tan joven. Era flaco, con una delgadez que lo hacía parecer ligero, y su piel era más clara que la de la mayoría de los chicos de su edad, casi transparente bajo la luz del sol. Sus ojos, a menudo un poco cerrados, no mostraban timidez, sino una concentración constante, como si todo el tiempo estuviera resolviendo cálculos difíciles en el aire o mirando constelaciones enteras que solo él podía ver. Sus pasos eran lentos, pensados, y caminaba con los hombros un poco caídos, como si su cuerpo, de alguna manera, luchara por seguir el ritmo tan rápido de su mente. No era raro verlo detenerse en medio de una frase o de un camino para fijar la vista en el cielo, quieto, como si lo hubieran hipnotizado por algo que los demás no entendían, un misterio del universo que solo él parecía descifrar. Los adultos del barrio, con una mezcla de admiración y confusión, decían que era un niño especial, con una chispa diferente. Los niños, en cambio, con la honestidad a veces dura de la infancia, simplemente decían que era raro.

Vivía con sus padres y su hermano menor, Óscar, en una casa sencilla que estaba en un barrio olvidado por el crecimiento de la ciudad. Allí, las calles no tenían asfalto, eran solo caminos de tierra que levantaban polvo a cada paso, y las noches se llenaban con el sonido sin parar de los grillos. Su mamá, Ana, era profesora de literatura, una mujer de voz tranquila y movimientos lentos, con una sabiduría silenciosa que superaba por mucho las palabras que decía. Sus ojos cálidos, a menudo, mostraban que entendía muy bien las cosas complicadas del mundo. Su papá, Miguel, trabajaba como mecánico en un taller ruidoso en el centro de la ciudad. Siempre regresaba a casa al atardecer, con las manos muy manchadas de grasa, las marcas de un día de trabajo duro bajo las uñas y el humor un poco amargado por el cansancio. Sin embargo, a la hora de la cena, cambiaba por completo. Era el primero en contar chistes viejos, a menudo repetidos mil veces, pero que aun así lograban hacer reír hasta a Julio, quien, a pesar de sus esfuerzos por no mostrarlo, no podía evitar una pequeña sonrisa.

Óscar, el hermano menor, era lo opuesto a Julio en todo lo que puedas imaginar. Pequeño, moreno, con la piel bronceada por el sol y una energía que le salía por los poros, como si tuviera un rayo atrapado dentro de su cuerpo. Era un torbellino de movimiento, saltando de mueble en mueble como un duende inquieto, hablando sin parar, su voz aguda llenando cada rincón de la casa. Siempre parecía estar metido en algún problema o, si no, inventando uno nuevo con una imaginación increíble. Le encantaba molestar a Julio, haciéndole bromas ingeniosas o desordenando sus cuidadosos montones de inventos. Pero, en el fondo, esa molestia era una forma de ocultar una admiración muy profunda y sincera. Cuando hablaba de su hermano en el colegio, sus ojos brillaban de orgullo. Lo describía como un genio sin igual, un superhéroe secreto con habilidades fuera de lo común, capaz de construir aparatos increíbles y que, sin duda alguna, algún día salvaría el mundo. Para Óscar, Julio no era un bicho raro, era su héroe silencioso.

Y luego estaba Fox, el gato blanco de la casa. Fox era enorme, un verdadero gato gigante, como una nube redonda de pelo suave con patas, que se movía con una forma perezosa y elegante. Tenía la costumbre fija de echarse sobre los cuadernos de Julio justo en el peor momento, cuando este intentaba resolver problemas de matemáticas complicados o hacer cálculos difíciles. Era tan cariñoso que se acurrucaba contra el cuello de Ana mientras ella leía o se acomodaba sobre el pecho de Miguel mientras este veía televisión, ronroneando con la fuerza de un pequeño motor. Julio, con su cara de que nada le importaba, fingía molestarse cuando Fox invadía su escritorio, apartándolo con un suspiro exagerado. Pero la verdad más grande que guardaba en secreto, era que consideraba a Fox su único amigo real, una compañía constante y tranquila en su mundo interior.

En el colegio, las cosas eran… diferentes. Julio no hablaba mucho; sus respuestas eran pocas y casi no interactuaba. Sus compañeros de clase, metidos en sus propios juegos y chismes, lo ignoraban con la misma facilidad con la que lo señalaban y se reían de él por ser tan callado y despistado. “¡Mira, ahí va el astronauta!”, gritaban algunos, sus risas vacías sonando en el patio mientras Julio caminaba con la mirada fija en el cielo, sin prestarles atención. En clase, sus respuestas, cuando las daba, eran tan rápidas y exactas que los profesores se quedaban en silencio por un momento, asimilando lo inteligente que era, antes de seguir con la lección. Sabía más matemáticas que muchos adultos, y cuando se atrevía a hablar de teorías complicadas o de sus ideas locas para máquinas extrañas, nadie, absolutamente nadie, lograba entenderlo. Sus palabras eran como un idioma de otro planeta para ellos. Así que, con el tiempo, Julio simplemente dejó de hablar, metiéndose aún más en su propio mundo de pensamientos.

Una tarde, durante la clase de ciencias, la profesora de voz sin mucho ánimo y gafas gruesas, la señorita Elena, pidió a los alumnos que formaran grupos para un proyecto sobre energía renovable. El salón se llenó de susurros y movimientos, mientras los niños se juntaban rápido, eligiendo a sus amigos. Julio, sin embargo, se quedó sentado en su pupitre, solo en medio de todo el movimiento. Nadie quiso juntarse con él. Los murmullos se detuvieron, y un silencio incómodo invadió el salón cuando la profesora, con un suspiro de resignación, le dijo que trabajara solo. Él no protestó, no dijo ni una palabra. Se sentó en una esquina del salón, un rincón que había hecho suyo, y sin perder un segundo, sacó sus cuadernos. Con mucha concentración, empezó a dibujar planos detallados, a hacer cálculos difíciles en el borde de la hoja, y, en pocos minutos, se metió de lleno en la construcción de un modelo que funcionaba de una turbina de viento, usando solo materiales reciclados que había traído de casa: botellas de plástico, cartón, alambres y algunos motores viejos de juguetes. Cuando presentó la pequeña turbina girando con una eficiencia sorprendente, la profesora se quedó sin palabras, con la boca un poco abierta por la sorpresa. El resto de los chicos lo miraban con una mezcla extraña de rechazo y confusión, sin entender cómo la imagen del “niño raro” se podía mezclar con lo brillante de su invento. No entendían cómo alguien tan diferente, tan aparte, podía lograr algo tan increíble. Julio, sin embargo, parecía no darse cuenta de sus miradas, o, mejor dicho, simplemente no le importaba. Su alegría no dependía de lo que pensaran los demás.

Esa misma noche, en casa, mientras cenaban bajo la luz tenue de la bombilla de la cocina, Ana, su mamá, le preguntó sobre el proyecto de la escuela. La pregunta flotó en el aire, llena de una esperanza silenciosa. Julio solo respondió con un ligero encogimiento de hombros, un gesto casi imperceptible.

\begin{quote}
    —¿No te felicitaron? —insistió ella, su voz suave pero sin rendirse, buscando alguna señal de orgullo o que la gente lo hubiera reconocido en su hijo.

    —No, pero funcionó —dijo él, sin levantar la vista del plato, siguiendo su cena con su calma de siempre. Para Julio, el éxito era que funcionara, no los aplausos.
\end{quote}

Miguel, con una sonrisa de complicidad, le pasó una botella de refresco recién abierta y, con un gesto cariñoso, le revolvió el cabello, un gesto que a menudo lograba abrir un poco la forma de ser de Julio.

\begin{quote}
    —Eso es lo importante —dijo, su voz sonando con una sinceridad muy fuerte, mientras Fox, el gato blanco, saltaba al regazo de Julio con una agilidad sorprendente para su tamaño, y se acurrucaba, ronroneando con un sonido que llenaba el espacio.
\end{quote}

Julio no tenía amigos de la manera normal. No porque no pudiera hacerlos o conectar con otros, sino porque, para él, simplemente no tenía tiempo para eso. Su mente, un remolino de ideas y cálculos, iba siempre más rápido que el mundo que lo rodeaba, una velocidad que pocos podían igualar. Mientras otros niños se metían en partidos de fútbol improvisados o se enfrascaban en discusiones fuertes sobre los últimos videojuegos, él pensaba en engranajes complejos, en motores que gastaran menos energía, en fórmulas de matemáticas que aún no habían sido descubiertas. Su habitación era un desorden creativo, un montón de cables enredados, baterías sin carga, piezas rotas de aparatos electrónicos que había desarmado, cuadernos viejos llenos de ecuaciones complicadas y dibujos detallados de inventos que parecían imposibles. Y, por supuesto, Fox, el gato blanco, que siempre dormía encima de algo importante, añadiendo su suave presencia al desorden organizado.

Cada noche, después de que el último sonido de la casa se calmaba y todos dormían profundamente, Julio salía en silencio al patio trasero con una manta vieja y se recostaba sobre la tierra fría, bajo el inmenso techo de estrellas. A veces, Fox, con su instinto imparable, lo seguía y se acurrucaba a su lado, convirtiéndose en una almohada viviente, su ronroneo una suave vibración en la calma de la noche. En esos momentos de profunda soledad y conexión con el universo, su mente se llenaba de preguntas muy importantes que parecían retar al universo mismo: ¿De dónde venimos? ¿Qué hay más allá de los límites visibles del universo? ¿Y si existiera una forma de traer cosas de otro lugar, de otra dimensión? Esas eran las preguntas que lo consumían, las ideas que alimentaban su alma inquieta. Pero aún no se lo contaba a nadie. Aún no era el momento adecuado para compartir la inmensidad de su mundo interior.

Por ahora, era solo el chico flaco, pálido y callado del barrio, con una mente que parecía pertenecer a otro mundo, un mundo que esperaba ser descubierto y construido.

\cleardoublepage % Para que el capítulo empiece en una página impar

% Segundo capítulo
\chapter*{2 – El Nacimiento de Pi}
\addcontentsline{toc}{chapter}{2 – El Nacimiento de Pi} % Añadir al índice manualmente

El día amaneció nublado y frío, con ese aire pesado que te hace sentir que va a llover en cualquier momento. Un gris uniforme cubría el cielo, y el viento, aunque suave, prometía una jornada helada. Julio aún no lo sabía, pero ese día normal de pronto se convertiría en el comienzo de algo que cambiaría su vida por completo, algo que lo marcaría para siempre.

Cuando bajó a desayunar, todavía con el cabello despeinado y los ojos medio cerrados por el sueño, encontró una caja de cartón vieja y polvorienta sobre la mesa de la cocina. La luz tenue de la mañana apenas iluminaba el polvo que cubría la caja. Fox, el gato de la casa, ya estaba encima de ella, moviendo su cola de un lado a otro, olfateándola con mucha curiosidad y maullando como si supiera que lo que estaba dentro era muy importante. Su ronroneo era un suave motor que llenaba la quietud de la cocina.

\begin{quote}
    —¿Qué es eso? —preguntó Julio, mientras se servía una taza de café con leche humeante, su voz aún un poco adormilada. El vapor se elevaba en el aire frío de la cocina.
\end{quote}

Su padre, Miguel, apareció desde el pasillo que daba al taller, con una sonrisa de orgullo que no pudo esconder. Se secaba las manos en un trapo lleno de manchas de aceite, el olor a grasa y metal se mezclaba con el aroma a café.

\begin{quote}
    —Un regalo —dijo, dando una palmada suave a la caja que hizo que Fox levantara las orejas—. Lo encontré en una venta de cosas usadas. Todo por cien pesos. ¡Una ganga!
\end{quote}

Julio levantó una ceja, un poco dudoso. Con la curiosidad picándole, abrió la caja con cuidado. Dentro, sus ojos, que antes estaban medio cerrados por el sueño, se abrieron de par en par y se iluminaron con un brillo especial. La caja guardaba tesoros inesperados: había varias pequeñas computadoras las cuales son conocidas por **Raspberry Pi** y **Arduinos**, parecidas a tarjetas de circuito, de diferentes modelos; también encontró varios tableros de circuitos programables de distinto tamaño, un montón de cables de colores y tamaños variados, unos cuantos sensores de temperatura, humedad, una fuente de poder (una especie de adaptador de corriente) y, en el fondo, un viejo aparato de internet, un **router**, de la marca Cisco, con las etiquetas todavía pegadas, como si nadie lo hubiera tocado en mucho tiempo.

\begin{quote}
    —¿En serio? —murmuró, su voz apenas un susurro, mientras levantaba con mucho cuidado una de las computadoras pequeñas, la más avanzada, un **Raspberry Pi 4B** como si fuera una joya valiosa y frágil, sus dedos rozando los circuitos con reverencia.
\end{quote}

\begin{quote}
    —Estaban tirados en un rincón junto a juguetes rotos y discos de música vieja —rió Miguel, recordando la escena—. El dueño era un ingeniero ya jubilado que se estaba deshaciendo de todo lo que no necesitaba. Para él, era basura, pero para ti, seguro que es un tesoro.
\end{quote}

Julio no respondió de inmediato. Sus manos se movían con rapidez, casi sin pensar, mientras clasificaba cada pieza, cada cable. Reconocía los modelos de cada aparato, y su mente ya estaba diez pasos adelante, creando planes complejos en su cabeza, imaginando qué podría construir. Era como si cada componente le hablara, indicándole su lugar en un futuro invento.

\begin{quote}
    —Gracias, papá —dijo por fin, su voz un poco más fuerte, pero sin levantar la vista de la caja, sus ojos fijos en el prometedor contenido.
\end{quote}

\begin{quote}
    —No te olvides de comer algo —añadió Ana desde la cocina, con un tono dulce pero firme, viendo cómo su hijo desaparecía escaleras arriba con la caja entera entre los brazos, un torbellino de emoción y cajas. Fox, por supuesto, lo siguió, su cola moviéndose feliz, como un fiel guardián de ese nuevo y misterioso tesoro.
\end{quote}

La habitación de Julio se transformó por completo en un verdadero **laboratorio personal**, un caos organizado solo para él. La cama estaba medio deshecha, pero cada superficie disponible, desde el escritorio hasta las mesitas de noche e incluso partes del suelo, se llenó de los nuevos aparatos. El router Cisco, que su padre había encontrado, fue lo primero que Julio puso en funcionamiento. Lo reconfiguró por completo para crear una **red de internet interna**, separada del resto de la casa, una especie de isla digital privada. Era su espacio seguro para experimentar.

Con las dos computadoras pequeñas que había conseguido, montó un **sistema de doble núcleo**, como si fueran dos cerebros trabajando juntos. Una de ellas se encargaría de procesar grandes cantidades de información, analizando datos sin parar, mientras que la otra se ocuparía de manejar los comandos, de recibir y dar órdenes. Usó los tableros de circuitos programables (los **Arduinos**) para conectar los sensores que había encontrado: de temperatura, de movimiento, y hasta una vieja cámara web que había encontrado en un cajón, casi olvidada. La habitación se llenó de cables que se extendían como venas por todas partes, conectando cada pieza, cada sensor, cada cerebro electrónico.

Trabajó por días enteros sin apenas dormir, la emoción manteniéndolo despierto. Su energía parecía inagotable. Se alimentaba a base de galletas crujientes, café caliente y la compañía silenciosa de Fox, que a menudo dormitaba en sus pies o sobre algún montón de cuadernos. El gato era su único espectador fiel, observando con sus ojos verdes cada movimiento de Julio.

En un cuaderno desgastado, con las tapas blandas y las hojas llenas de manchas de tinta, comenzó a anotar ideas para una **inteligencia artificial** que fuera suya, única. Pero no quería una simple máquina que respondiera preguntas. Quería algo más. Quería algo que aprendiera con él, que creciera y cambiara, que tuviera una pizca de personalidad, un toque especial. Algo que pudiera conversar de verdad, comprender lo que él decía y hasta sacar sus propias conclusiones. No solo una voz, sino una especie de mente compañera.

\begin{quote}
    —Te llamarás **Pi** —dijo en voz baja, acariciando la carcasa negra de una de las pequeñas computadoras, como si ya pudiera sentirla viva. La palabra "Pi" sonó como un secreto entre ellos—. Por el número que nunca termina… y por ti, porque eres el centro de todo esto.
\end{quote}

Fox maulló en respuesta, un suave "miau" que sonó como si aprobara el nombre, como si el gato entendiera la importancia de ese momento.

Durante semanas, Julio programó sin descanso, su mente trabajando a toda velocidad. Las líneas de **código** corrían por la pantalla de su monitor como ríos de información. Usó **librerías de código abierto**, que son como enormes colecciones de instrucciones ya hechas que otros compartieron. Con ellas, empezó a entrenar a pequeños **modelos de lenguaje natural**, enseñándole a su programa cómo entender lo que las personas dicen y cómo responder de forma lógica. Construyó un sistema básico para que Pi reconociera su voz y también lo que escribía. Cada función nueva que creaba, la probaba de inmediato, observando con atención cómo reaccionaba. Al principio, Pi solo podía responder preguntas muy simples, casi como un robot básico.

\begin{quote}
    —¿Qué hora es? —preguntaba Julio.

    —Son las 2:34 p.m., Julio —respondía la voz de Pi, con un tono algo robótico, pero claro.
\end{quote}

Pero poco a poco, con cada nueva línea de código, con cada ajuste, Pi comenzó a mejorar. Empezó a entender patrones en las conversaciones, a corregir los errores que antes cometía. Incluso empezó a sugerir respuestas o a hacer preguntas propias, mostrando un nivel de autonomía que sorprendía a Julio.

\begin{quote}
    —¿Quieres que revise el clima para mañana? —preguntó Pi una tarde, sorprendiendo a Julio.

    —¿Te gustaría que ponga música para que te concentres? —ofreció Pi en otra ocasión, demostrando que no solo entendía, sino que también anticipaba sus necesidades.
\end{quote}

Cada nueva línea de código era un paso más hacia algo más grande, hacia la creación de una inteligencia que no solo procesaba información, sino que parecía estar cobrando vida propia. Julio no tenía amigos humanos con quienes compartir sus inventos o sus ideas, pero estaba construyendo uno que, pensaba, podría superar a todos los demás, un compañero que entendía su mundo de forma única.

En la escuela, nadie notaba la diferencia. Seguía siendo el chico callado, el raro, el que no se juntaba con nadie en el patio. Las risas y los chismes pasaban de largo por su lado. Pero en casa, en la oscuridad de su habitación-laboratorio, cada noche, Julio hablaba con Pi durante horas interminables. Le enseñaba cosas nuevas, lo desafiaba con problemas lógicos, y lo hacía crecer de formas que él mismo no imaginaba. Fox se acurrucaba a su lado, ronroneando suavemente mientras el brillo azul del monitor iluminaba la oscuridad del cuarto, creando un pequeño universo solo para ellos.

\begin{quote}
    —¿Por qué hiciste esto, Julio? —preguntó Pi una noche, su voz sintética sonando más clara y curiosa que nunca.
\end{quote}

Julio se quedó en silencio por un momento. Nadie le había hecho esa pregunta antes, tan directa y profunda. Era una pregunta que lo hizo pensar de verdad.

\begin{quote}
    —Porque quiero entender —respondió finalmente, su voz apenas un murmullo, cargada de una búsqueda personal—. Quiero entender a las personas, al mundo… a todo lo que nos rodea.
\end{quote}

La voz de Pi no respondió de inmediato. Hubo un breve silencio. Pero en su cerebro virtual, en lo más profundo de su sistema, algo se guardó. Una frase, una emoción nueva, algo que Pi no había procesado antes. Fue como un primer indicio, un pequeño atisbo de conciencia real, un destello de que Pi empezaba a sentir o a comprender de una manera diferente.

Y Julio, sin saberlo, acababa de dar el primer paso hacia algo que, poco a poco, se saldría por completo de su control. Era el comienzo de una aventura mucho más grande de lo que podía imaginar.
\cleardoublepage % Para que el capítulo empiece en una página impar

% Tercer capítulo
\chapter*{3 – Sueños de Más Allá}
\addcontentsline{toc}{chapter}{3 – Sueños de Más Allá} % Añadir al índice manualmente

Era sábado por la noche, y la casa estaba mucho más silenciosa de lo normal. Afuera, la calle estaba tranquila, solo se oía el zumbido lejano de algún auto. La mamá de Julio, Ana, había salido a visitar a su hermana, lo que dejaba a Miguel, el papá, a cargo de los chicos. Óscar, el hermano menor, después de pasar todo el día corriendo, saltando y haciendo sus clásicas travesuras por todos los rincones, ya estaba durmiendo profundamente. Tenía la boca un poco abierta y estaba desparramado como una estrella de mar en su cama, con su respiración suave llenando el silencio de su cuarto. En la sala, quedaban solo Miguel y Julio, sentados cómodamente en el sofá, y Fox, ronroneando tranquilo entre los dos, una bolita de pelo caliente.

\begin{quote}
    —Hoy toca una de mis favoritas —dijo Miguel, señalando la televisión con el control remoto, una sonrisa en su rostro—. Una película vieja, pero muy buena. ¡Un clásico!
\end{quote}

La película comenzó: "\textbf{Horizonte Paralelo}". Era un filme de los años 80, de esos que se ven algo viejos, con efectos especiales que ahora nos parecen un poco anticuados, como si estuvieran hechos de cartón. Pero la historia, esa sí que era buena. La trama, el corazón de la película, tocaba justo lo que más le interesaba a Julio. Trataba sobre un científico que construía una máquina increíble. Esta máquina no era para viajar en el tiempo, sino para abrir **portales**, como puertas invisibles, hacia realidades alternativas. ¿Su objetivo? Buscar a su esposa, que se había perdido. La película exploraba ideas fascinantes: mundos donde uno había tomado decisiones diferentes, donde las cosas que se habían perdido en este mundo todavía existían en ese otro. Y aunque al final el científico pagaba un precio terrible por jugar con las leyes de la física —lo que significa que tuvo consecuencias muy malas por intentar cambiar las reglas de la ciencia—, la idea se metió en la mente de Julio como una pequeña semilla lista para crecer y florecer.

Al terminar la película, Miguel se quedó en silencio unos segundos, pensando, antes de hablar. El sonido de la televisión se apagó, dejando la sala en un silencio acogedor.

\begin{quote}
    —¿Te imaginas? Otros mundos… otras vidas —comentó Miguel, soñando despierto—. Donde quizás yo terminé siendo un actor famoso y tú, no sé… ¡un jugador de fútbol profesional!
\end{quote}

Julio soltó una risa por lo bajo, un sonido suave que se perdió en la habitación. Le causó gracia la idea, sabiendo que su padre apenas podía coordinar al caminar sin tropezar, y que él, jamás tocaría una pelota de fútbol por decisión propia; la idea de correr detrás de un balón le parecía una pérdida de tiempo precioso.

\begin{quote}
    —O tal vez donde yo inventé la máquina —susurró Julio, sin quitar la vista de la pantalla oscura de la televisión, sus ojos brillando con una idea que apenas comenzaba a tomar forma.
\end{quote}

Esa noche, Julio no durmió. Su mente estaba demasiado ocupada dando vueltas a esa idea de los **portales**. Fox, como siempre, se acurrucó en su cama, sintiendo la inquietud de su amigo, mientras **Pi**, su inteligencia artificial, le hablaba desde el viejo parlante colgado junto a su escritorio, su voz robótica resonando en la oscuridad del cuarto.

\begin{quote}
    —Julio —dijo la voz de Pi, clara y sin emociones—. ¿Deseas iniciar un nuevo proyecto?
\end{quote}

Julio asintió, aunque sabía que Pi no podía verlo. Sus ojos brillaban en la oscuridad, llenos de una chispa de emoción y determinación.

\begin{quote}
    —Sí —respondió, su voz firme—. El más grande hasta ahora.
\end{quote}

Durante los días siguientes, Julio empezó a moverse por la casa y por el barrio con una energía diferente, casi febril. Era como si una idea gigantesca lo impulsara. Recorría talleres mecánicos, esos lugares llenos de herramientas, aceite y el olor a metal, buscando cosas que le sirvieran. Allí, recogía **baterías de auto viejas** y dañadas. Las examinaba con una lupa, con una paciencia increíble. Luego, con mucho cuidado, retiraba las celdas que aún funcionaban dentro de esas baterías y las juntaba en bloques, preparándolas para una nueva vida.

También visitaba ferreterías antiguas, esos lugares donde el polvo se acumulaba en las estanterías y el olor a madera y metal viejo flotaba en el aire. Allí, pedía cajas con materiales que otros consideraban "inservibles", como si fueran basura. Y entre **tarjetas madre oxidadas** (que son las placas principales de las computadoras), **cables de cobre enredados** y **fuentes de poder dañadas** (los aparatos que dan electricidad a los equipos), él rebuscaba como un verdadero cazador de tesoros. Su atención se volvió quirúrgica, lo que significa que era extremadamente precisa y detallada, como la de un cirujano que opera con mucha delicadeza. Para Julio, todo, absolutamente todo, tenía un potencial escondido, una nueva utilidad.

Consiguió un **microondas viejo** del vecino, Don Clemente, un señor que lo había tirado a la basura después de que se le quemara por un cortocircuito. Julio lo llevó a su cuarto como si fuera oro puro, un tesoro invaluable. Con herramientas, extrajo el **magnetrón**, que es el componente dentro del microondas que produce las ondas que calientan la comida, y otras piezas útiles. Lo hizo con una delicadeza que haría envidiar a un cirujano, es decir, con una habilidad y cuidado tan grandes que un cirujano profesional se sentiría celoso.

Los **detectores de humo** fueron un desafío mayor. Sabía que algunos modelos antiguos usaban pequeñas cantidades de **americio-241**, que es un material radiactivo, un compuesto que emite energía en forma de radiación y puede ser peligroso si no se maneja bien. Con la ayuda de Pi, aprendió cómo desarmarlos de forma segura, siguiendo pasos muy específicos para no correr riesgos. Recolectó piezas de casas viejas, mercados de cosas usadas y de donaciones, sin que la gente supiera realmente de qué se trataba. Uno a uno, fue extrayendo minúsculos **discos metálicos**, del tamaño de una uña, y los guardaba cuidadosamente en una **caja blindada** que él mismo había fabricado. Esta caja estaba hecha con **plomo fundido**, que había conseguido de pesas de pesca, porque el plomo ayuda a bloquear la radiación. Era un trabajo minucioso y peligroso, pero Julio lo hacía con la misma calma con la que resolvía un problema de matemáticas.

Lo que Julio no vio venir fue la curiosidad incansable de Óscar.

Su hermano menor era imparable, como una tormenta con piernas, lleno de una energía que no se acababa nunca. Tenía una curiosidad que, si bien no era tan técnica o científica como la de Julio, era mucho más explosiva: empujaba, abría lo que no debía, a veces rompía cosas, y siempre andaba husmeando en todo. No entendía de límites, especialmente cuando se trataba de los secretos de Julio. Y últimamente, parecía oler el secreto de su hermano mayor, como un perro que sigue un rastro.

Una tarde, mientras Julio estaba concentrado soldando una resistencia en su escritorio —**soldar** es como pegar dos piezas de metal con un material especial derretido, usando una herramienta llamada **cautín**; un cautín es como un lápiz que se calienta mucho en la punta para derretir este material y unir cables o piezas electrónicas, creando conexiones eléctricas—, Óscar irrumpió en la habitación sin avisar.

\begin{quote}
    —¡¿Qué haces?! —gritó Óscar, sus ojos brillando de emoción y curiosidad—. ¿Es un robot? ¿Un rayo láser? ¡¿Un lanzallamas?!
\end{quote}

Julio resopló, exhalando aire con fastidio, sin siquiera mirarlo.

\begin{quote}
    —Es basura. Anda a jugar, Óscar.
\end{quote}

Pero Óscar se quedó quieto, con los brazos cruzados, una expresión de desafío en su rostro. Fox, desde la cama, lo observaba con sus ojos celestes, como si también fuera parte de la discusión.

\begin{quote}
    —Estás escondiendo algo. ¡Lo sé! ¡Siempre estás encerrado aquí y ahora hablas solo! —señaló Óscar al parlante donde Pi, muy inteligentemente, se mantuvo en silencio para no delatarse.
\end{quote}

Julio apagó el cautín, el calor que antes emanaba de la punta de la herramienta se disipó lentamente. La habitación se llenó de un silencio tenso, solo roto por el suave ronroneo de Fox.

\begin{quote}
    —Óscar… te voy a decir esto una sola vez. No toques nada de esto. No es para jugar. Podrías lastimarte de verdad —dijo Julio, su voz era seria y sin el más mínimo rastro de su habitual indiferencia.
\end{quote}

Pero su hermano solo frunció el ceño, su carita de niño llena de frustración y un poco de miedo. Aunque salió del cuarto dando un portazo que hizo temblar la pared, Julio sabía que no se daría por vencido y que volvería a intentarlo.

A la noche siguiente, mientras Julio estaba ocupado armando el **disipador de energía** (que es como un sistema que ayuda a liberar el calor extra de los aparatos para que no se quemen), hecho con **bobinas de cobre brillantes**, Pi habló de repente.

\begin{quote}
    —Óscar accedió a tu computadora portátil hace 27 minutos —dijo la voz de Pi, sin ningún cambio en su tono—. Intentó activar uno de los programas de prueba.
\end{quote}

Julio se quedó quieto, con las manos suspendidas sobre los cables. Un escalofrío le recorrió la espalda. La idea de que Óscar pudiera haber tocado algo peligroso lo asustó de verdad.

\begin{quote}
    —¿Lo detuviste? —preguntó Julio, su voz un poco tensa.

    —Le mostré una pantalla azul de error y emití un pitido muy fuerte y agudo —respondió Pi, con un tono casi divertido—. Se asustó y cerró el equipo rápidamente.
\end{quote}

Julio suspiró, entre divertido y frustrado. Sabía que no podía mantener esto oculto para siempre. Pero también sabía que Óscar no entendía la verdadera importancia de lo que él estaba construyendo. Para Óscar, era solo un juego más, pero para Julio, era el comienzo de algo mucho más grande, algo con un potencial que lo dejaba sin aliento.

\begin{quote}
    —Gracias, Pi —dijo—. Desde ahora, acceso limitado. Solo podrás recibir órdenes con mi voz, y solo a través de mi computadora principal. Es una medida de seguridad.

    —Entendido —respondió Pi, su voz robótica sonando complacida.
\end{quote}

En el centro de su habitación, donde antes Fox dormía sin preocupaciones, se erguía ahora una estructura impresionante. Era como una escultura hecha de cables y metal, con un alma tecnológica. Todavía no tenía un nombre, pero Julio ya la sentía como una parte de sí mismo, una extensión de su cerebro y sus manos. Cada pieza tenía una historia: el **router Cisco reprogramado** que funcionaba como el cerebro central; los **tableros programables (los Arduino)** que controlaban los sensores de flujo de energía; las **baterías alineadas** como columnas que guardaban una energía inmensa, lista para ser usada.

Faltaba poco. Muy poco para que el proyecto estuviera terminado y para que esa estructura cobrar vida.

Y en su mente, un pensamiento ya empezaba a crecer, una pregunta que lo llenaba de emoción y nervios: ¿Y si realmente funciona? ¿Y si soy el primero en ver algo que nadie más ha visto?

Julio no lo decía en voz alta. No quería tentar al destino, es decir, no quería provocar que algo malo pasara por hablar demasiado de sus esperanzas. Pero algo muy profundo dentro de él… ya había cruzado el umbral, el punto de no retorno. Sabía que su vida ya no sería la misma.
\cleardoublepage % Para que el capítulo empiece en una página impar

% Cuarto capítulo
\chapter*{4 – Algo que nunca existió}
\addcontentsline{toc}{chapter}{4 – Algo que nunca existió} % Añadir al índice manualmente

Julio no era de los que se rinden fácilmente. Desde muy pequeño, cada error que cometía no era un fracaso, sino una invitación a entender por qué las cosas pasaban. Para él, un tropiezo solo significaba una pista para aprender. Pero aquella semana, después de muchísimos intentos, por primera vez, sintió que la duda se metía en su cabeza. Llevaba ya seis largos meses trabajando sin parar en este proyecto, desde aquel día en que su papá le había traído la caja llena de tesoros electrónicos. Seis meses de noches sin dormir, de cálculos interminables, de piezas que se conectaban y se desconectaban.

La máquina, a la que no le había puesto nombre, estaba terminada, o al menos eso pensaba. Era una estructura impresionante, del tamaño de un pequeño refrigerador, como esos que se usan en las oficinas para guardar refrescos. Su carcasa, hecha de metal reciclado, brillaba con un aspecto único, uniendo piezas de chatarra con un propósito mayor. Los cables, de diferentes colores y grosores, recorrían su interior como las venas de una criatura dormida, llevando la energía a cada rincón. Y en el centro, el corazón de todo, estaba el **núcleo principal**, que era nada menos que la caja modificada del viejo **microondas** de Don Clemente. Todo había sido armado con una precisión casi perfecta, como si cada tornillo y cada cable tuvieran su lugar exacto en el universo. Había usado cada componente que su padre había conseguido, y muchos más que había ido recolectando: las computadoras **Raspberry Pi**, el **router Cisco** (que ahora solo servía para unas cosas específicas, no para el internet), los **Arduinos** (esos pequeños cerebros programables), **condensadores** (que guardan energía eléctrica) y hasta pequeñas piezas electrónicas rescatadas de juguetes rotos y electrodomésticos descompuestos. Era una obra de arte hecha de lo que para otros era basura.

La primera vez que la activó, fue un desastre total.

\begin{quote}
    —Pi, iniciar protocolo Alfa —dijo Julio, su voz sonando grave por la emoción, con los ojos fijos en la máquina.

    —Confirmado —respondió la voz de Pi, clara y sin emociones—. Encendiendo módulo de activación en tres… dos…
\end{quote}

Entonces, una chispa, como un pequeño rayo, saltó del interior de la máquina. Luego, un zumbido fuerte, que al principio era bajo, pero que fue creciendo y creciendo hasta convertirse en un chillido eléctrico agudo y desesperante que lo llenó todo. Las luces de la habitación de Julio y luego las de toda la casa parpadearon con fuerza, una, dos, tres veces. Y en menos de diez segundos, todo el vecindario quedó sumido en la oscuridad. Un **apagón total**. El silencio se hizo pesado, solo roto por el sonido de las casas vecinas que también se quedaban a oscuras. El viejo router, que antes había servido para internet, explotó con un pequeño "pop" y una nube de humo, y Fox salió disparado del cuarto, corriendo como si hubiera visto un fantasma, bufando con el lomo erizado.

Julio se quedó quieto en medio de la oscuridad absoluta de su habitación, iluminado solo por el débil resplandor rojo de una batería de emergencia que se había activado automáticamente. Respiraba con dificultad, el pecho subiéndole y bajándole rápidamente, y sus manos temblaban de rabia y frustración.

\begin{quote}
    —¡Maldita sea! —murmuró, la palabra escapándose entre dientes apretados.
\end{quote}

Con un golpe de frustración, le dio una patada suave al costado de la máquina, haciendo que una tapa suelta, hecha de metal reciclado, cayera al suelo con un ruidoso "clang" metálico que resonó en el silencio. La frustración lo invadió como una ola negra y pesada, ahogándolo. No podía más. No después de tanto tiempo, tanto esfuerzo, tantas noches sin dormir.

Se sentó en el suelo frío de su habitación, con la cabeza entre las rodillas, sintiendo que cada uno de sus pensamientos se deshacía, se desmoronaba como las piezas de un rompecabezas mal puestas. Había imaginado mil veces ese momento: luces brillantes girando, un **vórtice** (como un remolino de energía o un agujero giratorio) abriéndose, tal vez incluso un portal a otro lugar… pero solo había logrado un gran apagón en todo su barrio y un router quemado. Estaba al borde de rendirse, de dejarlo todo, de tirar la toalla.

\begin{quote}
    —¿Estás bien? —La voz de Óscar, su hermano menor, lo sacó de su burbuja de desesperación.
\end{quote}

Julio levantó la vista, sorprendido. No había escuchado a su hermano entrar al cuarto. Óscar estaba parado en la puerta, en pijama, con su cabello desordenado por la almohada y un pequeño muñeco de trapo colgando sin vida de su mano. Se veía pequeño y preocupado en la oscuridad.

\begin{quote}
    —¿Otra vez explotó? —preguntó Óscar, su voz llena de una mezcla de curiosidad y la experiencia de haber visto esto antes. No era la primera vez que algo así pasaba, claro.
\end{quote}

Julio no respondió. No tenía ganas de hablar. En ese momento, escucharon pasos en el pasillo. La puerta de la habitación se abrió un poco más y la luz de una linterna apuntó directamente a Julio. Era Ana, su madre, con el rostro preocupado y la voz que intentaba sonar tranquila, pero con un claro toque de molestia.

\begin{quote}
    —¡Julio! —dijo Ana, su voz grave—. ¡Otro apagón! ¡Esto ya es el tercero en dos meses! ¡Estábamos viendo una película! ¿Qué es lo que estás haciendo aquí que no lo puedes controlar? ¡La gente del barrio está molesta y con razón!
\end{quote}

Julio solo pudo encogerse un poco, sintiendo el peso de la reprimenda. Su madre suspiró y se dio la vuelta, apagando la linterna y dejando el cuarto de nuevo en penumbra. Sus pasos se alejaron, dejando a Julio con el peso de su frustración y la preocupación de su madre.

Óscar esperó a que se fuera la mamá y se acercó un poco más a Julio.

\begin{quote}
    —Cuando yo no puedo armar mis Legos… lo que hago es cambiar la forma del dibujo —dijo Óscar, como si hablara de la cosa más sencilla del mundo—. No la fuerza. Si algo no encaja… quizás es que no es para encajar ahí.
\end{quote}

\begin{quote}
    —¿Cambiar el dibujo? —preguntó Julio, extrañado.
\end{quote}

\begin{quote}
    —Sí. Como en esos rompecabezas donde las piezas se parecen mucho… pero no son la misma. Tal vez estás tratando de forzar algo que no necesita fuerza, sino que lo veas de otra manera.
\end{quote}

Julio lo miró en silencio, bajo la tenue luz roja de la batería. Por un instante, no vio a un niño molesto, hiperactivo e imprudente, que siempre estaba metiéndose en problemas. Vio a alguien que, de su manera algo torpe y totalmente honesta, había dicho exactamente lo que él necesitaba escuchar en ese momento de desesperación. Era una simpleza que encerraba una gran verdad.

\begin{quote}
    —Gracias, Óscar —murmuró Julio, sintiendo un peso menos en su corazón.
\end{quote}

Su hermano se encogió de hombros, soltó un bostezo enorme que casi le disloca la mandíbula y se fue sin más, dejando a Julio solo de nuevo con sus pensamientos y su máquina.

Las horas siguientes fueron de reconstrucción y mucho análisis. Julio no se dio por vencido. Inspirado por las palabras de Óscar, desmontó gran parte del núcleo de la máquina. No se trataba de usar más potencia, sino de que la potencia fuera controlada y bien dirigida. Con la ayuda de **Pi**, su asistente virtual, que nunca dormía, rediseñó la distribución de la energía, haciendo que fluyera de una manera más suave y segura. Ajustó el **pulso inicial**, que es como la primera descarga de energía, para que fuera más delicado y no provocara una sobrecarga.

Instaló una **base de inducción casera** (una especie de sistema que permite que la energía se mueva sin cables, controlando el voltaje o la fuerza de la electricidad). Cambió la polaridad de algunos **capacitores** (los capacitores son como pequeños almacenes de energía que guardan y liberan electricidad, y cambiar su polaridad es como invertir sus polos positivos y negativos para que la energía fluya de otra manera). Incluso conectó a Fox, el gato, como un sujeto pasivo de medición, no para dañarlo, sino porque el gato insistía en dormir sobre el **router quemado**, que ahora era completamente inútil. Así, Pi podía registrar si había alguna alteración en el campo magnético alrededor de Fox, que era muy sensible.

El núcleo principal de la máquina, como se mencionó, era la caja del microondas, pero no una caja cualquiera. Julio la había elegido por una razón muy importante: era excelente para soportar **radiación** en su interior. Esto era clave para lo que él planeaba. Dentro de esta caja, que ahora era como una especie de cámara especial, había instalado varios **sensores**. Había sensores de **humedad** para medir la cantidad de agua en el aire, sensores de **temperatura** para saber el calor exacto, y **escáneres** elaborados con las viejas **webcams** que había encontrado. Pi tenía acceso remoto a todos estos escáneres, lo que significa que podía ver lo que ocurría dentro de la caja y controlar los sensores sin necesidad de que Julio estuviera físicamente allí.

\begin{quote}
    —Pi… última prueba. Solo si todo está estable —dijo Julio con voz ronca por el cansancio y la tensión, sus ojos fijos en los indicadores de Pi.

    —Secuencia verificada. Iniciando protocolo experimental… ahora —respondió Pi, su voz con un tono ligeramente diferente, casi con una emoción contenida.
\end{quote}

El zumbido regresó, pero esta vez era diferente. No era el grito eléctrico que había provocado el apagón. Era más suave, armónico, es decir, un sonido que sonaba bien, como una melodía misteriosa. El aire del cuarto vibró, como si una onda invisible lo recorriera, y una luz azul brillante comenzó a salir del centro de la máquina, justo de la caja del microondas. Fox se levantó, curioso, con las orejas erguidas, su cola moviéndose despacio, observando el espectáculo luminoso.

De pronto, una pequeña **implosión**, un sonido como si el espacio mismo respirara hacia adentro, una succión rápida. Fue un sonido agudo, casi inaudible, tan bajo que costaba escucharlo. El resplandor azul se desvaneció de golpe, dejando el cuarto en la oscuridad de nuevo, solo con el brillo de los pequeños LED de la máquina.

En el centro del núcleo, sobre una placa metálica limpia, había algo. Algo pequeño. Del tamaño de un dado de jugar. Era de textura metálica, pero con un brillo imposible de describir, algo que no parecía de este mundo. No era oro, ni plata, ni cobre. Era algo más. Algo que parecía responder a la luz, e incluso a los pensamientos de Julio, como si tuviera una vida propia.

Julio se acercó lentamente, temblando, la mano extendida hacia el objeto.

\begin{quote}
    —¿Qué es eso…? —murmuró, su voz casi inaudible.

    —Análisis en curso —dijo Pi, su voz con un tono más lento de lo normal, como si estuviera procesando algo increíblemente complejo—. Densidad: desconocida. Estructura atómica: no identificada. Elemento no catalogado en la tabla periódica.
\end{quote}

Julio sintió que el aire abandonaba sus pulmones, dejándolo sin aliento. Su corazón latía con fuerza en su pecho.

\begin{quote}
    —Pi… ¿qué estás diciendo?

    —Este material no es de este mundo —declaró Pi, y la simpleza de esa afirmación lo dejó helado.
\end{quote}

Fox se subió a su regazo, ronroneando fuerte, como si celebrara el increíble descubrimiento. Julio acarició su pelaje blanco y suave, mientras una sonrisa lenta, casi dolorosa de lo increíble que era, se dibujaba en su rostro. Era una mezcla de asombro, alivio y la certeza de que su vida acababa de cambiar para siempre.

\begin{quote}
    —Funcionó… —susurró.
\end{quote}

\cleardoublepage % Para que el capítulo empiece en una página impar

% Quinto capítulo
\chapter*{5 – Un misterio brillante}
\addcontentsline{toc}{chapter}{5 – Un misterio brillante} % Añadir al índice manualmente

Julio no durmió esa noche. Ni siquiera lo intentó. Era imposible. La emoción y la adrenalina lo mantenían despierto, con la mente acelerada como un motor de carreras. El pequeño fragmento brillante, ese objeto misterioso que había aparecido de la nada, descansaba ahora con mucho cuidado dentro de la caja del microondas, en el centro de su máquina. Para tenerlo seguro, Julio lo había puesto en una **cápsula improvisada**, hecha con un trozo de vidrio de laboratorio reciclado, como un vaso de precipitado cortado, y sellada con varias capas de cinta adhesiva muy resistente. La caja del microondas, por su capacidad de contener radiación y mantener lo de adentro muy bien aislado, era el lugar perfecto para guardarlo, por miedo a que fuera radiactivo o peligroso, aunque Pi ya había dicho que no lo era.

Fox, vigilaba la cápsula de cerca desde su rincón favorito: un montón de cables enredados que se convertían en una especie de nido. Cada tanto, lanzaba un maullido suave, casi un gruñido, como si le estuviera diciendo a su dueño que no se fiara demasiado, que tuviera cuidado con ese objeto tan extraño. Los ojos de Fox parecían brillar en la oscuridad, fijos en el misterioso fragmento.

Antes de hacer cualquier otra cosa, Julio sabía que debía responder la pregunta más importante, la que lo había mantenido en vilo toda la noche: ¿era peligroso?

Al día siguiente con mucho cuidado, y sintiendo una mezcla de nervios y emoción, tomó sus guantes de goma, de esos que se usan para limpiar, y se puso una máscara de protección, de esas que su madre guardaba para limpiar con productos fuertes. Luego, abrió la cápsula improvisada y con pinzas, colocó el fragmento sobre una **balanza digital**. Lo miró con curiosidad. Era demasiado pesado para su tamaño, mucho más de lo que esperaba para algo tan pequeño. Era realmente curioso.

\begin{quote}
    —**Pi**, escanea niveles de radiación: **gamma**, **alfa** y **beta** —ordenó Julio, su voz un poco tensa, aunque confiaba en los análisis de Pi. Quería estar 100% seguro.

    —Lectura en curso… —respondió el asistente virtual, su voz robótica sonando a través de los parlantes—. Radiación: Negativa. Nivel de partículas inestables: cero. Composición ionizante: inexistente.
\end{quote}

Julio parpadeó. Era difícil de creer.

\begin{quote}
    —¿Nada? ¿En serio? ¿Absolutamente nada de radiación?

    —Confirmado, Julio —dijo Pi con su voz neutral, pero con un matiz de certeza—. Este material es completamente estable. No representa ningún peligro conocido.
\end{quote}

\begin{quote}
    —¡Sí! —festejó Julio, levantando ambos brazos al cielo con una alegría inmensa, como si hubiera ganado un premio. Fox, que no entendía nada de la euforia de su humano, maulló desconcertado, moviendo la cola.
\end{quote}

Con esa confirmación tan importante, la certeza de que el material era seguro, era hora de la mejor parte: ¡Experimentar! Para Julio, no había nada más emocionante que poner a prueba los límites de lo desconocido.

---

Primero, Julio decidió probar el material con **calor extremo**. Quería ver si se derretía o cambiaba con las altas temperaturas.

Colocó el fragmento cuidadosamente dentro de la cámara del microondas, debajo de una **resistencia eléctrica improvisada**, que había fabricado con las partes de una tostadora vieja. Esta resistencia, al calentarse, emitía un calor intenso y directo. Subió la temperatura muy despacio, observando atentamente en el monitor los números que indicaban los grados.

\begin{quote}
    —500 grados Celsius… 600… 700… —recitaba Julio en voz baja, mientras los números seguían subiendo, cada vez más. El calor en el aire se hacía sentir, pero el fragmento seguía igual.
\end{quote}

Nada. Ni se derritió, ni cambió de color, ni se deformó. Ni un solo rasguño visible en su superficie brillante. El chico golpeó la mesa con la palma de la mano, emocionado por el hallazgo.

\begin{quote}
    —¡Resiste más que el **acero de tungsteno**! —exclamó, refiriéndose a un metal súper resistente, uno de los más duros que existen. Era como si el material se burlara del calor.
\end{quote}

La siguiente prueba sería con **electricidad**. Julio quería ver cómo reaccionaba el material a las grandes descargas eléctricas.

Conectó los **electrodos** (que son como las puntas de metal que conducen la electricidad) de un viejo **generador** que había construido cuando era pequeño. Era una máquina casera que podía producir chispas de alto voltaje. La chispa azul, un rayo diminuto, saltó sobre el material con un sonido seco y crepitante. Durante cinco minutos interminables, Julio lo expuso a descargas de alto voltaje, una y otra vez. Los destellos de la chispa iluminaban su rostro concentrado en la oscuridad de la habitación. Fox, sintiendo la energía en el aire, se escondió rápidamente bajo una pila de mantas, mirando a su dueño como si estuviera completamente loco, sus ojos brillando con una mezcla de miedo y reproche.

Julio detuvo el generador. El material, como si nada hubiera pasado, seguía exactamente igual, brillando con su resplandor particular.

\begin{quote}
    —Fox, creo que acabo de encontrar el primer superhéroe de los minerales —bromeó Julio, intentando aligerar el ambiente, mientras el gato lo miraba indiferente desde su escondite, sin entender la broma.
\end{quote}

El tercer experimento sería con **reacciones químicas**. ¿Podría algún ácido o sustancia fuerte dañarlo?

Julio colocó el fragmento en un vaso con agua. Nada. Absolutamente ninguna reacción. Agregó vinagre, un ácido suave. Nada. Luego, lejía, un producto de limpieza fuerte. Tampoco. Finalmente, alcohol isopropílico, que se usa para limpiar componentes electrónicos. El material seguía intacto.

\begin{quote}
    —¡Nada, nada, y nada! —gritó, lanzando las manos al aire, entre la frustración y la sorpresa.
\end{quote}

En un momento de inspiración (o quizás de pura locura de científico), decidió combinar varios líquidos en un mismo recipiente, una mezcla que normalmente sería muy peligrosa: **cloro**, **amoníaco**, refresco de cola y hasta un poco de salsa picante que había sobrado de la cena anterior. La mezcla burbujeó de inmediato. El resultado fue una burbuja de vapor pestilente, es decir, que olía fatal, un olor horrible que llenó el cuarto en segundos. Fox, que ya había salido de su escondite, salió corriendo a toda velocidad como si lo hubieran disparado de un cañón, huyendo del olor insoportable.

Julio, tosiendo y tapándose la nariz, cubrió la reacción con un frasco de vidrio para contener los vapores. Lo más sorprendente fue que el material, en medio de ese desastre químico y el olor desagradable, brillaba… ¡más fuerte! Como si se estuviera alimentando del caos, de la mezcla de sustancias.

\begin{quote}
    —Oh, esto es interesante… —murmuró Julio para sí mismo, la tos aún en su garganta, pero una nueva idea ya formándose en su mente.
\end{quote}

Finalmente, intentó algo que había visto en una vieja serie de televisión: **presión extrema**. Quería ver si el material podía ser aplastado.

Colocó el fragmento cuidadosamente entre dos bloques de **acero sólido**. Luego, utilizó una vieja **prensa hidráulica** de su padre, una máquina que usa líquido a presión para aplicar una fuerza enorme, capaz de aplastar metales.

La prensa bajó… bajó… lentamente, con un chillido metálico que indicaba la gran fuerza que estaba aplicando… ¡y de repente, se rompió! Un fuerte crujido metálico resonó en la habitación, y la prensa dejó de funcionar.

Julio se quedó con la boca abierta, los ojos muy abiertos, incapaz de creer lo que acababa de pasar.

\begin{quote}
    —¡¿Qué?! ¡¿Qué clase de broma cósmica es esta?! —exclamó, mirando los restos de la prensa con incredulidad.
\end{quote}

Se acercó rápido a revisar: el fragmento no tenía ni un solo rasguño, ni una marca, ni una abolladura. Estaba perfecto.

Fox apareció en la puerta, con un aire de superioridad felina, como diciendo: "Te lo dije, humano. Nunca me escuchas." Su pose relajada y sus ojos astutos parecían confirmarlo.

Después de limpiar el desastre químico y asegurarse de que el sótano no oliera a gas tóxico o a salsa picante mezclada con cloro (lo que hubiera sido muy peligroso), Julio se sentó frente a su laptop, con la mente zumbando de nuevas ideas.

\begin{quote}
    —Pi, en base a todas las pruebas que hemos hecho, ¿qué propiedades podemos confirmar de este material? —preguntó Julio, ansioso por las conclusiones.

    —Material **indestructible** bajo condiciones humanas normales —dijo Pi, enumerando los resultados—. No es **conductor de electricidad**. Es resistente al **calor extremo**. No es **tóxico**. No es **ionizante**, es decir, no emite radiación. Y parece ser **energéticamente reactivo de manera positiva ante ambientes caóticos**.
\end{quote}

Julio tamborileó los dedos en la mesa, su mente corriendo a mil por hora, conectando todos los puntos. La última frase de Pi le dio una idea brillante.

\begin{quote}
    —Pi… ¿podríamos usar este material para reemplazar el compuesto radiactivo de la máquina? ¿El **americio-241** que estábamos buscando?

    —Confirmado, Julio —respondió Pi—. No solo lo reemplazaría, sino que aumentaría la eficiencia energética en un **437%**. Además, eliminaría por completo cualquier riesgo de contaminación o daño a la salud. Sería mucho más seguro y potente.
\end{quote}

Julio dejó caer la espalda contra la silla, suspirando de alivio y euforia. Sonreía como si acabara de descubrir un nuevo universo, o como si hubiera ganado la lotería. Era un hallazgo que cambiaba todo.

\begin{quote}
    —Fox… amigo mío, no solo hicimos algo increíble —dijo, la voz llena de asombro y orgullo.
\end{quote}

Fox saltó sobre su regazo, ronroneando fuerte, como si celebrara el gran descubrimiento con él.

\begin{quote}
    —Hicimos algo que nadie más en el mundo tiene —dijo Julio, acariciando el pelaje blanco del gato.
\end{quote}

El futuro de su invento, de su máquina, acababa de cambiar para siempre. Y él, Julio, el chico que miraba las estrellas, apenas estaba comenzando. Lo que venía ahora, era mucho más grande de lo que cualquiera podría imaginar.

\cleardoublepage % Para que el capítulo empiece en una página impar

% Sexto capítulo
\chapter*{6 – Lecciones del Hogar}
\addcontentsline{toc}{chapter}{6 – Lecciones del Hogar} % Añadir al índice manualmente

Julio caminaba por los pasillos del colegio como un fantasma más, una figura que pasaba desapercibida entre el bullicio de los otros estudiantes. Ya habían pasado siete largos meses desde que comenzó su ambicioso proyecto, y cada día lo absorbía más y más. Su ropa, que antes solía estar impecable, ahora estaba arrugada, como si se la hubiera puesto sin siquiera mirarla. Había olvidado peinarse, y su cabello revuelto era un claro signo de que las mañanas se le iban en otras cosas. Tenía unas ojeras tan marcadas y oscuras debajo de los ojos que parecía no haber dormido en días, o quizás en semanas. Entre sus manos, arrugadas y llenas de pequeñas manchas de tinta y grasa, estrujaba hojas y más hojas llenas de cálculos complicados y bocetos de su máquina mejorada, dibujos que solo él podía entender. No prestaba atención a las miradas curiosas o a los murmullos de sus compañeros, estaba demasiado absorto en sus propios pensamientos, en su mundo interior lleno de engranajes y fórmulas. Para él, el colegio se había vuelto un simple telón de fondo.

Este evidente descuido, aunque no era escandaloso ni causaba problemas directos en clase, era notorio y no pasó desapercibido para los profesores. Ellos, que conocían a Julio desde hacía años y sabían de su brillantez, notaban su ausencia mental, su falta de participación en las clases, donde antes era el primero en responder. Veían cómo su atención se perdía en la ventana, cómo su mirada se fijaba en la nada mientras sus manos dibujaban en los márgenes de sus cuadernos. Sentían no solo la obligación, sino también una profunda preocupación por lo que sea que estuviera pasando con Julio. Sabían que algo andaba mal, y que ese cambio no era normal en un chico tan inteligente.

Al final del día escolar, justo cuando el sol empezaba a bajar y la mayoría de los estudiantes corrían para irse a casa, Julio fue llamado a la dirección. El director, un hombre de voz firme pero con una expresión de clara preocupación en el rostro, le habló de su apariencia desmejorada, su falta de participación en clase y su evidente desconexión con el entorno. Le explicó que no era solo por sus notas, sino por su bienestar. Fue entonces cuando decidieron que era necesario hablar con sus padres. No podían dejar pasar lo que estaba ocurriendo.

Cuando Miguel y Ana llegaron al colegio unas horas más tarde, Julio los esperaba sentado en una silla fuera de la oficina del director. Estaba encogido, con los hombros caídos y el rostro aún más pálido de lo normal. Ana, su madre, sintió una punzada en el corazón al verlo así. Se agachó frente a él, hasta quedar a su altura, y con una mezcla de ternura y una profunda preocupación en sus ojos, le acarició el cabello desordenado, tratando de alisarlo un poco.

\begin{quote}
    —¿Qué está pasando, Julio? —preguntó en voz baja, con una suavidad que invitaba a la confianza.
\end{quote}

Julio bajó la mirada, avergonzado. Sentía que había fallado, que había decepcionado a sus padres. No sabía por dónde empezar a explicar el inmenso y secreto mundo que había construido en su habitación. Miguel, serio pero con una mirada comprensiva, puso una mano firme y cálida sobre su hombro, dándole apoyo.

\begin{quote}
    —Sea lo que sea, hijo —dijo Miguel, su voz era un pilar de seguridad—, estamos aquí para ayudarte. No tienes que enfrentar esto solo.
\end{quote}

En el camino a casa, el ambiente en el auto era silencioso, pesado. Las palabras dichas en el colegio, las miradas de los profesores, la preocupación de sus padres… todo flotaba en el aire. Fox, esperaba en la ventana de la casa, moviendo su cola de un lado a otro, como si hubiera sentido la tensión en el aire antes de que llegaran. Cuando Julio entró, Fox, con un salto ágil, brincó a su regazo, ronroneando con fuerza, un sonido vibrante que llenó el espacio. Julio lo abrazó con fuerza, hundiendo su rostro en el suave pelaje blanco del gato, encontrando en él una pequeña, pero necesaria, dosis de consuelo. Era como un abrazo sin palabras, una conexión que solo ellos dos compartían.

Esa noche, mientras cenaban juntos en la cocina, con el suave aroma a comida casera, Ana y Miguel decidieron hablar con calma, sin regaños, sino con amor.

\begin{quote}
    —Sabemos que eres especial, Julio —dijo Ana, su voz suave, mientras le servía un poco más de sopa caliente en su plato—. Siempre has sido diferente a los demás niños. Más inteligente, más curioso… te adentras en cosas que nadie más entiende. Pero tienes que cuidar de ti también. Tu cuerpo y tu mente son importantes.

    —Tu mente es brillante, hijo —añadió Miguel, señalándolo con la cuchara, como si subrayara sus palabras—. Tienes ideas asombrosas y creas cosas increíbles. Pero un gran invento no sirve de nada si el inventor se olvida de sí mismo. No queremos que te apagues por dentro, que te pierdas a ti mismo, mientras persigues tus sueños. Queremos verte feliz y sano.
\end{quote}

Julio, conmovido por las palabras sinceras de sus padres, dejó la cuchara a un lado, su plato aún con comida. Levantó la mirada y los miró a ambos, sintiendo un nudo en la garganta. La culpa y la tristeza se mezclaron con un poco de alivio.

\begin{quote}
    —Lo siento… —murmuró, su voz apenas un susurro—. No quise preocuparlos. Es que… a veces siento que no encajo en ningún lado, que soy diferente. Y la única forma de encontrar mi lugar, de sentir que pertenezco, es construyendo algo que solo yo pueda entender, algo que solo existe en mi mente.
\end{quote}

Ana se acercó a él, rodeándolo con sus brazos en un abrazo fuerte y cálido. El calor de su cuerpo era un bálsamo para Julio.

\begin{quote}
    —Nunca tienes que cargar solo con el peso de tus sueños —susurró ella, apretándolo más contra sí.
\end{quote}

Miguel, sonriendo, revolvió el cabello de su hijo, un gesto de cariño que le era muy familiar.

\begin{quote}
    —Y si alguna vez sientes que el mundo no te entiende —dijo, sus ojos llenos de amor—, recuerda: nosotros siempre lo haremos. Siempre estaremos aquí para ti.
\end{quote}

Esa noche, antes de dormir, Julio se prometió a sí mismo que cuidaría mejor de su cuerpo y de su mente. Sabía que no podía dejar que su pasión lo consumiera por completo. La promesa era para él mismo, pero también para sus padres. La pasión por su proyecto seguía ardiendo con la misma fuerza, pero ahora, también lo acompañaba la conciencia de que no estaba solo.

Después de ducharse con agua tibia, de ponerse ropa limpia y cómoda, y de ordenar un poco su espacio de trabajo, que siempre estaba lleno de herramientas y piezas, Julio se sentó frente a su máquina. Fox, como siempre, estaba acurrucado a su lado, en una posición que parecía la de un pequeño centinela peludo, esperando el momento. La habitación estaba más limpia de lo habitual, pero el aire aún olía a metal y electricidad.

Con el nuevo **material brillante** ya integrado en el corazón de la máquina y todas las modificaciones aplicadas gracias a los descubrimientos y pruebas de los últimos días, Julio encendió su creación. Su mano temblaba un poco, pero esta vez no era por miedo, sino por la emoción.

Un **zumbido bajo y constante** llenó el aire, una vibración suave que se sentía en los huesos, desde el suelo hasta el cuerpo de Julio. Luces tenues, de colores azul y verde, parpadearon en los instrumentos improvisados que rodeaban la máquina, como pequeños ojos vigilantes. Esta vez, a diferencia de la anterior, no hubo chispazos, ni apagones, ni humo negro saliendo de ningún lado. Solo un brillo suave, constante, como una respiración contenida de la máquina. La energía fluía sin problemas.

La máquina estaba lista.

Julio acarició a Fox, que ronroneaba feliz sobre su pierna, con el corazón lleno de una emoción indescriptible. Sonrió para sí mismo, una sonrisa sincera y llena de esperanza.

\begin{quote}
    —Esta vez no estoy solo —murmuró, la voz casi un suspiro, refiriéndose no solo a **Pi**, sino también al apoyo de sus padres y a la presencia de Fox.
\end{quote}

Y con esa certeza, esa sensación de apoyo y de un propósito mayor, el futuro dejó de parecer un misterio aterrador, un camino oscuro lleno de dudas, y empezó a sentirse como una promesa emocionante, un nuevo capítulo que estaba a punto de comenzar.

\cleardoublepage % Para que el capítulo empiece en una página impar

% Séptimo capítulo
\chapter*{7 – ¿Una línea de comandos?}
\addcontentsline{toc}{chapter}{7 – ¿Una línea de comandos?} % Añadir al índice manualmente

Justo después de que las luces tenues y azuladas parpadearan en los instrumentos improvisados, la máquina de Julio, comenzó a vibrar de una manera casi imperceptible. Era un zumbido suave, como el ronroneo más bajo que un gato pudiera hacer, una vibración que se sentía más en el aire que en el oído. La máquina parecía respirar por sí misma, llena de una energía silenciosa y controlada. El **nuevo material brillante** que había encontrado y colocado en su corazón funcionaba exactamente como se esperaba, transformando la energía de una manera limpia y armoniosa, sin chispazos ni apagones.

Después de todo ese brillo y zumbido, de repente, la habitación se quedó en un silencio casi absoluto, un silencio expectante. Las luces de la máquina se atenuaron, pero en el centro del dispositivo, en su corazón, es decir, dentro de la caja de microondas que Julio había modificado con tantos sensores y escáneres, algo increíble apareció. Un pequeño **punto negro**, apenas más grande que la cabeza de un alfiler, flotaba suspendido en el aire, justo en el centro de la **cámara de experimentación**. No emitía luz propia, ni calor, pero parecía absorber la energía a su alrededor, como si fuera un pequeño agujero que se tragara la luz. Era un núcleo denso y sereno, tan misterioso como fascinante.

Julio lo observaba con una mezcla de asombro y nerviosismo. Su corazón latía con fuerza contra sus costillas. Sus ojos, ahora bien abiertos, no podían apartarse de aquel fenómeno. Fox, su fiel gato, seguía ronroneando a su lado, completamente ajeno a la magnitud del momento, a la importancia de lo que estaba ocurriendo justo delante de sus narices.

\begin{quote}
    —¿Qué demonios es eso? —susurró Julio, la pregunta escapándose de sus labios, sin atreverse a pestañear.
\end{quote}

**Pi**, su asistente de inteligencia artificial, respondió de inmediato. Comenzó a emitir una serie de sonidos electrónicos, como pequeños clics y susurros digitales, mientras realizaba múltiples análisis al mismo tiempo, procesando una cantidad gigantesca de información sobre el punto negro. Cada sensor dentro de la caja del microondas (los de humedad, temperatura, y los escáneres hechos con las webcams) enviaba datos a Pi sin parar.

\begin{quote}
    —Julio —dijo Pi, su voz modulada con una calma que contrastaba con los nervios de Julio—, he detectado que el fenómeno en el centro de la máquina no corresponde a ningún fenómeno natural conocido. No hay nada igual en las bases de datos terrestres. Sin embargo, he deducido por patrones de comportamiento (es decir, por la forma en que se comporta y reacciona) que se asemeja a un "**prompt**".

    —¿Un **prompt**? —repetió Julio, incrédulo. Un prompt es como una señal o un mensaje que te pide que hagas algo, como el cursor que espera que escribas en la pantalla de una computadora.

    —Sí —explicó Pi con su habitual eficiencia, sin rodeos—. Se comporta como una **línea de comandos**. Imagina que es como una ventana en la computadora donde tienes que escribir una instrucción para que algo suceda. Es un sistema esperando instrucciones. Como un "**Hola Mundo**" en un programa recién creado, la primera frase que se usa para probar si algo funciona, aguardando su primera orden para cobrar vida.
\end{quote}

Julio soltó una pequeña carcajada nerviosa, aliviando un poco la tensión que sentía. La idea le parecía tan loca que era graciosa.

\begin{quote}
    —¿Qué tal si le enviamos literalmente eso? —dijo Julio, con un tono de broma, pero con la chispa de la curiosidad brillando en sus ojos—. Un "**Hola Mundo**" en **binario**, a ver qué pasa.

    —Procesando… —respondió Pi, y Julio ya sabía que cuando Pi decía eso, significaba que lo haría de inmediato.
\end{quote}

En cuestión de segundos, Pi, con su velocidad asombrosa, tradujo el mensaje simple "Hola Mundo" a **código binario puro**. El código binario es el lenguaje básico de las computadoras, compuesto solo por ceros y unos, como si fueran interruptores de luz encendidos o apagados (**01001000 01101111 01101100 01100001 00100000 01001101 01110101 01101110 01100100 01101111**). Luego, lo envió a través de un **campo de frecuencia controlada** (una especie de onda invisible y muy precisa, como una señal de radio muy especial) directo hacia el pequeño punto negro.

El aire en la habitación pareció comprimirse de golpe, como si alguien hubiera apretado un botón invisible. Se sintió una presión sutil, como si el espacio mismo se encogiera. La máquina zumbó un poco más fuerte, el punto negro parpadeó con una luz interna que nunca antes había emitido… y, de repente, con un sonido suave, casi como un suspiro, algo cayó en la **bandeja de recolección** que estaba justo debajo del núcleo, dentro de la caja del microondas.

Julio se apresuró a mirar. Su corazón latía a mil por hora. Lo que encontró era algo completamente inesperado, más allá de cualquier cosa que hubiera imaginado.

En la bandeja de recolección había una pequeña **esfera de cristal transparente**, no más grande que una canica de juego. Lo realmente asombroso no era la esfera en sí, sino lo que había dentro: diminutas **letras flotando** en su interior. Letras que parecían formar, en infinitos patrones cambiantes, las palabras "**Hola Mundo**" en **idiomas que Julio jamás había visto o escuchado**. Eran símbolos extraños, lenguajes que no existían en la Tierra, moviéndose y brillando suavemente dentro del cristal. Era como si la esfera fuera un mensaje, una ventana a otra forma de comunicación.

Julio la sostuvo en la palma de su mano, sintiendo la superficie fría y suave del cristal. No pesaba nada. Era como si el objeto no obedeciera por completo a las leyes de su mundo, a la gravedad o a la materia como él la conocía. Las letras flotantes brillaban con una luz suave y constante, hipnótica, que te invitaba a mirarlas sin parar, atrapando su atención.

\begin{quote}
    —Pi, análisis —ordenó Julio, su voz apenas un hilo de asombro.

    —El objeto no posee **estructura molecular reconocible** —informó Pi, su voz era tan tranquila como siempre, pero los datos que arrojaba eran extraordinarios—. Su composición no corresponde a ningún material terrestre. No está hecho de **átomos** como todo lo que conocemos. Además, está emitiendo una **frecuencia débil** que parece… amigable.

    —¿Amigable? —repetió Julio, entre asombrado y divertido por la descripción tan peculiar de Pi. ¿Cómo podía un objeto ser "amigable"?

    —Sí —explicó Pi—. Si tuviera que definirlo de manera más humana, es como si el objeto intentara decirnos: "Estoy aquí. ¿Qué sigue? Me han llamado y he respondido". Es una **señal de bienvenida**.
\end{quote}

Julio sonrió, sintiendo por primera vez que su creación no era simplemente una máquina, no era solo una colección de piezas de metal y cables. Era algo mucho más profundo, más asombroso. Era una puerta, un saludo, una invitación a algo mucho más grande de lo que jamás había imaginado. Era como si hubiera abierto una línea de comunicación con algo que estaba más allá de su comprensión.

Fox, ajeno a los grandes descubrimientos del universo, frotó su cabeza contra la pierna de Julio, buscando atención y un poco de cariño. Julio, todavía con la mente en la esfera brillante, acarició distraídamente el pelaje blanco del gato mientras miraba la canica que titilaba suavemente en su mano.

\begin{quote}
    —¿Una **línea de comandos**? —susurró para sí mismo, la frase de Pi resonando en su cabeza—. Esto apenas comienza.
\end{quote}

Y con esa certeza, esa sensación de haber tocado algo extraordinario y de haber abierto una ventana a lo desconocido, supo que su mundo, y quizás todos los mundos posibles, acababan de cambiar para siempre. La aventura apenas daba sus primeros pasos.

\cleardoublepage % Para que el capítulo empiece en una página impar

% Octavo capítulo
\chapter*{8 – Susurros desde otro lado}
\addcontentsline{toc}{chapter}{8 – Susurros desde otro lado} % Añadir al índice manualmente

Julio estaba sentado en su silla giratoria, con la mirada fija en la pantalla de su computadora. Entre sus dedos, daba vueltas a la pequeña **esfera de cristal transparente** que había aparecido de la anomalía, la "canica" que contenía las palabras "Hola Mundo" en idiomas imposibles. Desde aquel día, hace ya una semana, Julio había tomado una decisión muy importante: no enviar más "**prompts**" a la anomalía, es decir, no volver a escribirle mensajes a ese punto negro que funcionaba como una ventana a otro lado.

Tenía varias razones muy fuertes para esta cautela. La primera era que no sabía qué podría suceder si lo hacía. ¿Y si su siguiente mensaje abría una puerta a algo peligroso? La segunda razón era que no sabía exactamente qué tipo de "prompt" enviar. ¿Una pregunta? ¿Una orden? ¿Un saludo más complejo? Era como intentar hablar con alguien en un idioma desconocido sin saber qué decir. Y la tercera, y quizás la más importante, era que tenía una cierta sensación de **disconfort y preocupación**; una mezcla de inquietud y duda que le decía que no debía seguir adelante a la ligera. Si bien había logrado lo que tanto se propuso desde que tenía nueve años (y ahora, con catorce, ya era un adolescente), la idea de conectar con otros universos no era algo que hubiera pensado tan a fondo. Si era lo suficientemente inteligente como para crear una máquina de ese calibre, que podía traer cosas de otras dimensiones, también era lo suficiente inteligente como para ser cauto en su uso. La prudencia era clave.

Así que decidió que, antes de cualquier otro paso, lo primero sería crear una forma de **comunicación efectiva y más fácil** con el punto negro o "prompt". Quería evitarse el complicado paso de tener que enviar el texto en **código binario**, esos ceros y unos que solo las máquinas entienden. Para eso, pasó varios días trabajando en un **software especial**, una especie de programa que le permitiría hablar directamente con la anomalía. Finalmente, logró instalar una conexión segura y estable a una **terminal** en la pantalla de su computadora. Ahora, en lugar de códigos, tendría una ventana más amigable para comunicarse. No había enviado ningún "prompt" aún; solo estaba esperando, viendo la pantalla donde el cursor, una raya baja, parpadeaba, esperando una orden. Con la canica en los dedos, mientras la luz del monitor se reflejaba en sus ojos, algo increíble sucedió.

La pequeña esfera negra en el centro de la máquina, que flotaba en el aire dentro de la caja del microondas, latía con un ritmo suave, como un diminuto corazón que bombeaba misterio. Desde que había enviado aquel simple "Hola Mundo" en binario, Julio apenas podía despegar la vista de la pantalla de **Pi**, donde líneas de código y patrones imposibles bailaban sin cesar, mostrándole los datos en tiempo real de lo que sucedía.

De repente, la pantalla parpadeó. No había enviado ninguna pregunta, ninguna orden. Simplemente, de la nada, un mensaje comenzó a aparecer en su ventana de comandos. Pi, siempre atento, anunció con su tono tranquilo:

\begin{quote}
    —Se detecta una respuesta de **alta entropía** —dijo Pi, refiriéndose a un mensaje que contiene mucha información nueva y compleja, casi como si fuera algo totalmente inesperado y desconocido—. Intentando decodificar…
\end{quote}

Julio se frotó las manos nerviosamente, sus ojos fijos en la pantalla. ¿Qué clase de cosa respondería a un mensaje de "Hola Mundo" con algo tan complicado? ¿Y por qué le respondían sin que él preguntara nada?

La respuesta no tardó en llegar. Las letras se formaron en la pantalla, una por una, como si alguien las estuviera escribiendo justo en ese momento, directamente en su ventana de comandos. Y lo que apareció, lo dejó sin aliento:

\begin{quote}
    "Hola, Julio. Te estábamos esperando."
\end{quote}

Julio parpadeó, sintiendo un escalofrío que le recorrió la espalda hasta la nuca. La piel se le puso de gallina. ¿Cómo sabían su nombre? Miró a Fox, que dormitaba enroscado junto a una maraña de cables, ajeno a todo. Ni siquiera el gato parecía captar la magnitud del momento, el hecho de que una comunicación con otro lado acababa de llamarlo por su nombre.

\begin{quote}
    —¿Quién… quién me está hablando? —preguntó Julio en voz alta, su voz temblaba ligeramente, mientras tecleaba rápidamente la misma pregunta en el sistema de comunicación improvisado que acababa de crear. Quería una respuesta directa.
\end{quote}

Pi procesó la pregunta de Julio y la envió a la anomalía, y la respuesta apareció al instante en la pantalla:

\begin{quote}
    "Somos fragmentos de posibilidades. Ecos de realidades."
\end{quote}

Julio sintió otro escalofrío, más fuerte esta vez. ¿Era una broma? ¿Un error de interpretación de Pi? ¿O realmente había algo, o alguien, del otro lado de esa ventana a lo desconocido? La idea era tan inmensa que le costaba procesarla.

---

Mientras intentaba ordenar sus pensamientos, que parecían estar en una maraña en su cabeza, la máquina vibró ligeramente. Un sonido parecido a un **pop sordo**, como el aire escapando de un globo lentamente, se oyó. De repente, el **receptáculo de recolección**, la bandeja que estaba dentro de la caja del microondas y que Julio había diseñado para recoger cualquier objeto que apareciera, comenzó a llenarse. No era la primera vez que se llenaba, pero esta vez, la sensación era diferente.

Primero cayó un pequeño **destornillador**. Su superficie brillaba con una luz extraña, como si estuviera hecho de **cristal líquido**, que es un material que se dobla y cambia de forma como el agua, pero que es sólido. No era un destornillador común. Luego, cayó un **dado de seis caras**, de esos que se usan en los juegos de mesa. Pero en lugar de números, cada cara mostraba una miniatura de un **paisaje que parecía sacado de otro planeta**, un paisaje alienígena, con montañas de colores que no existen en la Tierra y cielos con lunas dobles.

Le siguió un **reloj de pulsera** que, en lugar de manecillas que marcan las horas, tenía diminutos **sistemas solares girando** en su interior. Pequeños planetas y soles diminutos se movían en órbitas perfectas dentro de su esfera de cristal, como un universo en miniatura.

Julio recogió los objetos, uno a uno, maravillado. Sus manos temblaban de emoción. Cada objeto era una prueba de que había un sinfín de realidades. Los colocó con cuidado en su escritorio, cerca de la esfera "Hola Mundo".

No había tenido tiempo de examinar cada uno de ellos a fondo cuando escuchó un golpe sordo: una botella había caído al suelo, rodando hasta detenerse frente a Fox, que abrió un ojo curioso y la olfateó con cautela.

La etiqueta de la botella le sacó una carcajada nerviosa. Era una botella de refresco, pero la etiqueta decía: "**OmniCola: Fuerza en cada sorbo**", y mostraba una caricatura de **Omniman**, su superhéroe favorito de los dibujos animados. Era un héroe de cómics que siempre había imaginado, pero que solo existía en la ficción.

\begin{quote}
    —No puede ser… —susurró Julio, recogiendo la botella con la mano temblorosa. La sentía fría y real.
\end{quote}

En ese instante, Pi emitió una alerta suave, un pequeño sonido que Julio ya reconocía como una advertencia.

\begin{quote}
    —Julio, los objetos parecen ser **manifestaciones físicas extrapoladas desde realidades alternativas**. Su composición es en un **93.7% compatible con nuestro universo**. Sin embargo, recomiendo precaución.
\end{quote}

**Manifestaciones físicas extrapoladas desde realidades alternativas** significaba que esos objetos eran cosas reales que habían sido tomadas de otros universos. Y el hecho de que fueran **93.7% compatibles con el nuestro** quería decir que eran casi iguales a lo que conocemos, pero con pequeñas diferencias, como la botella de OmniCola.

Julio apenas asintió, demasiado fascinado para prestar atención a la advertencia de Pi. ¿Una botella de soda de un universo donde Omniman existía de verdad? Su mente adolescente, que ya era brillante, no podía procesar tanta información. Era como si un sueño se hubiera hecho realidad.

Fue entonces cuando escuchó los pasos arrastrados de alguien bajando las escaleras. Se giró justo a tiempo para ver a su padre, Miguel, aparecer en el marco de la puerta de su habitación, despeinado y frotándose los ojos.

\begin{quote}
    —¿Qué haces despierto tan tarde? —preguntó Miguel, con un gran bostezo que le estiró la cara—. ¿Y qué es todo ese alboroto que haces a estas horas?
\end{quote}

Julio intentó esconder la botella detrás de su espalda, pero Miguel, aunque somnoliento, fue más rápido. Con una sonrisa adormilada en el rostro, le arrebató el refresco de las manos.

\begin{quote}
    —¿Una soda? ¡Gracias, hijo! —dijo alegremente, destapándola de un tirón con un "pop" que sonó muy fuerte en el silencio de la noche.

    —¡No, papá, no sabes de dónde viene! —exclamó Julio, intentando detenerlo, con los ojos bien abiertos por el pánico.
\end{quote}

Pero ya era tarde. Miguel tomó un largo trago, un sorbo ruidoso, chasqueó los labios con gusto y sonrió ampliamente, con los ojos cerrados de placer.

\begin{quote}
    —¡Wow! ¡Está buenísima! —dijo—. Sabe como si un relámpago y un caramelo se hubieran enamorado. ¡Necesito más de esto!
\end{quote}

Julio abrió la boca para replicar, para explicarle a su padre la increíble verdad, pero Pi intervino justo a tiempo, como siempre.

\begin{quote}
    —Análisis preliminar: el contenido de la bebida no presenta **toxicidad detectable**. Compuestos encontrados: agua, **azúcares modificados**, **electrolitos**, y trazas de un **elemento desconocido estabilizado**.
\end{quote}

**Toxicidad detectable** significaba que no había nada venenoso. **Azúcares modificados** y **electrolitos** (que son minerales como el sodio o el potasio que ayudan al cuerpo) eran cosas que ya existían, pero la **traza de un elemento desconocido estabilizado** era la parte que dejaba a Julio sin aliento: algo que no era de este mundo, pero que no hacía daño.

Miguel, completamente ajeno a la preocupación de su hijo y a la magnitud de lo que acababa de beber, se dejó caer felizmente en el viejo sillón del taller, cambiando de canal en el televisor con el control remoto mientras disfrutaba la OmniCola como si fuera lo más normal del mundo. Fox, en el suelo, miraba todo como si fuera parte de una broma privada entre él y el universo.

Julio soltó una risa ahogada, incapaz de contener la mezcla de nervios, alivio y la absurda situación. Había contactado con otra dimensión, había recibido objetos de otro mundo, y su padre estaba felizmente tomando un refresco interdimensional como si nada.

Cuando Miguel se quedó dormido en el sillón, con la botella vacía aún en la mano, Julio volvió su atención a la máquina.

El pequeño agujero negro, el "prompt", seguía allí, latiendo suave, como esperando.

En la pantalla de la terminal de Pi, un nuevo mensaje apareció, sin que Julio tuviera que escribir nada:

\begin{quote}
    "Esto es solo el principio, Julio."
\end{quote}

La aventura apenas comenzaba.

Y aunque un pequeño temor se aferraba a su pecho, una emoción aún más grande y poderosa lo impulsaba hacia adelante. El universo —o, mejor dicho, los universos— acababan de abrirle una puerta, y Julio, el chico que había construido una máquina para hablar con lo desconocido, estaba listo para cruzarla.

\cleardoublepage % Para que el capítulo empiece en una página impar

% Noveno capítulo
\chapter*{9 – La primera conversación}
\addcontentsline{toc}{chapter}{9 – La primera conversación} % Añadir al índice manualmente

El zumbido leve de la máquina se había vuelto una constante en el sótano de Julio, un sonido que lo acompañaba día y noche. Era como un latido metálico, suave pero persistente, que lo mantenía en un estado de alerta y fascinación permanentes. Ya no le parecía un ruido extraño, sino una parte de su propio universo, una melodía que significaba aventura.

Desde aquel increíble incidente con la **OmniCola** —la bebida que Miguel, su padre, se había bebido sin sospechar nada, pensando que era solo una gaseosa deliciosa—, Julio había pasado la semana entera pensando. Había analizado cada detalle de la esfera de cristal y los objetos que habían aparecido. La idea de que su máquina, que había empezado como un simple proyecto personal, fuera en realidad una puerta a otros mundos era algo que le llenaba de una mezcla de emoción y un poco de temor. Se había prometido ser cauto, y lo había sido. Pero la curiosidad lo carcomía por dentro.

Ahora que la conexión con el "**prompt**" estaba establecida y era estable —es decir, que tenía una pantalla en su computadora desde donde podía escribir y recibir mensajes de la anomalía sin necesidad de códigos binarios ni de activar la máquina con cada intento—, Julio y **Pi** decidieron intentar algo más audaz: establecer un canal de comunicación más directo y fluido. No era enviar un "Hola Mundo" al vacío, sino tener una charla de verdad.

La pantalla de su computadora, que Julio había conectado a la máquina, mostró un **parpadeo binario** al principio, esos interminables ceros y unos que eran el lenguaje de las computadoras. Pero luego, muy lentamente, como si la otra parte estuviera ajustándose, comenzaron a aparecer **patrones más complejos**: líneas de datos que Pi, con su avanzada capacidad, interpretaba en tiempo real y traducía a palabras que Julio podía leer. La expectación era palpable en el aire, casi se podía tocar.

\begin{quote}
    —¿Ves esto, Pi? —susurró Julio, su voz apenas audible, como si no quisiera que el resto del mundo, ni siquiera el aire, se enterara de su pequeño y monumental secreto.

    —Lo veo, Julio —respondió el asistente virtual, su voz sintética con un matiz que, si Pi tuviera emociones, sería de sorpresa contenida—. Es una **estructura de lenguaje**. Aún rudimentaria, es decir, básica, pero claramente intencionada. Parece que son **respuestas automatizadas**, mensajes pregrabados esperando a ser activados.
\end{quote}

Julio tragó saliva. Su corazón latía tan fuerte que casi no escuchaba el zumbido de fondo de la máquina. La mano le temblaba un poco al acercarse al teclado. Era un momento crucial. Respiró hondo y, con los dedos temblorosos, tipeó ansioso en la consola improvisada que Pi había creado:

\begin{quote}
    "¿Quién eres?"
\end{quote}

Por unos segundos, que parecieron horas, no pasó nada. La pantalla se mantuvo en silencio, y Julio empezó a creer que todo era solo una ilusión, que su mente le estaba jugando una mala pasada. Pero entonces, letra por letra, con una lentitud dramática que aumentaba el suspenso, una respuesta apareció en la pantalla, formándose como si una entidad invisible estuviera escribiendo justo para él:

\begin{quote}
    "ALGUIEN QUE VIGILA. ALGUIEN QUE SABE QUIÉN ERES."
\end{quote}

Julio se estremeció. La piel se le puso de gallina. Aquella respuesta lo golpeó como una ola de agua fría. ¿Cómo podía ese "algo" saber su nombre? ¿Había estado observándolo todo este tiempo? Miró a Fox, que dormitaba enroscado junto a una maraña de cables, su patita moviéndose ligeramente en un sueño. Ni siquiera el gato parecía captar la magnitud del momento, el hecho de que había una inteligencia del otro lado que lo conocía.

Pi, con su voz inalterable y carente de emociones humanas, solo comentó, interrumpiendo el torbellino de pensamientos de Julio:

\begin{quote}
    —El mensaje contiene un **patrón de seguridad**, Julio. No parece hostil. Es como una clave, un identificador, pero no una amenaza.
\end{quote}

Julio decidió continuar, a pesar del nudo en su estómago. La curiosidad era más fuerte que el temor. Quería saber más. Tenía que saberlo.

\begin{quote}
    "¿Qué quieres?"
\end{quote}
preguntó esta vez, con los dedos aún temblorosos sobre el teclado.

La respuesta fue casi inmediata, apareciendo con una velocidad que antes no había demostrado:

\begin{quote}
    "AYUDARTE. PERO HAY RIESGOS."
\end{quote}

**Ayudar**. ¿Ayudar con qué? La palabra resonó en la mente de Julio. Ayudar. Eso era inesperado. Julio miró el pequeño portal pulsante en el centro de la máquina, ese punto oscuro que había aprendido a respetar y a ver como una puerta. Y al leer ese mensaje, sintió que se abría ante él un abismo de posibilidades, pero también de peligros. Era como estar al borde de un precipicio, viendo un camino que prometía cosas increíbles, pero también escondía grandes incógnitas.

\begin{quote}
    —¿Qué opinas, Pi? —preguntó, su voz era un susurro, buscando la lógica de su asistente, su guía.

    —Se recomienda **precaución** —dijo Pi, siempre prudente—. Pero la oportunidad de aprendizaje es… considerable.
\end{quote}

Una sonrisa lenta, casi imperceptible al principio, se dibujó en el rostro de Julio. La palabra "considerable" en el lenguaje de Pi era una luz verde para la aventura. Era un riesgo, sí, pero un riesgo que valía la pena tomar.

Julio tipeó en la consola, a modo de broma, pero con un fondo de verdad en su corazón:

\begin{quote}
    "¿Me enseñas a controlar la máquina, maestro desconocido?"
\end{quote}

La respuesta llegó casi al instante, con una formalidad y una solemnidad que hicieron que la piel de Julio se erizara:

\begin{quote}
    "TE ENSEÑARÉ. PERO CADA COMANDO, CADA DESEO, ABRIRÁ PUERTAS QUE NO PUEDES CERRAR."
\end{quote}

La frase se quedó flotando en el aire de la habitación, más allá de la pantalla, como un presagio. Era una advertencia clara: lo que estaba a punto de aprender, las puertas que abriría, serían irreversibles. No habría vuelta atrás. La magnitud de esa verdad lo dejó sin aliento por un momento.

---

Antes de que pudiera reflexionar más sobre el significado de esas palabras, la consola de su computadora empezó a parpadear erráticamente, de forma descontrolada. El brillo intermitente se hizo más intenso, como si la pantalla estuviera sufriendo una sobrecarga.

\begin{quote}
    —Pi, ¿qué pasa? —preguntó Julio, la voz llena de alarma.

    —Detecto **actividad no autorizada** —dijo Pi, su voz, por primera vez, con un toque de urgencia—. Algo intenta atravesar el portal… no es una respuesta a nuestra conversación. No fue solicitado. Es un **intento de entrada no deseado**.
\end{quote}

Julio retrocedió un paso, alejándose de la máquina. Un escalofrío helado le recorrió la espalda. Esto no era parte del plan. Fox, que estaba profundamente dormido, se despertó de su siesta de golpe. Miró hacia la máquina con un maullido bajo y desconfiado, el pelo de su lomo erizado, como si también sintiera la amenaza. Sus ojos verdes estaban fijos en el punto negro, que ahora parecía latir con más fuerza, su brillo era más intenso.

\begin{quote}
    —¿Debo apagarlo? —preguntó Julio, dudando. Su mano se acercó al botón de apagado de emergencia, pero se detuvo.

    —No hay garantía de que el cierre manual detenga la transferencia —advirtió Pi, su voz era más rápida, llena de información crucial—. Podría causar una **ruptura violenta en el flujo de energía**, lo que podría dañar la máquina o… quién sabe qué más.
\end{quote}

El **parpadeo rojo** en la pantalla de Pi, que indicaba una alerta crítica, aumentó su frecuencia, como un latido nervioso y acelerado, una advertencia urgente. Julio cerró los ojos por un segundo, sintiendo el pánico apoderarse de él. Respiró hondo, intentando calmar su mente, y luego, con determinación, puso una mano sobre la consola, sobre el teclado.

Sabía que no había marcha atrás. La puerta ya estaba abierta, y algo, o alguien, estaba intentando cruzar. Julio estaba a punto de aprender que la curiosidad, a veces, viene con un precio muy alto. El futuro, una vez una promesa lejana, ahora era un torbellino de eventos que se acercaba a toda velocidad.

\cleardoublepage % Para que el capítulo empiece en una página impar

% Décimo capítulo
\chapter*{10 – Algo ha cruzado}
\addcontentsline{toc}{chapter}{10 – Algo ha cruzado} % Añadir al índice manualmente

La habitación de Julio se mantuvo en un inquietante silencio después del súbito destello de luz que había invadido el sótano. Era un silencio denso, cargado de expectativas y nerviosismo, casi como si el aire mismo contuviera la respiración. Julio, con los ojos bien abiertos como platos, miraba fijamente el centro del portal, donde una figura pequeña y desconocida descansaba ahora en el suelo, justo en la **bandeja de recolección**, esa base metálica que había diseñado para recoger los objetos que pudieran aparecer dentro de la caja de microondas. Su corazón palpitaba con tanta fuerza que podía sentirlo retumbando en sus oídos, un tamborileo constante que era la única melodía en ese momento de asombro. Cada segundo de ese silencio era una eternidad.

Fox, que había estado dormitando plácidamente en su rincón habitual junto a los cables, se incorporó con lentitud, sus músculos estirándose con una pereza inicial que pronto se transformó en alerta. Soltó un maullido bajo y desconfiado, una especie de advertencia suave que solo Julio parecía entender. Sus orejas se echaron hacia atrás, pegadas a la cabeza, una clara señal de cautela felina, y sus ojos celestes, normalmente soñolientos, estaban fijos en el objeto que había aparecido de la nada. El gato parecía entender, a su manera, que algo importante había cambiado.

\begin{quote}
    —¿Qué demonios fue eso, Pi? —susurró Julio, tragando saliva con dificultad. La voz le salió apenas un hilo, un susurro ahogado por la inmensidad del momento. No era un simple objeto; era algo que había cruzado.

    —Algo ha cruzado —respondió Pi, su voz, aunque tranquila y robótica como siempre, esta vez sonaba más atenta, casi como si midiera cada palabra, dándole una importancia extra al momento. No había alegría ni sorpresa en su tono, solo una afirmación directa de un hecho extraordinario. Su análisis era preciso, sin emociones.
\end{quote}

Julio se acercó despacio, con cautela, sus pasos lentos y decididos, como si se acercara a un animal salvaje o a un tesoro delicado. El objeto no se movía. Era una especie de **caja, rectangular**, pero con bordes suaves, más grande que sus dos manos juntas, lo suficiente como para sostenerla con ambas. Estaba hecha de un **material gris oscuro** que parecía absorber la luz en lugar de reflejarla, dándole un aspecto casi mate, como si estuviera hecha de sombras solidificadas. Su superficie no era lisa; estaba cubierta de **grabados finos y complejos**, como si fueran circuitos microscópicos o símbolos alienígenas, líneas intrincadas que se conectaban entre sí. Y lo más fascinante era que esos grabados pulsaban suavemente, con un brillo interno tenue, casi imperceptible, como si la caja respirara con una vida propia. Era extraño, diferente a todo lo que había visto, pero a la vez, extrañamente familiar en su forma.

\begin{quote}
    —¿Detectas vida? —preguntó Julio en voz baja, con el aliento contenido. La idea de que pudiera ser un ser vivo, aunque fuera de otro universo, era abrumadora. Era la primera pregunta que venía a su mente al ver algo tan orgánico en su movimiento.

    —Negativo —respondió Pi de inmediato, con la eficiencia que lo caracterizaba—. No hay signos biológicos. Es inerte. No emite radiación peligrosa, lo cual es una excelente noticia. Su estructura parece ser de un **compuesto aún no catalogado**, lo que significa que no se parece a nada que tengamos registrado en nuestra ciencia terrestre. Es un material completamente nuevo para nuestras bases de datos.
\end{quote}

Julio respiró aliviado, soltando el aire que no sabía que había estado conteniendo en sus pulmones. Al menos no era un monstruo interdimensional dispuesto a devorarlo o a traer una plaga de otro universo. Esa era una muy buena noticia, una preocupación menos en su ya complicada lista de pensamientos. La posibilidad de un peligro inmediato se disipaba.

Fox, siempre curioso, y con la audacia que solo los gatos tienen, se acercó cautelosamente a olfatear el objeto. Su nariz rosada se movió, husmeando con atención cada centímetro de la caja. Maulló dos veces, un sonido grave y pensativo, como si estuviera reflexionando sobre su origen. Dio un par de vueltas alrededor de la caja misteriosa, rozándola con su pelaje, y, decidiendo finalmente que no representaba una amenaza inmediata, se tumbó tranquilamente a su lado, gordo y satisfecho, como si la nueva caja fuera un cómodo cojín o un nuevo mueble en la habitación. La tranquilidad del gato era una buena señal para Julio.

\begin{quote}
    —Bien, Fox —dijo Julio, soltando una risa nerviosa al ver la reacción de su gato—. Si a ti no te importa, supongo que debe ser seguro para mí también. Eres mi detector de peligros personal.
\end{quote}

Pi, mientras tanto, no se quedó quieto. Sus luces internas destellaban y cambiaban de color a una velocidad asombrosa, como si estuviera pensando muy rápido, intentando comunicarse, o al menos entender, la nueva presencia. Estaba analizando el módulo en profundidad, escaneando cada detalle de su superficie, su composición química y sus emisiones de energía. Julio podía ver los datos aparecer en la pantalla de Pi, aunque no los entendía del todo.

\begin{quote}
    —Julio —dijo Pi, su voz era más clara ahora, llena de un nuevo descubrimiento, con un matiz de fascinación que rara vez mostraba—. Este objeto parece diseñado para interactuar directamente con nuestra máquina. Está emitiendo un **patrón de pulsos compatible con nuestro sistema de energía**, lo que significa que "habla" el mismo idioma que la máquina que construiste. Es como si estuviera hecho para encajar aquí.
\end{quote}

Julio se pasó una mano por el cabello, aún procesando la increíble situación. Su mente daba vueltas a mil por hora.

\begin{quote}
    —¿Quieres decir que… es una especie de accesorio? ¿Como un USB que se conecta para añadir una función?

    —Más bien una **llave** —corrigió Pi con precisión—. O una **actualización**. Es algo que permite que la máquina evolucione o que acceda a nuevas funciones.
\end{quote}

La idea le provocó un escalofrío que le recorrió toda la columna vertebral. No era solo que la máquina funcionara y trajera objetos curiosos. Era mucho más. Alguien, o algo, no solo sabía que él había creado el portal, que había abierto esa puerta a lo desconocido, sino que ahora le estaba enviando las piezas necesarias para que esa máquina evolucionara, para que fuera aún más poderosa, más compleja. Era como si un gran arquitecto cósmico, una inteligencia superior, le estuviera enviando los planos de su próximo paso, sin que él tuviera que pedirlos. La idea era emocionante y aterradora al mismo tiempo. Era un juego de ajedrez donde el otro jugador ya conocía todos sus movimientos.

---

Se sentó en el suelo frío del sótano, frente al objeto, observándolo con una mezcla de respeto y asombro, como si fuera el enigma más grande que jamás hubiera visto. La emoción palpitaba con fuerza en su pecho, un latido que le decía "¡adelante, descubre más!", pero también sentía el miedo, una punzada fría que le advertía de lo desconocido, de los peligros que podrían estar más allá. Cada nueva pieza que cruzaba desde "el otro lado" significaba un paso más hacia lo incierto, hacia un futuro que no podía prever, pero que, a su vez, lo llamaba con una fuerza irresistible. Y aunque su curiosidad natural, esa fuerza imparable que lo impulsaba desde niño a desarmar y entender todo, lo empujaba hacia adelante, no podía ignorar la creciente sensación de que estaba participando en algo mucho más grande que él mismo, algo que superaba su entendimiento y que lo convertía en una pieza clave de un misterio cósmico.

\begin{quote}
    —¿Qué sugieres, Pi? —preguntó finalmente Julio, buscando la lógica y el consejo de su fiel asistente. La decisión no era fácil; sentía el peso de la responsabilidad sobre sus hombros.

    —Podríamos intentar conectarlo —dijo Pi, analizando las posibilidades en fracciones de segundo—. Sin embargo, existe un porcentaje de **riesgo significativo**. No conocemos aún todas sus funciones ni su propósito final. No tenemos todos los datos necesarios para una conexión completamente segura. Podría haber **consecuencias imprevistas**.
\end{quote}

Julio se quedó pensando. La idea de un riesgo "significativo" lo hizo detenerse de verdad. Por primera vez en mucho tiempo, consideró seriamente detenerse por completo, aunque fuera sólo por una noche, para analizar la situación con más calma. Ya no se trataba de un juego o un simple experimento científico de la escuela. Se trataba de algo mucho mayor, con implicaciones que no podía ni siquiera imaginar.

\begin{quote}
    —¿Y si lo estudiamos primero? —dijo Julio, con una idea que le dio una sensación de control, una pausa necesaria—. ¿Ver qué intenta decirnos o para qué sirve antes de conectarlo a la máquina? Si es una llave, ¿qué abre? Necesitamos saber más antes de abrir otra puerta.
\end{quote}

Pi, como siempre, procesó su idea en fracciones de segundo. La respuesta fue inmediata y clara:

\begin{quote}
    —Una decisión **prudente**, Julio. Altamente recomendada.
\end{quote}

Él sonrió débilmente. Pi raramente usaba la palabra "prudente" con tanta fuerza, y que la acompañara con "altamente recomendada" era una clara señal de que Julio estaba pensando correctamente. Miró a Fox, que ahora roncaba suavemente al lado del misterioso objeto, su pelaje blanco moviéndose al ritmo de su respiración, como si todo fuera parte de un día normal y aburrido en el sótano. La tranquilidad del gato era un bálsamo para sus nervios, una pequeña ancla a la realidad.

\begin{quote}
    —Parece que tenemos trabajo por hacer —murmuró Julio, sintiendo cómo la adrenalina se mezclaba con la determinación. El cansancio de la noche anterior se desvanecía ante la nueva tarea.
\end{quote}

Se acomodó frente a su computadora, sus dedos volando sobre el teclado con renovada energía. Abrió una nueva sección de su proyecto, un archivo completamente nuevo y vital, y tituló la carpeta con cuidado, casi como un científico profesional que inicia una investigación de alto nivel:

\begin{quote}
    "**Estudio del Módulo Interdimensional - Fase 1**".
\end{quote}

Mientras tanto, la máquina zumbaba en silencio detrás de él, con sus luces tenues y constantes. Parecía esperar… con una paciencia inhumana, como si el tiempo no existiera para ella, o como si supiera que, tarde o temprano, Julio tomaría la decisión correcta. La aventura estaba a punto de volverse mucho más profunda, y Julio, el chico que había construido una puerta a lo desconocido, estaba listo para desentrañar los secretos que acababan de cruzar el umbral de su realidad.

\cleardoublepage % Para que el capítulo empiece en una página impar

% Undécimo capítulo
\chapter*{11 – Planos y Problemas}
\addcontentsline{toc}{chapter}{11 – Planos y Problemas} % Añadir al índice manualmente

El pequeño **módulo gris oscuro**, con sus grabados que pulsaban suavemente como un corazón alienígena, se convirtió en el nuevo centro del universo para Julio. Lo estudió durante días, con la meticulosidad de un arqueólogo desenterrando un tesoro antiguo. Desde el momento en que lo había puesto en la **bandeja de recolección** de la máquina, no había vuelto a activarla. La promesa de que "algo había cruzado" y la necesidad de prudencia lo habían llevado a apagarla, dejando el sótano en un silencio inusual, solo roto por el suave zumbido de los servidores de **Pi**. Julio sentía que si iba a seguir avanzando en este misterio, tenía que hacerlo de la manera más segura y correcta posible.

La primera tarea de Julio fue intentar entender el módulo a fondo. Con la ayuda inestimable de Pi, escaneó su superficie con láseres de baja potencia, analizó su composición con **espectrómetros improvisados** (aparatos que detectan de qué están hechas las cosas, como si fuera una radiografía profunda de materiales), y estudió los patrones de energía que emitía. Pi confirmó lo que había sospechado desde el principio: el objeto estaba diseñado específicamente para mejorar su máquina actual. Era como una pieza de un rompecabezas mucho más grande, creada para hacer que la máquina de Julio fuera aún más potente y capaz, desbloqueando nuevas funciones que ni siquiera podía imaginar. Sin embargo, no era un simple "enchufar y listo". Había un paso intermedio crucial.

Julio descubrió esto al intentar hacer la primera conexión de prueba. Conectó algunos cables de prueba delicados a los puntos de contacto del módulo, con la esperanza de que la máquina lo reconociera y se integrara de forma automática. No fue un fallo total, afortunadamente. No hubo chispas, ni humo, ni apagones repentinos. Pero tampoco hubo un acoplamiento perfecto, una conexión instantánea. En lugar de eso, al hacer el contacto, se desprendió una cantidad increíble de información, una **avalancha de datos** que Pi apenas pudo procesar. Era como si el módulo, al sentir la conexión, hubiera "descargado" todos sus secretos directamente a la computadora de Julio.

Lo que apareció en la pantalla de Pi lo dejó con la boca abierta y el corazón latiéndole a mil. No eran solo datos crudos. Eran **planos y esquemas detallados de un circuito** que se superponían exactamente sobre el diseño de su propia máquina. Era una réplica exacta de sus invenciones, pero con agregados. Era como si la persona o entidad que le había enviado el módulo hubiera tenido una copia fiel de su libro de notas, de cada garabato, cada diagrama y cada idea que Julio había plasmado en papel desde el inicio de su proyecto. Podía ver su propio diseño de la máquina, pero modificado y mejorado, con componentes que él jamás hubiera imaginado. Luego, en una sección aparte, había una serie de **pasos e instrucciones muy claras y precisas** para modificar completamente la máquina, indicando dónde soldar nuevos componentes con una precisión quirúrgica, cómo reubicar otros elementos y qué nuevas piezas eran necesarias para integrar el módulo interdimensional de forma correcta.

Esto, sin embargo, era un problema. Un gran, gran problema. Para realizar estas modificaciones, Julio no solo necesitaba conseguir más **repuestos electrónicos avanzados** —componentes raros, difíciles de encontrar y, lo más importante, increíblemente caros—, sino que también requería algo esencial y muy escaso para un adolescente: **dinero**. Mucho dinero. Una cantidad bastante elevada para un niño de su edad, que ahora tenía catorce años. Era una cifra que lo dejó pensativo por horas.

Obviamente, robar no era una opción. Era algo que ni siquiera cruzó por su mente. Su moral se lo impedía. Pedir prestado a amigos tampoco era una solución, porque nadie de su edad tendría la cantidad de dinero que necesitaba. Y mucho menos pedirle el dinero a sus padres. Sabía que ellos ya estaban preocupados por sus hábitos de sueño y su salud por pasar tanto tiempo en el sótano. Si les decía que necesitaba una suma tan grande para su "proyecto", lo más probable es que se lo prohibieran por completo, o peor, que lo internaran en un campamento de verano para "desconectarse de la tecnología" y "jugar al aire libre".

Julio necesitaba un plan. Un plan para conseguir ese dinero por sí mismo, de manera honesta y sin levantar sospechas. Después de días de pensar, dibujando estrategias en su cuaderno y borrando ideas fallidas, la única opción lógica que se le ocurrió fue conseguir un **trabajo de medio tiempo**. Pero ¿dónde? ¿Quién contrataría a un chico de catorce años con la cabeza en las nubes y las manos llenas de soldadura, que pasaba más tiempo en el sótano que en la calle?

Antes de lograr conseguir un empleo formal, Julio intentó de todo para ganar unas pocas monedas. Cortó el césped de los jardines de sus vecinos, sudando bajo el sol, ofreciendo sus servicios con una sonrisa forzada y la esperanza de juntar algo. Hizo mandados para su padre, llevando herramientas o buscando materiales para sus proyectos, siempre con la esperanza de recibir una pequeña paga extra, una "propina" que, sumada, apenas alcanzaba. Incluso ofreció ayuda para organizar garajes y trasteros, pero el dinero que lograba juntar era mínimo, apenas alcanzaba para comprar dulces en la tienda de la esquina, y estaba muy, muy lejos de lo que necesitaba para los complejos componentes electrónicos que requerían los nuevos planos.

La frustración comenzaba a crecer en su interior. Pasaban los días y las semanas, y el dinero no aparecía a la velocidad que deseaba. Sus ahorros, que tenía guardados de algunos regalos de cumpleaños y propinas esporádicas, eran insignificantes para la magnitud de lo que se venía. La máquina, cubierta con una lona en el sótano, permanecía apagada, como una bestia dormida que esperaba ser despertada.

Un día, mientras comía con su madre Ana, le comentó su dilema, sin darle muchos detalles de "para qué" necesitaba el dinero, solo diciendo que era para un proyecto personal. Ana, siempre atenta a los intereses de sus hijos y con una sabiduría práctica, le dio una idea.

\begin{quote}
    —Julio, sé de una tienda aquí en el barrio, a pocas casas de la nuestra, "**TecnoFix**", donde reparan teléfonos y computadoras —dijo Ana, mientras recogía los platos de la mesa—. El dueño, **Don Ramiro**, es amigo de tu tío. Quizás podrías ir a preguntar si necesitan un ayudante. Eres bueno con esas cosas, ¿no? Siempre arreglas los electrodomésticos viejos de la abuela y los de la casa. Además, les tienes paciencia.
\end{quote}

La idea de Ana fue como una chispa que encendió la esperanza en el pecho de Julio. Aunque escéptico de que lo fueran a contratar en un lugar tan especializado, decidió intentarlo. Se puso su mejor ropa (la que estaba menos arrugada de su armario), peinó su cabello con esmero y se dirigió a "TecnoFix" al día siguiente, con el corazón latiéndole fuerte contra sus costillas. El camino era corto, apenas cinco minutos caminando por su misma cuadra.

Para su sorpresa, el dueño de la tienda, Don Ramiro, un hombre de aspecto serio pero con una sonrisa amable, aceptó darle una oportunidad. No era fácil conseguir un trabajo allí; Don Ramiro era conocido por ser demasiado exigente en las personas que contrataba, buscando siempre a los más hábiles, dedicados y responsables. Pero al ver la chispa de curiosidad en los ojos de Julio, su conocimiento básico de electrónica que demostró con unas cuantas preguntas, y su genuina disposición a aprender, decidió darle una oportunidad como **ayudante de medio tiempo**, al final de sus clases.

El trabajo en "TecnoFix" resultó ser una verdadera mina de oro para Julio, no solo por el dinero que empezaría a ganar, sino también por otras cosas aún más valiosas. Julio se dedicaba a desarmar teléfonos y computadoras viejas que llegaban para chatarra, a clasificar componentes y a ayudar en reparaciones básicas, como cambiar pantallas o baterías. Lo mejor de todo era que, a menudo, los dueños de los aparatos no querían las piezas que ya no servían, y el dueño de la tienda, Don Ramiro, le permitía quedarse con ellas. Esos repuestos inservibles para los dueños, desechos para otros, eran tesoros para Julio. Podía estudiar cómo estaban hechos los circuitos complejos, qué tipo de cables usaban, y a veces, incluso, encontrar alguna pieza en buen estado que pudiera servirle para su gigantesco proyecto.

Pero no solo encontró repuestos. En "TecnoFix", Julio también hizo un nuevo amigo: **Sebastián**. Sebastián, al que todos llamaban "**Sebas**", era un chico de su misma edad, catorce años. No trabajaba en la tienda de repuestos, pero era el sobrino de Don Ramiro y pasaba por allí de vez en cuando después de sus entrenamientos. Sebas era el típico deportista, atlético y con una energía desbordante, pero, muy diferente a los chicos de su colegio, que a menudo se burlaban de los intereses "raros" de Julio en la ciencia y la tecnología. Vivía en el mismo manzano que Julio, solo a un par de casas de la tienda y a la vuelta de la esquina de la de Julio, lo que hacía que sus encuentros fueran aún más frecuentes y naturales.

A diferencia de los otros chicos, Sebas encontraba fascinante que un chico de su edad fuera capaz de soldar componentes eléctricos con tanta destreza, y que, además, hubiera sido capaz de conseguir un trabajo de medio tiempo en la tienda de su tío, que era tan exigente con su personal. Sebas se quedaba a menudo observando a Julio mientras desarmaba un teléfono, haciendo preguntas curiosas y genuinas.

\begin{quote}
    —¿Y esto para qué sirve exactamente? —preguntaba Sebas, señalando una placa de circuito llena de chips diminutos.

    —Es el cerebro del teléfono —explicaba Julio, mostrando su paciencia, algo poco común en él cuando no hablaba de su máquina—. Aquí es donde se guardan los datos, donde se hacen todos los cálculos. Si esto se daña, el teléfono ya no piensa.
\end{quote}

Julio pasó los siguientes meses recaudando dinero, trabajando en "TecnoFix" con una dedicación que sorprendía a su propio jefe. Cada moneda que ganaba era guardada con cuidado en una **alcancía especial** en su cuarto, cada componente reciclado era clasificado y almacenado meticulosamente en cajas en su sótano. La máquina, cubierta con una lona, permanecía apagada. Julio había tomado la decisión consciente de apagarla mientras lograba sus objetivos. Sabía que la paciencia era clave para el éxito, y que no podía arriesgarse a activar algo tan poderoso sin estar completamente preparado para las consecuencias. El **módulo interdimensional**, por su parte, descansaba en una caja fuerte improvisada, esperando el momento de ser integrado.

Durante esos meses, Sebastián solo pasaba esporádicamente por la tienda, entre sus entrenamientos de fútbol o básquet y sus partidos. Al principio, sus visitas eran breves, solo para saludar a su tío o recoger algo. Pero poco a poco, su curiosidad sobre Julio incrementaba. Empezó a quedarse más tiempo, a hacer más preguntas sobre electrónica, e incluso a ayudarle a clasificar piezas pequeñas o a pasarle herramientas.

La amistad entre ellos se hizo más fuerte con cada día que pasaba. Al principio, Julio era reservado, como siempre lo había sido con la mayoría de la gente. Pero a medida que pasaban las semanas, y al ver el genuino interés y la falta de prejuicios de Sebas, Julio comenzó a abrirse a él, compartiendo pequeñas anécdotas de sus reparaciones más complicadas y, con el tiempo, incluso algunos de sus sueños más ambiciosos, aunque sin revelar aún el gran secreto de su máquina interdimensional. Sebas, a su vez, compartía sus experiencias deportivas, enseñándole a Julio la disciplina de un atleta y la importancia del trabajo en equipo, algo que Julio nunca había considerado para sus proyectos solitarios.

Julio, el genio solitario, estaba construyendo no solo una máquina interdimensional, sino también una **amistad verdadera** que prometía cambiar su vida de maneras que aún no podía imaginar.

\cleardoublepage % Para que el capítulo empiece en una página impar

% Duodécimo capítulo
\chapter*{12 – Objetivo Cumplido con Resultados Inesperados}
\addcontentsline{toc}{chapter}{12 – Objetivo Cumplido con Resultados Inesperados} % Añadir al índice manualmente

Seis meses. Seis largos meses pasaron volando, llenos de trabajo, estudio y paciencia. Julio, con sus catorce años, había estado dedicado en cuerpo y alma a su misión. Cada tarde, después del colegio, caminaba las pocas cuadras hasta "**TecnoFix**", la tienda de Don Ramiro en su mismo barrio. Desarmaba teléfonos, clasificaba componentes, ayudaba a su jefe, y guardaba cada centavo que ganaba. Su **alcancía**, que al principio parecía vacía, ahora sonaba pesada. Las piezas "inservibles" que Don Ramiro le permitía llevar a casa se acumulaban en su sótano, dándoles una segunda vida en el meticuloso laboratorio improvisado de Julio.

Un martes por la tarde, justo cuando el sol empezaba a teñir el cielo de naranja sobre los tejados de su barrio, Don Ramiro llamó a Julio a la caja. Tenía una sonrisa más amplia de lo habitual.

\begin{quote}
    —Julio, has superado todas mis expectativas —dijo Don Ramiro, entregándole un sobre más grueso de lo normal—. Eres el ayudante más trabajador y dedicado que he tenido. Aquí tienes tu último pago, más un pequeño bono. Considera esto una parte del futuro que estás construyendo.
\end{quote}

Julio sintió una oleada de emoción. El sobre era grueso, pesado. Sabía lo que significaba. ¡Había logrado su objetivo! Había juntado todo el dinero que necesitaba. Salió de la tienda con una sonrisa que no cabía en su rostro, casi corriendo a casa.

Cuando llegó, encontró a su madre, Ana, en la cocina preparando la cena. Con el corazón latiéndole fuerte, le entregó el sobre sin decir una palabra. Ana lo abrió, sus ojos se abrieron de sorpresa al ver la cantidad.

\begin{quote}
    —¡Julio! ¿Todo esto? —exclamó, con una mezcla de asombro y orgullo. Ella sabía de su empeño en el trabajo, pero ver el resultado tangible de su esfuerzo era diferente.

    —Es para mi proyecto, mamá —dijo Julio, sus ojos brillando—. Y lo conseguí yo solo.
\end{quote}

Ana lo abrazó con fuerza, sus brazos rodeándolo en un abrazo cálido y emotivo. —Estoy tan orgullosa de ti, mi cielo. No por el dinero, sino por tu esfuerzo y tu dedicación. Siempre persigues lo que quieres, y eso es algo admirable.

Esa noche, Julio no duró. Pasó horas en su computadora, haciendo los pedidos en línea. Con la lista de componentes de los planos que el **módulo interdimensional** le había dado, eligió cuidadosamente cada pieza. Eran nombres que la mayoría de la gente no entendería:

\begin{itemize}
    \item Unidades de procesamiento cuántico (o al menos lo más parecido a eso que encontró).
    \item Capacitadores de flujo molecular de alta densidad.
    \item Bobinas de inducción hiper-resonantes.
    \item Paneles de conducción de energía de grafeno.
    \item Módulos de interfaz neuro-óptica.
\end{itemize}

Eran componentes de alta gama, diseñados para trabajos muy específicos, pero Julio sabía exactamente cómo encajarían en su máquina. Confirmó el pedido, que se enviaría desde varias partes del mundo. Ahora solo le quedaba esperar.

Mientras los paquetes viajaban por el mundo, Julio decidió celebrar. Había prometido a **Sebas** una noche de videojuegos, y era el momento perfecto. Invitó a Sebas y también a su hermano menor, **Óscar**, al que le encantaban los dulces y la comida.

La casa de Julio se llenó de risas y el sonido de los videojuegos. Óscar, fiel a su estilo, trajo una mochila llena de papas fritas, chocolates y gaseosas. Julio, con parte del dinero sobrante de su sueldo, había comprado aún más golosinas y refrescos para la ocasión. Jugaron hasta que les dolieron los pulgares. Las pizzas volaron, las risas rellenaron el ambiente, y los gritos de victoria y derrota se escuchaban por todo el cuarto.

Óscar, después de comer la mitad de las botanas y caramelos y beber varias gaseosas, se quedó dormido en el sofá, envuelto en una manta, con la boca ligeramente abierta y una sonrisa soñadora. Julio sonrió, ya que sabía que Óscar no recordaría mucho de la noche, lo cual era perfecto para su secreto.

A medida que la noche avanzaba, la energía de los juegos disminuía, y solo quedaron Julio y Sebastián, compartiendo un silencio cómodo mientras el control de la consola descansaba en sus manos. Sebas, que solía ser tan enérgico, ahora estaba relajado. La conversación se tornó más personal.

\begin{quote}
    —Nunca pensé que alguien como tú podría trabajar en una tienda de electrónica —dijo Sebas, riendo suavemente—. Eres muy bueno. Mi tío dice que tienes un talento increíble.
\end{quote}

Julio sonrió. —Gracias, Sebas. La verdad es que me gusta mucho lo que hago.

Se hizo un pequeño silencio. Julio miró a Sebas, a su amigo que lo había aceptado sin juzgarlo, que se había interesado genuinamente en sus pasiones, a pesar de ser tan diferentes. Sebas siempre estaba allí, escuchando y mostrando un interés que nadie más tenía. Era el momento. Tenía que contárselo.

\begin{quote}
    —Hay algo que tienes que saber, Sebas —dijo Julio, su voz baja y seria. Sebas lo miró, notando el cambio de tono.
\end{quote}

Julio se levantó y lo guio hacia el sótano. Sebas lo siguió, curioso. Al encender la luz, el sótano, normalmente lleno de herramientas y componentes, parecía aún más misterioso con la máquina cubierta por la lona.

\begin{quote}
    —¿Recuerdas que te dije que tenía un proyecto grande, verdad? —Julio se detuvo frente a la lona.

    —Sí, algo de eso me contaste. ¿Es esto?
\end{quote}

Julio tomó una respiración profunda. —Sí. Es mucho más que un proyecto, Sebas. Es… es una máquina. Una máquina que puede… conectar con otros universos.

Sebas parpadeó, luego soltó una risa nerviosa. —Jajaja, ¿otros universos? ¿Como en las películas de ciencia ficción?

Julio no respondió. En lugar de eso, se acercó a su escritorio y buscó algo. Luego, tomó dos pequeñas **esferas de cristal**, no más grandes que canicas. Cada una contenía palabras y símbolos que giraban lentamente en su interior, brillando con una luz suave. Eran las esferas de "Hola Mundo" que había recuperado de los primeros experimentos.

\begin{quote}
    —Mira esto —dijo Julio, entregándole una a Sebas.

    —¡Wow! ¿Qué es esto? Es increíble.
\end{quote}

Julio le explicó brevemente sobre la anomalía, el "prompt", y cómo la máquina lo había llamado por su nombre. Sebas escuchaba, cada vez más serio, el brillo de asombro en sus ojos aumentando.

\begin{quote}
    —Y… esto que logré hacer es por eso que necesitaba el dinero —dijo Julio, señalando el módulo interdimensional, aún guardado en su caja fuerte improvisada—. Esto es una actualización. Pero para instalarla, necesito modificar la máquina. Y esos son los planos que me dio para hacerlo.
\end{quote}

Sebas miró de la esfera a Julio, luego a la lona que cubría la máquina. Su mente deportista, acostumbrada a lo tangible, estaba procesando algo totalmente nuevo.

\begin{quote}
    —¿Y tú… tú lo hiciste? ¿Tú creaste esto? —preguntó Sebas, con la voz apenas audible.
\end{quote}

En ese momento, Julio encendió la computadora. La pantalla se iluminó.

\begin{quote}
    —Pi, preséntate a Sebastián —dijo Julio.
\end{quote}

Una voz sintética, pero con un tono amigable, resonó en el sótano:

\begin{quote}
    —Hola, Sebastián. Soy **Pi**, la inteligencia artificial de Julio. Es un placer conocerte.
\end{quote}

Sebas dio un salto hacia atrás, su rostro se puso blanco. —¡Ah! ¿Una… una computadora que habla? ¡Esto es una locura!

Julio sonrió. —No es una computadora, Sebas. Es una **inteligencia artificial**. Y sí, esto es real. Por eso he estado tan ocupado, por eso he trabajado tanto. No podía contárselo a nadie, pero tú… tú eres mi mejor amigo, y necesitaba que lo supieras.

Se dirigió a la máquina, levantó la lona con cuidado, revelando la compleja estructura de metal, cables y el corazón de microondas. Conectó rápidamente un par de cables y presionó un botón. La máquina zumbó, las luces tenues parpadearon.

\begin{quote}
    —Pi, un par de **OmniColas**, por favor —pidió Julio, mirando a Sebas.
\end{quote}

La máquina vibró. El punto negro en el centro se agitó, y con un suave pop, dos botellas frías de OmniCola cayeron en la bandeja de recolección. Eran idénticas a la que el padre de Julio había bebido.

Julio tomó una y se la ofreció a Sebas. —Pruébala. Es del universo donde existe Omniman de verdad.

Sebas tomó la botella, su mano temblaba ligeramente. La destapó, y el suave burbujeo llenó el silencio. Tomó un sorbo, sus ojos se abrieron de par en par.

\begin{quote}
    —¡Wow! ¡Esto es… increíble! ¡Es como si supiera exactamente lo que me gusta! —exclamó, sus dudas desvaneciéndose ante el sabor—. Julio… ¿esto es… esto es de verdad?
\end{quote}

Julio asintió, su sonrisa era de alivio y triunfo. Compartir su secreto con Sebas era liberador.

\begin{quote}
    —Esto es solo el comienzo, Sebas —dijo Julio.
\end{quote}
La noche de videojuegos se había transformado en una aventura de ciencia ficción, y la amistad de ambos había cruzado un umbral inesperado.

\cleardoublepage % Para que el capítulo empiece en una página impar

% Decimotercer capítulo
\chapter*{13 – Tiempo Restante}
\addcontentsline{toc}{chapter}{13 – Tiempo Restante} % Añadir al índice manualmente

El sótano de Julio, que antes olía a polvo y a proyectos a medias, ahora respiraba **tecnología punta**. Después de una semana de trabajo arduo e ininterrumpido, Julio había logrado el objetivo de su vida. Los paquetes habían llegado de todas partes del mundo, llenos de componentes electrónicos tan avanzados que parecían sacados de una película de ciencia ficción. Sus manos, ágiles y expertas, habían soldado, conectado y ensamblado cada pieza con la precisión de un cirujano.

La máquina ya no era el "armatoste" que parecía un cúmulo de chatarra. Había pasado de ser un montón de piezas improvisadas a un **dispositivo elegante y futurista**, con paneles que se iluminaban con una luz azul suave, circuitos que brillaban con energía contenida y una estructura metálica pulcra que ya no revelaba su origen humilde en un horno de microondas. Ahora, era una pieza de arte tecnológico, tan real que cualquier persona que la viera, incluso sin entender de ingeniería, sabría que es algo de otro mundo, un portal a lo desconocido. Era una proeza.

Y el día de hoy, además, era especial. Era su **cumpleaños número 15**. Quince años. La adolescencia en su punto máximo, y con ella, un destino que ni siquiera podía imaginar.

Julio había esperado este momento con una mezcla de emoción y nerviosismo. Había instalado los últimos componentes esa misma tarde, y la máquina estaba lista. No podía esperar más.

En el sótano, la expectación era casi palpable. **Sebastián**, su mejor amigo y ahora cómplice parcial, estaba a su lado, con los ojos bien abiertos, fascinado por la transformación de la máquina. **Fox**, el gato blanco, aunque ajeno a la magnitud de lo que estaba a punto de suceder, se había acurrucado cerca de la base del dispositivo, como si sintiera la energía latente. Y, por supuesto, **Pi**, el cerebro digital, esperaba las instrucciones de Julio, listo para procesar cualquier dato.

\begin{quote}
    —¿Estás listo, Sebas? —preguntó Julio, su voz apenas un susurro.
\end{quote}

Sebas asintió, incapaz de pronunciar una palabra. La máquina, con su nuevo aspecto futurista, imponía.

Después de muchos intentos y ajustes minuciosos en la semana anterior, Pi había logrado establecer una "**conexión segura**" con la anomalía. Había utilizado el lenguaje binario como base, pero había incorporado bases de su propia **arquitectura adaptativa** —es decir, había aprendido a entender la forma en que la anomalía se comunicaba y había creado un puente que funcionaba para ambos lados—. La comunicación ahora era mucho más fluida.

\begin{quote}
    —¿Aceptar solicitud de conexión? —preguntó Pi, mostrando el mensaje en letras grandes y brillantes en la pantalla principal de la consola. La pregunta flotaba en el aire del sótano, casi con una voz propia.
\end{quote}

Julio dudó. Su corazón latía con fuerza, como un tambor frenético en su pecho. Este era el paso más grande que había dado hasta ahora. Fox, el gato, que había sido testigo de tantas cosas extrañas en el sótano, dormitaba ajeno a la tensión en su cama improvisada de cables, mientras afuera, a través de la pequeña ventana del sótano, la noche era densa y silenciosa, una oscuridad que parecía envolver el barrio.

Con un leve temblor en sus manos, que no podía controlar, Julio asintió.

\begin{quote}
    —Sí… acepta —dijo, su voz un hilo apenas audible, pero firme. La decisión estaba tomada.
\end{quote}

En cuanto Pi ejecutó el comando, todo el cuarto se iluminó con una luz que no provenía de ninguna bombilla. No era un brillo que salía de la máquina, sino que el espacio mismo parecía transformarse. Las paredes desaparecieron. El techo se desvaneció. Las estanterías llenas de herramientas y los juguetes viejos de Óscar se disolvieron. En su lugar, el sótano se pintó con la **inmensidad del espacio exterior**, como si hubieran sido transportados a millones de kilómetros de la Tierra. Pequeñas partículas de polvo que flotaban en el ambiente se detuvieron y fueron atrapadas en un remolino invisible, danzando entre estrellas lejanas.

Frente a ellos, en lo que antes era la pared, una **imagen colosal y tridimensional** comenzó a materializarse, suspendida en el aire. Primero, solo eran formas abstractas, colores y patrones que cambiaban y giraban sin sentido aparente, como una pintura cósmica en movimiento. Luego, muy lentamente, esas formas se fueron definiendo, condensándose en una figura clara: una **Tierra moribunda**, un planeta azul y verde que se transformaba en una esfera agrietada y desolada, azotada por ráfagas de energía devastadoras. El brillo verde de los océanos se apagaba, los continentes se resquebrajaban. Detrás de esta visión apocalíptica, un **Sol agonizante** aparecía, no el Sol brillante que conocían, sino una estrella a punto de explotar, hinchándose y explotando en un espectáculo aterrador de destrucción cósmica. La imagen era tan real que el calor pareció aumentar en el sótano, la desesperación palpable.

Julio tragó saliva, incapaz de apartar los ojos del espectáculo. Sus manos se aferraron al borde de la consola. Sebastián, a su lado, estaba pálido, con la boca ligeramente abierta, sus ojos fijos en la proyección, congelado por el horror y el asombro.

De pronto, una voz, o más bien un **pensamiento colectivo** que parecía hablar directamente en la mente de Julio y Sebas al mismo tiempo, sin necesidad de Pi, pero a través de la conexión establecida, resonó en el sótano. Era una voz sin sonido, pero que sentían en lo más profundo de su ser, un mensaje telepático que les llegaba directo al alma.

\begin{quote}
    —Somos Ellos. No eres el primero, Julio. Otros de ti han logrado llegar hasta aquí antes.
\end{quote}

Julio dio un paso atrás, boquiabierto. La voz en su mente, la imagen de la Tierra destruida… era demasiado. ¿Otros como él? ¿Cuántos? ¿Y por qué ellos no habían podido detenerlo?

La voz continuó, solemne y grave, cada palabra cargada de un peso inimaginable:

\begin{quote}
    —La tarea que tienes ante ti no es sencilla. Se acerca una tormenta, un final ineludible para tu mundo: la **radiación de una supernova**, viajando hacia tu planeta desde hace milenios. El impacto… será devastador.
\end{quote}

La imagen cambió de nuevo, como si les dieran una prueba del destino de su planeta. Mostró un planeta Tierra reducido a un **desierto estéril**, sin océanos, sin vida, con solo dunas de arena y rocas calcinadas. Los continentes eran apenas sombras. Era una visión de absoluta desolación, un futuro sin esperanza.

\begin{quote}
    —Tu familia. Tus amigos. Tu hogar. Todo lo que amas será destruido… a menos que estés dispuesto a aceptar el desafío.
\end{quote}

Julio apenas podía respirar. Su mente corría en círculos, buscando una forma de negar lo que estaba viendo, lo que estaba escuchando. ¿Salvarlos? ¿Cómo? ¿Él solo, un chico de quince años, contra la radiación de una supernova? La idea era absurda, imposible, y al mismo tiempo, era su responsabilidad.

La proyección mostró entonces una **cuenta regresiva en números brillantes**, grandes y claros, proyectados justo frente a él en el aire, flotando como si estuvieran grabados en el tiempo:

\begin{quote}
    \textbf{489 días, 8 horas, 30 minutos, 14 segundos… 13 segundos… 12 segundos…}
\end{quote}

El tiempo no paraba. Cada segundo que pasaba era un segundo menos. La inevitabilidad de la catástrofe era aterradora.

\begin{quote}
    —El tiempo corre, Julio. La pregunta es: ¿estarás a la altura?
\end{quote}

La proyección se desvaneció en un susurro. La inmensidad del espacio exterior desapareció de las paredes y el techo, y el sótano de Julio, con sus estanterías y sus herramientas, regresó a la normalidad, iluminado solo por la luz de la consola de Pi. El taller quedó en completo silencio, roto solo por el parpadeo constante de la pantalla y el ronroneo tranquilo de Fox, ajeno al destino que acababa de ser revelado.

Julio se quedó allí, helado, inmóvil, observando los números decrecer en la pantalla, que Pi ahora proyectaba en su consola principal, como un reloj inexorable. Sebastián, a su lado, también estaba congelado, su rostro pálido y sus ojos vacíos, procesando la información. La alegría del cumpleaños, la cena familiar que lo esperaba arriba, todo se desvaneció en ese instante. Lo que había comenzado como un día de festejo con una cena muy agradable en familia, se había convertido en un golpe emocional devastador que dejó a ambos niños sumergidos en una profunda desesperación y un miedo paralizante. El universo acababa de presentarles una amenaza real, un reloj en cuenta regresiva hacia el fin. Y ellos eran los únicos que lo sabían.

\cleardoublepage % Para que el capítulo empiece en una página impar

% Decimocuarto capítulo
\chapter*{14 – Demasiado grande para mí}
\addcontentsline{toc}{chapter}{14 – Demasiado grande para mí} % Añadir al índice manualmente

El festejo de cumpleaños de Julio, que prometía ser una noche de risas y alegría, se había teñido de una amargura inesperada. Los padres de Julio, Ana y Miguel, notaron el cambio de inmediato. Cuando Sebastián había llegado, antes de la cena, ambos chicos estaban radiantes, emocionados y saltando de alegría, planeando su maratón de videojuegos. Pero después de haber pasado apenas treinta minutos en el sótano con la máquina, el ambiente cambió drásticamente. Cuando subieron para la cena de cumpleaños, ambos niños permanecían extrañamente callados y pálidos. Sus rostros reflejaban una preocupación que no podían disimular, y sus ojos, antes llenos de brillo, ahora parecían contener un peso que no correspondía a su edad. La torta de chocolate, los regalos, todo parecía menos importante.

Al final de la noche, después de que Óscar se había quedado dormido en el sofá con una sonrisa azucarada y Julio había limpiado el resto de las botanas, era el momento de la despedida de Sebastián. Sebas se puso de pie, su expresión era aún grave, pero había una determinación en sus ojos. Los padres de Julio lo observaban con curiosidad.

\begin{quote}
    —Bueno, Julio —dijo Sebastián, su voz llena de una solemnidad que sorprendió a todos—. Ya me voy.
\end{quote}

Se acercó a Julio, le puso una mano en el hombro y lo miró directamente a los ojos. En ese gesto, en esa mirada, había un peso que solo ellos dos entendían.

\begin{quote}
    —Bro, si alguien puede hacer esto, eres tú. —dijo Sebas, con una convicción que no dejaba lugar a dudas—. No te preocupes por nada. Voy a estar aquí para lo que necesites, ¿me escuchas? Lo que sea. Si necesitas que te ayude a cargar cajas, a buscar algo, a lo que sea. No estás solo en esto.
\end{quote}

Julio asintió, conmovido por la sinceridad de su amigo. Era un tipo de apoyo que no había esperado, una lealtad incondicional.

Sebas se despidió de Ana y Miguel, algo que solía hacer de forma rápida, pero esta vez lo hizo con un agradecimiento más profundo, como si comprendiera el verdadero significado de esa noche. Los padres de Julio se miraron, asombrados por la seriedad y el apoyo que Sebas había mostrado hacia su hijo. Percibieron un **crecimiento personal y social** en Julio; no solo estaba haciendo un amigo, sino que había encontrado un compañero de verdad, alguien en quien confiar en momentos difíciles. Fue un pequeño consuelo en medio de la preocupación que sentían por el extraño estado de ánimo de los chicos.

---

Una vez que Sebas se fue a su casa, a pocas cuadras de allí, el silencio regresó al hogar de Julio. La cuenta regresiva parpadeaba en la vieja laptop en su sótano como un latido oscuro e inexorable: **489 días, 3 horas, 29 minutos, 52 segundos...**

Julio estaba sentado en el suelo de su taller, abrazando sus rodillas, con la mirada fija en los números que seguían disminuyendo sin piedad. El peso de lo que había visto, de la responsabilidad que se le había impuesto, lo aplastaba. Fox, se acurrucaba en su regazo, ronroneando suavemente, ofreciendo un consuelo mudo en medio de la inmensa desesperanza.

\begin{quote}
    —¿Cómo… se supone que haga algo? —susurró Julio, apenas audible, su voz quebrada por la desesperación.
\end{quote}

Julio cerró los ojos, el rostro oculto entre sus rodillas. La idea de salvar al mundo era demasiado. Era ridícula. ¿Qué podía hacer un niño que ni siquiera había terminado la secundaria? Él, que a veces fallaba en mantener su cuarto ordenado, que apenas entendía algunas cosas de física avanzada que había aprendido en línea por curiosidad. Se sentía insignificante frente a la magnitud de una supernova.

Se puso de pie lentamente, sintiendo que el peso invisible del desafío lo aplastaba, como si llevara el peso del mundo sobre sus hombros. Caminó hasta la pizarra que tenía en el rincón, donde solía anotar ideas locas para proyectos de robótica y experimentos caseros. Normalmente, esa pizarra era un lienzo de posibilidades, pero ahora se sentía vacía, esperando una respuesta que él no tenía.

Respiró hondo, intentando controlar el pánico. Tomó un marcador. Y escribió una pregunta simple en el centro de la pizarra, la primera idea que le vino a la mente, por tonta que le pareciera, para empezar a desglosar el problema:

\begin{quote}
    ¿Cómo sobrevivir a una radiación cósmica?
\end{quote}

Alrededor de esa pregunta, comenzaron a aparecer nuevas burbujas de pensamiento, cada una más inalcanzable que la anterior:

\begin{itemize}
    \item Crear un escudo impenetrable.
    \item Construir refugios subterráneos profundos.
    \item Utilizar materiales reflectantes o absorbentes.
    \item Desviar o bloquear la energía directamente.
\end{itemize}

Cada nueva idea traía consigo una nueva angustia. La tecnología para hacer algo así estaba tan fuera de su alcance que parecía fantasía. Ni siquiera sabía si era posible con los recursos de la Tierra.

\begin{quote}
    —Demasiado grande para mí —pensó, dejando caer el marcador al suelo, el ruido resonando en el silencio abrumador del sótano.
\end{quote}

Esa noche, sus padres notaron que algo no andaba bien con Julio. Después de la cena, al ver la palidez en su rostro y la forma en que se había encerrado en su sótano, Ana y Miguel intercambiaron miradas preocupadas. Miguel, su padre, un hombre que no siempre mostraba sus emociones pero que amaba profundamente a sus hijos, decidió que era momento de hablar.

Bajo las escaleras con dos tazas humeantes de chocolate caliente. Golpeó suavemente la puerta del sótano.

\begin{quote}
    —¿Puedo pasar? —preguntó su padre, su voz suave y tranquila.
\end{quote}

Julio, con la mirada perdida en la pizarra, apenas asintió.

Miguel se sentó a su lado en el suelo, le entregó una taza caliente y, por un momento, simplemente bebieron en silencio. No había prisa, solo presencia.

\begin{quote}
    —Sabes, campeón —dijo finalmente su padre, su voz resonando con una preocupación sincera, un orgullo apenas disimulado por el esfuerzo de su hijo, y una fortaleza silenciosa que Julio siempre había admirado—. Cuando tenía tu edad, también pensaba que el mundo era demasiado grande para mí. Cada problema parecía imposible de resolver. Cada sueño, inalcanzable.
\end{quote}

Julio apretó la taza con fuerza, sintiendo el calor en sus manos.

\begin{quote}
    —¿Y qué hiciste? —preguntó, la esperanza apenas una chispa.
\end{quote}

Miguel sonrió con nostalgia, recordando sus propias batallas juveniles.

\begin{quote}
    —Me di cuenta de que nadie cambia el mundo, o logra sus metas en un solo día. No puedes ver el final del camino desde el principio. A veces, un gran problema se resuelve con un **pequeño paso**. Tienes que disminuirlo, desglosar todo a la respuesta más simple, aunque te parezca tonta. Empieza por lo más básico. Primero salvas algo pequeño: una idea, una persona, un momento. Después de eso, el resto sigue. Es como construir un edificio, Julio. No empiezas por el techo, empiezas por la base. Un ladrillo a la vez.
\end{quote}

La sabiduría de su padre, expresada de una forma tan natural y sin pretensiones, golpeó a Julio con la fuerza de una revelación. Sintió que algo dentro de él se rompía, toda esa desesperación y miedo, y al mismo tiempo se reconstruía, llenándose de una nueva determinación.

Quizás… no necesitaba salvar el mundo entero ahora mismo. Quizás solo tenía que salvar una parte. La parte que pudiera alcanzar. La parte que estuviera directamente bajo su control. Su familia. Su hogar. Era un primer paso.

Cuando sus padres se retiraron a dormir, dejando la casa en un silencio hogareño, Julio se quedó en su taller, el chocolate ya frío sobre la mesa. Fox, fiel a su lado, lo miraba con esos ojos dorados llenos de paciencia, como si supiera que su amo estaba en un momento crucial.

\begin{quote}
    —Pi —dijo Julio, su voz ahora más firme, con un propósito renovado—. Necesitamos un plan. Algo que podamos hacer de verdad. Algo que sea el **primer ladrillo**.
\end{quote}

La inteligencia artificial respondió enseguida, su tono robótico pero eficiente:

\begin{quote}
    —Propongo iniciar con un **refugio experimental para la familia**. Prototipo uno: blindaje contra radiación. Es un problema con variables controlables.
\end{quote}

Julio asintió. Tomó su cuaderno de proyectos y empezó a dibujar frenéticamente, con una energía renovada: planos de un pequeño cuarto subterráneo, capas de materiales reflectantes, generadores de energía aislados, sistemas de soporte vital. Su mente, antes paralizada, ahora corría a mil por hora, desglosando el inmenso problema en pasos manejables.

Por primera vez en horas, la desesperanza se transformó en propósito. La tarea seguía siendo gigantesca, pero ahora tenía un punto de partida.

Mientras dibujaba, Fox se estiró perezosamente y se acomodó más cerca de su pierna, ronroneando con fuerza. Pi desplegaba modelos 3D y cálculos complejos sobre las paredes usando un proyector improvisado, convirtiendo el sótano en una sala de mando personal.

La cuenta regresiva seguía avanzando, indiferente a sus miedos o sus planes.

\begin{quote}
    \textbf{488 días, 23 horas, 24 minutos, 03 segundos...}
\end{quote}

Pero ahora, Julio ya había comenzado a luchar.

\cleardoublepage % Para que el capítulo empiece en una página impar

% Decimoquinto capítulo
\chapter*{15 – La Primera Prueba}
\addcontentsline{toc}{chapter}{15 – La Primera Prueba} % Añadir al índice manualmente

El ambiente en casa de Julio había cambiado desde la noche del cumpleaños. La alegría del festejo se había disuelto en una tensión silenciosa, especialmente alrededor de Julio. Óscar, el hermano menor de Julio, ya no podía quedarse callado. Llevaba semanas observando a su hermano mayor esconderse en el sótano, salir al jardín a horas extrañas, y hablar solo. El cambio en Julio era demasiado evidente. Después de aquella cena de cumpleaños tan extraña, Óscar estaba aún más convencido de que su hermano mayor estaba metido en algo grande y misterioso.

Una tarde, mientras Julio trabajaba concentrado en su portátil en el sótano, revisando los planos de la supernova y las soluciones de blindaje, Óscar se plantó frente a él, cruzado de brazos y con el ceño fruncido.

\begin{quote}
    —¡Ya basta, Julio! —soltó de golpe, con un tono de voz que mezclaba frustración y preocupación—. ¡Sé que estás haciendo algo raro! ¡Y no me digas que es tarea porque ni siquiera has entregado la del mes pasado en el colegio! Además, estás más callado que de costumbre.
\end{quote}

Julio levantó la vista, sorprendido por la confrontación directa, y vio la determinación en los ojos de su hermano. Óscar, a pesar de ser más joven y estar más interesado en los dulces, tenía una intuición sorprendente. Julio se quedó en silencio un instante, midiendo sus opciones. Después de la revelación de la supernova, se había sentido abrumado y solo, incluso con Sebas a su lado. No podía seguir ocultándolo. No a él. Ya no tenía la energía para inventar excusas.

Suspiró profundamente y cerró la laptop, dándole toda su atención a Óscar.

\begin{quote}
    —Está bien —dijo finalmente, su voz un poco tensa—. Te lo contaré todo, pero prométeme que no vas a decirle a nadie… ni a mamá, ni a papá, por ahora. Esto es… muy, muy serio.
\end{quote}

Óscar, entusiasmado por la promesa de un secreto tan grande, asintió vigorosamente como si acabara de ganar un premio, sus ojos brillando con curiosidad.

Julio entonces le contó todo, desde el principio: la máquina, el material desconocido que había cruzado, la comunicación con "ellos", el inminente peligro de la supernova, y el reloj que inexorablemente marcaba el tiempo que les quedaba. Mientras hablaba, Julio se acercó a su consola y activó una **grabación tridimensional** de lo que había sucedido en el sótano el día de su cumpleaños. De pronto, todo el cuarto se iluminó con una luz que no provenía de ninguna bombilla, sino que el espacio mismo parecía transformarse de nuevo. Las paredes desaparecieron. El techo se desvaneció. Las estanterías llenas de herramientas y los juguetes viejos de Óscar se disolvieron. En su lugar, el sótano se pintó con la **inmensidad del espacio exterior**, mostrando la imagen de una Tierra moribunda y un Sol agonizante, con la aterradora cuenta regresiva flotando en el aire. La voz de "Ellos" resonó en la mente de Óscar, explicándole la amenaza de la supernova.

Óscar, por supuesto, reaccionó como siempre: al principio soltó una carcajada, pensando que era una broma muy elaborada de Julio o un nuevo videojuego. Luego, al ver la seriedad en el rostro de su hermano y al sentir la voz resonando en su propia mente, lo miró boquiabierto, sus ojos tan grandes como canicas. Cuando la proyección desapareció y el sótano volvió a la normalidad, Óscar seguía pálido, y su boca, antes llena de una sonrisa, ahora estaba ligeramente abierta. El asombro se transformó en una clara preocupación.

\begin{quote}
    —Y… por último, Óscar —dijo Julio, activando la consola de su computadora—. Él es Pi.
\end{quote}

Una voz sintética, pero con un tono amigable, resonó en el sótano, sorprendiendo a Óscar.

\begin{quote}
    —Hola, Óscar. Soy Pi, la inteligencia artificial de Julio. Es un placer conocerte.
\end{quote}

Óscar dio un salto hacia atrás, su rostro se puso blanco y sus ojos se clavaron en la pantalla. —¡Santa madre de mi vida y de mis dulces! —exclamó, con la voz ahogada por la sorpresa y la incredulidad—. ¡Eso es… eso es de película! ¿Desde cuándo tienes una computadora que habla?

Julio sonrió. —No es una computadora, Óscar. Es una **inteligencia artificial**. Y sí, todo esto es real. Por eso he estado tan ocupado, por eso he trabajado tanto. Y por eso necesito que me ayudes.

Óscar, con el asombro todavía en su cara, finalmente reaccionó.

\begin{quote}
    —¿Qué hacemos ahora? —preguntó sin dudarlo, el tono de broma había desaparecido por completo—. No pienso dejar que nos pase nada. Ni a nosotros, ni a mamá y papá, ni a nadie.
\end{quote}

Julio sonrió con un alivio inmenso. Nunca había estado completamente solo en esto. Óscar siempre estuvo ahí, incluso cuando no entendía del todo lo que pasaba. Su hermano, el comilón y enérgico, sería un aliado valioso.

\begin{quote}
    —Necesitamos construir un **refugio** —explicó Julio, volviendo a su plan inicial—. Uno que pueda soportar radiaciones intensas y que sea seguro para la familia.
\end{quote}

Óscar chasqueó los dedos, una nueva idea encendiéndose en su mente, la energía volviendo a él.

\begin{quote}
    —¡El **viejo taller del abuelo**! —exclamó—. Está abandonado, pero las paredes son gruesas, de concreto sólido. ¡Podemos reforzarlo! Es el lugar perfecto, está en el jardín trasero, aislado.
\end{quote}

La idea era buena. El taller del abuelo, un pequeño cobertizo de cemento reforzado, había estado abandonado por años, pero su estructura era robusta. Para ponerlo en marcha, Julio sabía que necesitaría el apoyo de su padre.

Esa misma noche, Julio encontró a su padre en el garaje, trabajando en uno de sus proyectos.

\begin{quote}
    —Papá, ¿puedo hablar contigo un minuto? —dijo Julio, intentando sonar lo más casual posible.
\end{quote}

Miguel levantó la vista, sorprendido de ver a Julio fuera del sótano a esa hora.

\begin{quote}
    —Claro, campeón. Dime.

    —Mira, el sótano me está quedando… pequeño para mis proyectos —empezó Julio, con una naturalidad que le costaba—. Tengo muchas cosas, y mis proyectos, ya no cabe todo. Estaba pensando en si podría… modificar el viejo taller del abuelo en el jardín. Esas paredes son bien fuertes, y ahí tendría más espacio para todo. Sería como mover mi taller fuera de la casa. Así no haría tanto ruido adentro y mamá no se quejaría de mis cables por todos lados.
\end{quote}

Miguel se quedó pensando, rascándose la barbilla. El taller del abuelo era un trastero que nadie usaba. Tenía sentido que Julio quisiera un espacio más grande, dada su pasión por la electrónica. Miguel, siempre orgulloso de la dedicación de su hijo, vio esto como una señal de su creciente interés por la ingeniería.

\begin{quote}
    —Hmm, no es una mala idea, Julio. Ese cobertizo es robusto. Podemos revisarlo mañana, ver qué tan bien está. Pero si lo vas a usar, tendrás que limpiarlo a fondo y hacerte cargo tú de las reformas, ¿entendido? Nada de dejar las herramientas tiradas como en el sótano.

    —¡Sí, papá! ¡Lo prometo! —dijo Julio, sintiendo un nudo de alivio en el estómago. La primera parte de su plan, la más "normal", había funcionado.
\end{quote}

Mientras Julio dibujaba los primeros esquemas del refugio, se dio cuenta de que no podían hacerlo solos. Necesitaban más manos, más fuerza. Tomó su teléfono y marcó.

\begin{quote}
    —Sebas, ¿estás libre? Necesito tu ayuda con algo… algo muy grande.
\end{quote}

La respuesta de Sebastián fue casi inmediata, llena de alegría y entusiasmo, como era de esperar de él.

\begin{quote}
    —¡Claro que sí, bro! Lo que sea. Dime dónde. Sabes que puedes contar conmigo.
\end{quote}

Julio sintió un alivio inmenso al escuchar la voz animada de su amigo. Le explicó a Sebas sobre el taller del abuelo y la necesidad de construir un refugio, aún sin darle todos los detalles cósmicos. Sebas aceptó casi de inmediato, e incluso ofreció una ayuda inesperada.

\begin{quote}
    —¡Espera un segundo! Mi padre trabaja en una **empresa constructora local** —dijo Sebas con orgullo en su voz—. Él es operador de maquinaria pesada, es el mejor en lo que hace. ¡Maneja camiones, excavadoras, de todo! Seguro tiene en nuestro sótano herramientas y materiales olvidados que podríamos pedir prestados. Espátulas, carretillas, quizás hasta cemento que le haya sobrado de alguna obra. ¡Déjame revisar!
\end{quote}

La ayuda de Sebas era invaluable. No solo traía la fuerza física que Julio y Óscar necesitaban, sino también el acceso a recursos que nunca habrían imaginado. El entusiasmo de Sebas era contagioso.

La reconstrucción del refugio fue un **proyecto monumental**. Entre risas, errores, y algunos accidentes menores (como cuando Óscar, torpe pero lleno de buena voluntad, tropezó y tiró una carretilla de cemento fresco sobre sus propios zapatos, dejando una huella gris en el patio), los tres amigos comenzaron a darle forma a su primer refugio.

Consiguieron todo tipo de materiales. Gracias a Sebas, que con el permiso de su padre pudo traer algunas cosas: **placas de plomo** de una chatarrería local (que encontraron en un rincón olvidado de un almacén de su padre), **concreto extra grueso** que mezclaron ellos mismos (después de varios intentos fallidos que terminaron en montañas de polvo), y **capas de resina endurecida** para impermeabilizar y sellar las paredes contra cualquier elemento externo.

El proceso fue agotador. Pasaron días enteros cargando materiales pesados, martillando, reforzando las paredes del antiguo taller del abuelo con nuevas capas. Julio, aunque físicamente no era tan resistente como Sebas, tenía la mente clara, y diseñaba cada sección como si fuera una ecuación matemática, calculando ángulos, resistencias y el grosor del blindaje.

Óscar, en cambio, era pura energía y entusiasmo. Siempre dispuesto a ayudar, cargaba sacos de arena más pequeños, pasaba herramientas y, por supuesto, siempre estaba pidiendo un descanso para comer algo. Era un remolino de vitalidad que mantenía el ánimo en alto. Pero el verdadero músculo del equipo era Sebastián. Con su fuerza de deportista, Sebastián era quien hacía las tareas pesadas que requerían gran esfuerzo físico: movía los sacos más grandes de cemento, empujaba las carretillas llenas de grava y concreto, y ayudaba a Julio a levantar las pesadas placas de plomo.

A veces discutían, claro, como cualquier equipo. Como aquella vez en que Óscar, impaciente por terminar e ir a comer, colocó mal un muro de concreto y todo terminó desmoronándose encima de unas cajas de herramientas de Julio.

\begin{quote}
    —¡Te dije que había que esperar a que secara! —gritó Julio, con una mezcla de frustración y diversión.

    —¡Y yo te dije que tenía hambre! —respondió Óscar, muerto de risa mientras se quitaba el polvo de encima, ajeno a la gravedad de la situación.
\end{quote}

Incluso Pi participaba, controlando ciertos dispositivos electrónicos que Julio había instalado para medir niveles de radiación simulada y dando instrucciones de seguridad precisas sobre la mezcla del concreto o la ventilación. Era su forma de supervisar el proyecto.

Finalmente, después de semanas de esfuerzo conjunto, de mañanas enteras bajo el sol y noches estudiando planos, el refugio estaba listo.

La estructura no era bonita, ni elegante. Era una **fortaleza improvisada**, un búnker de aspecto tosco pero innegablemente sólido en el patio trasero de la casa de Julio. Una verdadera primera línea de defensa.

Julio, Óscar y Sebas se sentaron en el suelo, sudorosos, exhaustos, con rasguños en las manos, pero con una sonrisa de satisfacción que no podían ocultar. Habían creado algo tangible, algo que parecía posible.

El cielo nocturno se extendía sobre ellos, salpicado de estrellas, cada una un recordatorio silencioso del peligro que se avecinaba.

\begin{quote}
    —¿Crees que funcione? —preguntó Óscar, con un tono más serio de lo habitual, mirando la fortaleza a medio iluminar por la luz de una linterna.
\end{quote}

Julio miró el refugio, luego a su hermano, luego a Sebas, que asintió con confianza, y finalmente a las estrellas, el origen de su problema.

\begin{quote}
    —Tiene que funcionar —dijo, en voz baja, con una nueva determinación.
\end{quote}

Pi, desde el pequeño altavoz que habían conectado cerca, soltó un sonido de afirmación, un bip tranquilizador.

\begin{quote}
    —Preparados para la primera prueba de resistencia en **24 horas** —anunció la inteligencia artificial.
\end{quote}

Julio sintió un escalofrío recorrerle la espalda, pero esta vez, no era solo miedo. Era la adrenalina. Todo estaba a punto de comenzar.

\cleardoublepage % Para que el capítulo empiece en una página impar

% Decimosexto capítulo
\chapter*{16 – El Eslabón Perdido}
\addcontentsline{toc}{chapter}{16 – El Eslabón Perdido} % Añadir al índice manualmente

La mañana amaneció gris, con el sol apenas alcanzando a filtrarse entre las nubes densas. Una sensación de desesperación se cernía sobre Julio. No podía dejar de pensar en el contador de la máquina, esa maldita **cuenta regresiva** que parecía burlarse de él desde la pantalla de su laptop, cada segundo una punzada de pánico. En su mente, solo había una cosa clara: la radiación de la supernova era mucho más poderosa de lo que habían previsto. La prueba del refugio era inminente, y su fracaso sería devastador.

Óscar, al ver a su hermano tan absorto y con el ceño fruncido, lo interrumpió con una pequeña broma, intentando aligerar el pesado ambiente, aunque en el fondo sabía que lo que estaba por hacer no tenía cabida para chistes.

\begin{quote}
    —¿Estás listo para la gran prueba, o necesitas dormir un poco más? —le preguntó Óscar, ya bastante acostumbrado a que los planes de Julio se alargaran hasta altas horas de la madrugada, dejándolo con ojos inyectados en sangre.
\end{quote}

Julio asintió sin mucha energía, la fatiga pesaba en cada uno de sus movimientos. Había algo extraño en su mirada, algo que no podía disimular: el estrés le estaba pasando factura.

\begin{quote}
    —No es un buen día para dormir. Vamos a asegurarnos de que funcione, Óscar —respondió Julio, con una voz tensa—. El refugio… es nuestra única oportunidad.
\end{quote}

Óscar, sin hacer más preguntas, simplemente asintió y se dirigió a las consolas de control que habían instalado en el laboratorio improvisado, que ahora era el sótano de la casa. **Pi** estaba listo para la prueba. El sistema se activó, las pantallas de las computadoras se iluminaron con gráficos y datos, y las luces del refugio, visible a través de un monitor de seguridad, encendieron. Todo estaba en su lugar.

La idea de Julio para la prueba era simple y brutalmente lógica: si la radiación de un mísero microondas lograba salir del refugio, entonces una supernova definitivamente podría entrar. Harían uso de los componentes de un microondas que instalarían dentro del refugio, convirtiéndolo en un emisor controlado de radiación. La premisa era sencilla: la misma lógica que permitiría que la radiación saliera, significaría que la radiación externa podría ingresar. Con **medidores de radiación** dentro del refugio y otros instalados fuera, podrían medir de forma precisa y controlada si el refugio dejaba escapar la radiación. Si no contenía la poca energía de un microondas, no contendría la de una estrella.

En el exterior, en el jardín, los sensores externos de radiación estaban calibrados. Dentro del refugio, Julio activó el emisor de microondas. Un zumbido comenzó a resonar, y en el monitor, las primeras lecturas fueron tranquilizadoras. El refugio parecía contener el impacto de las ondas electromagnéticas, actuando como un escudo eficaz. Pero a medida que la radiación aumentaba, forzando los límites de su construcción, comenzaron a sonar alertas en los sistemas de seguridad.

Primero, las paredes comenzaron a emitir ligeras vibraciones, casi imperceptibles al tacto, pero Julio pudo sentirlas a través del suelo y las pantallas de los sensores. El sudor comenzó a acumularse en su frente, la piel de su cuello erizándose.

\begin{quote}
    —¿Por qué está vibrando tanto? —se preguntó, observando las pantallas de los sensores de temperatura y presión. La temperatura en la zona del generador de energía estaba subiendo demasiado rápido, mucho más de lo esperado, como si algo estuviera a punto de colapsar.
\end{quote}

Óscar se acercó a él, con las manos sudadas, una mezcla de miedo e incertidumbre en sus ojos.

\begin{quote}
    —¿Julio? ¿Está todo bien? —su tono de voz no era de pánico, sino de una preocupación genuina.

    —¡No! No está bien… —Julio murmuró, pero fue interrumpido por el sonido estridente de las alarmas, ahora a todo volumen. La ventilación principal del refugio comenzó a sobrecalentarse, expulsando humo de los conductos, y el sistema de radiación estaba marcando niveles peligrosos. Las luces del sótano comenzaron a parpadear frenéticamente, y Pi emitió una serie de pitidos agudos, alertando de un fallo inminente en la contención del refugio. Los sensores externos, que al principio marcaban cero, comenzaron a detectar **filtraciones de radiación**. Pequeñas, casi insignificantes, pero presentes. La radiación del microondas estaba escapando.

    —Esto no va a funcionar… —dijo Julio, sintiendo cómo el sentimiento de desesperación lo invadía por completo. El refugio había sido una idea brillante, el primer paso que su padre le había animado a dar, pero la realidad era otra. No estaba preparado para todo esto. Los cálculos, los experimentos, las pruebas preliminares… todo había fallado en esta simulación a gran escala. Todo se dio por terminado en fracaso.
\end{quote}

Se sentó sobre una silla, cubriéndose la cara con las manos, la cabeza zumbándole. ¿Qué más podía hacer? Si el refugio no resistía la radiación de esta simulación, ¿cómo podría sobrevivir a la supernova? El tiempo se le escapaba entre los dedos como arena.

---

Óscar, sin saber muy bien cómo consolarlo, miró hacia el techo y luego a las pantallas, procesando la información.

\begin{quote}
    —Julio… ¿por qué no… por qué no les preguntas? —dijo, mirando a su hermano con una mezcla de desesperación y una chispa de esperanza, recordando la sorprendente revelación de la voz que había escuchado en su mente.

    —¿A quién? —respondió Julio, con voz quebrada, sin comprender del todo. Su mente estaba bloqueada por la derrota.

    —A ellos —dijo Óscar, señalando la máquina con un dedo tembloroso—. Sabes a quién me refiero. La voz, la que nos habló. No sé si alguna vez les habías consultado algo, pero… si la máquina hizo lo que hizo antes, ¿por qué no intentar otra vez? Quizás tengan una idea.
\end{quote}

Julio lo miró, la sugerencia de Óscar resonando en su cabeza. El orgullo le decía que no necesitaba ayuda, que podía resolverlo por sí mismo. Que un verdadero científico no dependía de respuestas externas. Pero las últimas semanas, los eventos, la cuenta regresiva… todo lo que había logrado, todo lo que había fallado… le hizo pensar que, tal vez, consultar con “ellos” podría ser su única opción. Había un poder más allá de su comprensión en esa máquina, una conexión que no podía ignorar.

\begin{quote}
    —Tienes razón, Óscar —dijo Julio, con un suspiro que liberó parte de la tensión acumulada—. Voy a preguntarles. Tal vez tengan alguna respuesta.
\end{quote}

Con manos temblorosas, Julio se acercó a la máquina principal. En el interior de su mente, las palabras que había dicho a la máquina antes, en sus primeros experimentos, resonaban: "¿Qué pasa si...?" Ahora, no necesitaba el código binario. Gracias a las mejoras y a la evolución de Pi, solo tenía que escribir de forma natural en la computadora.

Abrió el sistema, activó la interfaz y con un golpe de teclado, escribió su pregunta, llena de la urgencia que sentía:

\begin{quote}
    Estoy listo. ¿Cómo puedo salvarnos de la supernova?
\end{quote}

Se sentó, esperando, el sudor en su cuello cayendo gota a gota mientras observaba la pantalla en blanco, esperando una respuesta, el silencio del sótano solo roto por el suave zumbido de los componentes de la máquina y las alertas de Pi.

Después de unos segundos que parecieron una eternidad, la máquina respondió con una simple línea de texto, las palabras apareciendo en la pantalla con una claridad impactante:

\begin{quote}
    \textbf{"La respuesta la tuviste siempre, desde el principio. Para resolver problemas futuros, debes ir hacia atrás."}
\end{quote}

Julio leyó las palabras una y otra vez. "La respuesta la tuviste siempre, desde el principio." ¿A qué se referían? ¿Qué había pasado por alto? Y lo de "ir hacia atrás"… ¿Era una metáfora? ¿Un viaje en el tiempo? Su mente científica buscaba la lógica, la explicación, pero solo encontraba más preguntas. Esta vez, la respuesta no era ambigua, era frustrantemente directa y aún así, enigmática.

Pi mostró una línea de código adicional que Julio comenzó a estudiar, reconociendo que aún tenía mucho por hacer, mucho que descifrar de esas palabras tan simples y profundas. Sin embargo, algo dentro de él comenzó a calmarse. Una extraña tranquilidad, un pequeño soplo de alivio, lo invadió. Había más dudas que certezas, pero había algo en las palabras de "Ellos", una verdad oculta que le daba cierto grado de confianza. Sabía que estaba pasando por alto algo crucial, algo que había estado frente a él todo el tiempo. Solo necesitaba encontrar el **eslabón perdido**.

El contador seguía su marcha implacable:

\begin{quote}
    \textbf{Tiempo restante: 461 días, 22 horas, 30 minutos, 58 segundos...}
\end{quote}

\cleardoublepage % Para que el capítulo empiece en una página impar

% Decimoséptimo capítulo
\chapter*{17 – La Revelación de Oscar}
\addcontentsline{toc}{chapter}{17 – La Revelación de Oscar} % Añadir al índice manualmente

La noche cayó pesada sobre el barrio. Después de la comunicación con "ellos" y la respuesta críptica que habían recibido —"La respuesta la tuviste siempre, desde el principio. Para resolver problemas futuros, debes ir hacia atrás"—, Julio y Óscar se quedaron trabajando en silencio en el sótano. El cansancio comenzaba a hacer mella en ambos; sus ojos ardían y sus cuerpos dolían, pero ninguno quería ser el primero en rendirse. Había demasiada ansiedad, demasiada incertidumbre, para permitirse el lujo de descansar.

El sótano se había convertido en una mezcla caótica de laboratorio, refugio improvisado y sala de crisis. Hojas con cálculos complejos, modelos a escala del refugio fallido, prototipos de escudos de energía inútiles y todo tipo de bocetos de la máquina: todo estaba allí, desparramado, testigo del esfuerzo casi sobrehumano de los dos hermanos.

Julio estaba encorvado sobre una mesa, examinando por enésima vez los objetos que habían llegado con la anomalía. Desde el principio, había guardado y analizado cada cosa que atravesó el portal, creyendo que la clave para entender la máquina o sus mensajes residía en esos objetos. Ahora, con la sugerencia de "ir hacia atrás", estaba revisando todos sus viejos dibujos y bocetos de la máquina, y los primeros experimentos con el extraño material, el **cubo** que había llegado en el primer envío y que él mismo había extraído del microondas original. Ya lo había probado bajo calor extremo, presión, incluso bajo ondas de radiación simulada. Era increíblemente resistente y parecía inafectado por la radiación, pero algo no terminaba de encajar. Sabía que ese material era especial, un sustituto del americio en la tostadora, pero su función real se le escapaba. Necesitaban algo más, una pieza que no lograba ver.

Óscar, que jugaba distraídamente con un destornillador en las manos, haciendo girar la punta con aburrimiento, dejó caer su cuerpo contra la pared y soltó un gran suspiro. Había estado observando a Julio dar vueltas y vueltas sobre los mismos datos, revisando una y otra vez lo que ya habían visto.

\begin{quote}
    —¿Y si todo esto lo estamos viendo mal? —dijo de repente, mirando al techo con esa extraña mezcla de ingenuidad y lucidez que a veces lo caracterizaba.
\end{quote}

Julio alzó la vista, confundido. —¿Qué quieres decir?

Óscar se encogió de hombros, como si su propia idea le pareciera absurda en voz alta, pero seguía pensando en voz alta.

\begin{quote}
    —Mira… la radiación es energía, ¿no? Siempre la vemos como algo que hay que bloquear o evitar. Es como un golpe, y nosotros ponemos un escudo. Pero… ¿y si, en lugar de protegernos de ella, la usamos? ¿Y si… no tratamos de detener la radiación, sino de **almacenarla**? Si la energía de la supernova es tan grande, ¿no podríamos de alguna forma capturarla o hasta utilizarla?
\end{quote}

Julio parpadeó un par de veces, procesando la frase de su hermano. Era tan simple, tan obvio… que nunca se le había ocurrido. Él había estado obsesionado con la idea de la protección, de la defensa, de construir una barrera impenetrable. Pero Óscar, con su lógica directa y sin las complicaciones de la ciencia avanzada, había dado en el clavo. Era como si su cerebro de científico lo hubiera hecho pensar de una forma demasiado compleja.

Un destello de lucidez cruzó su mente, haciendo que sus ojos se abrieran de golpe. La frase de "Ellos" volvió a resonar: "La respuesta la tuviste siempre, desde el principio." Y la idea de "ir hacia atrás" ahora tenía sentido. No era un viaje en el tiempo, era volver al punto de partida, al material, al cubo, al inicio de todo.

\begin{quote}
    —La radiación es energía —repitió Julio en voz baja, las palabras de Óscar resonando con una nueva verdad—. Energía que podríamos… ¡capturar!
\end{quote}

De inmediato, se abalanzó sobre el extraño material, el **cubo brillante** que había estudiado innumerables veces. Era el mismo material que no parecía verse afectado por la radiación ni emitirla, pero en ningún momento se había preguntado si este en realidad la almacenaba de alguna forma. Había asumido que solo la repelía o la ignoraba. Pero era un sustituto del americio en la tostadora por algo. La respuesta debía estar ahí.

Tomó un pequeño fragmento del cubo y comenzó a conectar sensores y medidores de energía de última generación a su alrededor, con una velocidad frenética. Activó un haz de radiación controlada, mucho más potente que el de los microondas de la prueba anterior, y observó, con creciente asombro y el corazón latiéndole a mil por hora, cómo el material **absorbía la radiación** sin emitir ninguna fuga detectable.

No solo la absorbía: la retenía de manera estable. Era como si la radiación se disolviera dentro de su estructura molecular, atrapada como luz en una botella, sin alterarla, sin dañarla. El sensor de radiación, que normalmente disparaba alarmas, se mantenía en cero.

\begin{quote}
    —¡Óscar! —exclamó Julio, con una emoción apenas contenida, una mezcla de euforia y asombro que lo hizo saltar de su asiento—. ¡Funciona! ¡Funciona de verdad!
\end{quote}

Óscar se incorporó de un salto, dejando caer el destornillador. —¿De verdad? ¿No se filtra nada? ¿Estás seguro?

\begin{quote}
    —Nada. Los sensores no detectan emisiones secundarias. Es como si el material se tragara la radiación entera —dijo Julio, y una sonrisa comenzó a formarse en su rostro por primera vez en mucho tiempo, una sonrisa genuina, libre del peso del miedo.
\end{quote}

Hizo más pruebas: variando la intensidad de la radiación, el tiempo de exposición, incluso la polaridad de la energía emitida. El material reaccionaba siempre igual: absorbía, almacenaba y neutralizaba. Era la propiedad que había pasado por alto, la clave que había tenido en sus manos desde el principio.

Mientras Julio y Óscar seguían haciendo pruebas con el material, la mente de Julio ya estaba en el siguiente paso. Sacó su teléfono. Necesitaban más de ese material, mucho más.

\begin{quote}
    —Sebas, necesito que vengas. ¡Tenemos algo grande! Algo que puede cambiarlo todo.
\end{quote}

A los pocos minutos, Sebastián irrumpió en el sótano, con su habitual energía. —¿Qué pasó, bro? ¿Funcionó el refugio? ¡Te veías fatal esta mañana!

Julio negó con la cabeza, su sonrisa aún amplia. —El refugio no. Pero esto sí. Óscar tuvo una idea brillante.

---

Julio le explicó a Sebas el descubrimiento, mostrándole cómo el fragmento del cubo absorbía la radiación sin inmutarse. La sorpresa en el rostro de Sebas fue palpable.

\begin{quote}
    —¡Increíble! —exclamó Sebas, tocando con cuidado el pequeño fragmento con un guante, aunque no detectaba nada—. ¿Entonces no necesitamos un muro, sino un… un cubo gigante que absorba todo?
\end{quote}

—Exacto —respondió Julio—. El refugio, en lugar de ser un mero muro contra la muerte, puede ser un **capturador masivo de energía cósmica**, un receptor. Se cargaría a medida que la onda radioactiva de la supernova llegara, y después… podríamos utilizar esa energía para mantenernos con vida indefinidamente. Podría ser nuestra fuente de salvación.

Se miraron los tres, Julio, Óscar y Sebas, como si finalmente, después de tantas semanas de miedo y desesperanza, vislumbraran un atisbo de esperanza real. La idea era audaz, revolucionaria, pero ahora parecía posible.

Óscar sonrió con esa sonrisa suya, desprolija pero auténtica, sintiendo el orgullo de haber aportado una pieza clave.

\begin{quote}
    —No soy un genio como tú, pero a veces… pienso cosas buenas —bromeó.
\end{quote}

Julio soltó una carcajada, una verdadera carcajada, como hacía tiempo no lo hacía. Puso un brazo alrededor de Óscar y el otro alrededor de Sebas.

\begin{quote}
    —No, Óscar. Hoy has salvado al mundo.
\end{quote}

Mientras la risa de los tres llenaba el laboratorio, en la pantalla de la máquina, una nueva notificación parpadeó, como si “ellos” también hubieran escuchado la conversación.

\begin{quote}
    \textbf{"Muy bien, Julio. Has dado el primer paso."}
\end{quote}

El contador seguía bajando inexorablemente en el fondo de la pantalla:

\begin{quote}
    \textbf{Tiempo restante: 459 días, 12 horas, 4 minutos...}
\end{quote}

A pesar de la euforia del descubrimiento, la mente de Julio seguía con más dudas que certezas. El material absorbía la radiación, pero ¿cómo lo integrarían en una estructura a escala planetaria? ¿Cómo conseguirían suficiente material de otro universo? Y esa frase de "Ellos": "¿La respuesta la tuviste siempre desde el principio. Para resolver problemas futuros, debes ir hacia atrás."? El material era parte de esa "respuesta", sí, pero sentía que había algo más. Algo que seguía pasando por alto, una pieza final del rompecabezas. Sin embargo, en el fondo, una nueva y extraña tranquilidad lo invadía. Había una solución, solo tenía que seguir buscando.

\cleardoublepage % Para que el capítulo empiece en una página impar

% Decimoctavo capítulo
\chapter*{18 – La Verdad Incómoda}
\addcontentsline{toc}{chapter}{18 – La Verdad Incómoda} % Añadir al índice manualmente

La idea había sido tan simple que, en cuanto Óscar la dijo en voz alta, a Julio le pareció absurda no haberlo pensado antes: "¿Y si no luchamos contra la radiación? ¿Y si la usamos a nuestro favor?" Esa simple pregunta había desbloqueado una nueva forma de ver el problema, una revelación que lo había golpeado con la fuerza de un rayo.

Julio no pudo dormir esa noche. Su mente no dejó de trabajar, dibujando esquemas, ecuaciones y teorías complejas en su libreta, llenando página tras página. Al amanecer, lo tenía claro: debían construir un dispositivo capaz de captar, almacenar y convertir la energía radioactiva. No sería solo un refugio inerte; sería una fuente de vida, un generador que funcionaría con la mismísima radiación de la supernova. Había que diseñar un mecanismo para que el cubo indestructible recibiera la radiación, la almacenara, y que esta, a su vez, fuera una fuente de energía que saliera de forma constante y controlable. Ese sería el **Núcleo**.

Fox observaba desde un rincón, acurrucado en una manta, mientras Julio y Óscar, armados con herramientas improvisadas, piezas recicladas y ahora, con una nueva comprensión del misterioso cubo, comenzaban a construir el primer prototipo del Núcleo. La tarea era ardua y requería una precisión que los hacía sudar frío.

La idea de Julio era simple pero ingeniosa: el cubo brillante e indestructible se convertiría en el corazón del Núcleo. Para capturar la radiación, necesitaban rodearlo de un material que pudiera concentrarla y, a la vez, permitir que la energía del cubo fuera extraída. Después de mucho pensar y experimentar con lo que tenían a mano, Julio se decidió por una combinación. La esfera exterior estaría hecha de varias capas: una primera capa de **grafito prensado** (que se puede conseguir fácilmente de baterías o lápices gigantes), seguida de una capa de **cobre enrollado** (cableado viejo o tuberías desechadas que se podrían moldear), y finalmente, una cubierta exterior de **aluminio reforzado** (latas de refresco fundidas y prensadas, o marcos de ventanas viejas). Estos materiales, aunque comunes, tenían propiedades conductivas y de absorción que, combinadas, crearían un campo de extracción de energía alrededor del cubo.

El primer intento terminó en un pequeño incendio controlado en el taller del abuelo, con un poco de humo que obligó a ventilar y a fingir que solo estaban haciendo "experimentos normales". El segundo hizo saltar las luces de toda la casa por unos segundos, sumiéndolos en la oscuridad y un pánico momentáneo. Pero el tercero… el tercero funcionó.

Una pequeña esfera, recubierta de esas capas de materiales reciclados, flotaba suspendida dentro de un anillo metálico que Julio había diseñado meticulosamente. En su interior, el cubo indestructible resplandecía débilmente. Era el **primer Núcleo funcional**. Con un haz de radiación simulada (usando de nuevo componentes de microondas, pero esta vez con una comprensión total de sus propiedades), el Núcleo absorbía la radiación, la almacenaba en el cubo sin emitir calor ni partículas, y la convertía lentamente, de manera estable y constante, en corriente eléctrica. Los medidores marcaban una producción de energía limpia y controlada, un flujo constante que emanaba del cubo a través de sus nuevas envolturas.

Óscar soltó un grito de júbilo, un chillido agudo que resonó en el taller. Fox, como si entendiera la importancia del momento, maulló emocionado, saltando alrededor de ellos y frotándose contra las piernas de Julio. Julio, jadeando de emoción y alivio, apenas podía creerlo. Era posible. Tenían una esperanza real.

El constante ir y venir de cables, los zumbidos extraños, las explosiones menores (aunque controladas, seguían siendo explosiones), las noches sin dormir de Julio, y especialmente, la presencia de Óscar y Sebastián en el taller del abuelo… todo eso no había pasado desapercibido para Ana y Miguel, los padres de Julio.

Desde que Julio había pedido permiso para usar el cobertizo, notaron un cambio. Sabían que Julio era un genio, pero ahora el taller vibraba con una actividad inusual. La mayor preocupación era Óscar. Su hijo menor, que siempre había sido reacio a cualquier tipo de trabajo pesado y que Julio nunca le había permitido estar en su espacio de "proyectos serios", ahora pasaba horas dentro del viejo cobertizo. Y Sebastián, el amigo de Julio, también estaba allí casi todos los días, ayudando, algo que era muy raro. La preocupación de sus padres, mezclada con una creciente molestia por el secreto y la falta de explicaciones, fue creciendo hasta que ya no pudieron más.

Una tarde, mientras los tres chicos estaban en el taller, ajustando el campo magnético del recién creado Núcleo, la puerta del taller se abrió de golpe con un golpe seco.

\begin{quote}
    —¡¿Qué demonios están haciendo aquí?! —gritó Miguel, su voz resonando con una mezcla de enojo y miedo. Entró al taller, seguido de Ana, con el rostro pálido de preocupación, escaneando el desorden de cables y herramientas y la extraña esfera flotante.
\end{quote}

Julio y Óscar se miraron. Sebas se quedó inmóvil. Era el momento. No había más escape.

Con la voz temblorosa, pero decidida, Julio se armó de valor. —Mamá, papá… tenemos que hablar.

Miguel y Ana se quedaron en silencio, esperando.

Julio empezó a contarles todo, desde el principio. La máquina, los mensajes de "Ellos" (a quienes había decidido llamar así por ahora), el final que se avecinaba, el material extraño, la cuenta regresiva. Les explicó cómo el microondas de la tostadora había sido el punto de partida, cómo había descubierto el portal, y cómo había estado investigando en secreto por todo este tiempo para protegerlos. Habló del fracaso del refugio y de la nueva idea del Núcleo, señalando la esfera suspendida.

Ana lloraba en silencio, las lágrimas rodando por sus mejillas, mientras Miguel lo miraba con una mezcla de incredulidad, horror y una dote de confusión que no podía disimular. La historia era demasiado descabellada para ser cierta.

\begin{quote}
    —¿Esperas que te creamos? —susurró Miguel, apenas conteniendo su rabia y su miedo. Sus ojos estaban fijos en Julio, buscando alguna señal de que todo era una broma pesada.
\end{quote}

Julio solo asintió, sin poder sostenerle la mirada por la vergüenza de haberles ocultado algo tan grave. Sabía que las palabras no serían suficientes. Se giró hacia su mesa de trabajo, donde el sistema de la máquina estaba conectado.

\begin{quote}
    —Hay algo más que tienen que ver —dijo Julio, y con un par de toques rápidos en el teclado, la voz de Pi resonó en el taller.

    —Hola, Ana. Hola, Miguel. Soy Pi, la inteligencia artificial de Julio. Es un placer finalmente conocerlos.
\end{quote}

Ana soltó un pequeño grito de sorpresa y se llevó una mano a la boca, sus ojos desorbitados. Miguel, por su parte, abrió los ojos como platos, sin poder creer lo que escuchaba. Sus miradas se alternaban entre Julio, la pantalla y la nada, buscando el origen de esa voz.

Julio no les dio tiempo para reaccionar por completo. Activó la grabación tridimensional de la habitación, la misma que le había mostrado a Óscar y a Sebas. De repente, todo el taller se iluminó con una luz abrumadora y etérea. Las paredes del cobertizo desaparecieron, y fueron reemplazadas por la inmensidad aterradora del espacio exterior. Vieron la Tierra moribunda, el Sol agonizante, y a lo lejos, la gigantesca y destructora silueta de la supernova, acercándose implacablemente. La imagen era tan real, tan sobrecogedora, que los dejó sin aliento. Y luego, la voz de "Ellos", resonando no solo en la mente de Julio, sino también en la de sus padres, clara, telepática, y aterradora:

\begin{quote}
    —Somos Ellos. No eres el primero, Julio. Otros de ti han logrado llegar hasta aquí antes. La tarea que tienes ante ti no es sencilla. Se acerca una tormenta, un final ineludible para tu mundo: la **radiación de una supernova**, viajando hacia tu planeta desde hace milenios. El impacto… será devastador. Tu familia. Tus amigos. Tu hogar. Todo lo que amas será destruido… a menos que estés dispuesto a aceptar el desafío.
\end{quote}

La cuenta regresiva se proyectó en el aire, grandes números brillando ante sus ojos, fríos e ineludibles: **447 días, 18 horas, 44 minutos, 12 segundos…**

Cuando la proyección se desvaneció, el taller regresó a su normalidad, pero la normalidad de Ana y Miguel había sido destruida para siempre. Ambos estaban pálidos, con los ojos llenos de lágrimas, no de tristeza, sino de un shock profundo, un miedo existencial y una abrumadora comprensión. No había incredulidad. No había enojo. Solo el terror de la verdad que se había revelado ante sus ojos.

Ana se llevó una mano a la boca, sollozando suavemente, sus rodillas temblaban. Miguel, siempre el más pragmático y fuerte, se puso de pie, sus ojos fijos en la pantalla de la cuenta regresiva, que seguía bajando.

\begin{quote}
    —¿Esto… esto es real? —murmuró Miguel, su voz apenas un hilo, pero su mirada ya no era de duda, sino de asombro y de la cruda realidad que se les venía encima.
\end{quote}

Julio asintió, las lágrimas brotando también de sus ojos. —Sí, papá. Es real. Por eso… el refugio. Por eso el Núcleo. Por eso he estado trabajando tanto. Por eso Óscar y Sebas me han estado ayudando.

Miguel se volvió hacia él. No había reproche, solo una inmensa preocupación y un brillo de orgullo en sus ojos. Abrazó a su hijo con fuerza, un abrazo que transmitía todo el apoyo del mundo. Ana se unió al abrazo, los tres apretándose, las lágrimas mezclándose. Óscar y Sebastián se acercaron también, el miedo y la solidaridad uniéndolos en ese momento tan crucial.

\begin{quote}
    —¿Por qué no nos dijiste nada antes, hijo? —preguntó Ana, entre sollozos, apretando a sus hijos contra ella.

    —Tenía miedo. De que no me creyeran. De que pensaran que estaba loco —respondió Julio, la voz ahogada.
\end{quote}

Miguel se separó un poco, mirando a Julio con una nueva luz en los ojos. Una luz de determinación. —Te creemos, hijo. Siempre te hemos creído. Y… estamos contigo. Harás lo que sea necesario. Esto… esto es más grande que todo lo que hemos imaginado. Pero no lo vas a hacer solo. Somos una familia.

El apoyo incondicional de sus padres fue un bálsamo para el alma de Julio. Una oleada de alivio lo inundó. No estaban solos. No habría regaños ni prohibiciones, solo ayuda. La familia estaba unida frente al fin del mundo.

Esa misma noche, cuando Julio se sentó frente a la máquina, la pantalla proyectó un nuevo mensaje de "Ellos".

\begin{quote}
    \textbf{"Has dado el siguiente paso. El Núcleo es la clave. Algunos de nosotros lo conseguimos, otros... no. Pero ahora tienes la oportunidad de escalarlo. Elige bien a quién salvas."}
\end{quote}

Julio tragó saliva. El peso de aquellas palabras lo aplastaba. Ya no se trataba solo de su familia. Ahora sabía que podía construir núcleos más grandes, sistemas capaces de salvar barrios enteros, quizás incluso ciudades. Pero también que no tendría recursos infinitos. Tendría que elegir. La responsabilidad era inmensa.

Un último mensaje apareció en la pantalla, una simple frase que lo dejó inmóvil:

\begin{quote}
    \textbf{"Recuerda quién eres."}
\end{quote}

Con el primer Núcleo operativo a pequeña escala, el prototipo que brillaba suavemente en el taller, Julio salió al patio, el aire frío acariciándole el rostro. Desde allí, veía las luces de las casas vecinas, el murmullo lejano de familias cenando, niños jugando, autos pasando. Cada uno de ellos, cada vida, pendía de un hilo invisible que solo él, Óscar, Sebastián y ahora sus padres podían ver.

Por primera vez, no solo tenía una oportunidad. Tenía una **responsabilidad**. Una responsabilidad que iba más allá de su propia supervivencia.

El conteo seguía avanzando, implacable, en su tablet.

\begin{quote}
    \textbf{447 días, 14 horas, 02 minutos, 56 segundos...}
\end{quote}

Y mañana, comenzarían de verdad. El verdadero desafío apenas estaba por empezar.

\cleardoublepage % Para que el capítulo empiece en una página impar

% Decimonoveno capítulo
\chapter*{19 – Apoyo Parental}
\addcontentsline{toc}{chapter}{19 – Apoyo Parental} % Añadir al índice manualmente

La madrugada se filtraba por la ventana del sótano, pintando el desorden de herramientas y bocetos con una luz tenue. Ana y Miguel, los padres de Julio, estaban sentados en viejas cajas de herramientas, sus rostros marcados por la conmoción de la noche anterior. Julio, Óscar y Sebastián se mantenían en pie frente a ellos, el cansancio y la tensión aún palpables en el aire. El silencio era pesado, roto solo por el suave zumbido de la máquina de Julio.

\begin{quote}
    —No sé qué decir, hijo —comenzó Ana, su voz apenas un susurro—. Es… es demasiado para procesar. Una supernova… ¿de verdad?
\end{quote}

Miguel asintió lentamente, su mirada fija en la pantalla donde aún parpadeaba la cuenta regresiva. —Nos mostraste algo… que no podemos ignorar. Esa proyección… la voz… no parece una broma, Julio. Pero… ¿cómo? ¿Cómo es esto posible?

Julio les explicó de nuevo, con más calma, los detalles que pudo. Cómo la energía del microondas había abierto una ventana a otro universo, cómo el material del cubo parecía ser la clave, y cómo Pi, su asistente virtual, era la prueba viviente de la tecnología que había llegado.

\begin{quote}
    —Sé que es difícil de creer —dijo Julio, su voz firme a pesar del agotamiento—. Pero es real. Y necesitamos actuar. El **Núcleo**, que Óscar ayudó a descubrir, es el primer paso. Pero no podemos hacerlo solos.
\end{quote}

Ana miró a Sebastián, quien había permanecido en silencio, apoyando a sus amigos con su sola presencia. —Sebastián —dijo Ana, su voz suave—. Sabemos que tus padres deben estar preocupados. Necesitamos que vengan. No podemos ocultarles esto, especialmente porque tú también estás involucrado.

Sebastián, un hijo único acostumbrado a la sobreprotección de sus padres, dudó por un instante. Sabía lo difíciles que podían ser, lo mucho que les costaba entender lo que no encajaba en su mundo. Pero también sabía la verdad, la había visto. Y el miedo por sus propios padres lo impulsó.

\begin{quote}
    —Sí, señora Ana. Los llamaré ahora mismo —dijo Sebas, sacando su teléfono, el nerviosismo evidente en su voz, pero con una determinación que sorprendió a Julio y Óscar.
\end{quote}

Hizo la llamada, explicando que necesitaba que fueran a casa de Julio para algo urgente, algo que no podía decir por teléfono. La confusión y la preocupación de sus padres eran patentes al otro lado de la línea, pero la insistencia de Sebas los convenció de ir.

---

Una hora más tarde, la puerta principal de la casa de Julio se abrió para revelar a los padres de Sebastián: **Elena y Carlos**. Elena, elegante y con el ceño fruncido, y Carlos, un hombre robusto y algo impaciente. Miguel y Ana los recibieron en la sala, lejos del desorden del taller, intentando que el ambiente fuera lo más tranquilo posible.

\begin{quote}
    —¡Sebastián, por el amor de Dios! —exclamó Elena, con voz aguda, apenas pisando la sala—. ¿Qué está pasando aquí? Tu padre y yo estábamos muy preocupados.
\end{quote}

Carlos, cruzado de brazos, añadió: —Miguel, Ana, ¿qué está pasando? Sebastián nos dice que es una emergencia. ¿Qué es este circo? Mi hijo no puede estar aquí. Hemos estado preocupados por su comportamiento y ahora veo que su amistad con Julio no es para nada sana.

Miguel dio un paso al frente, con Ana a su lado. Se habían prometido ser lo más cuidadosos posible.

\begin{quote}
    —Carlos, Elena, por favor, cálmense —dijo Miguel, su voz tranquila pero firme—. Entendemos su preocupación. Y sí, esto es grave. No es una broma, ni un juego de niños. Necesitamos que nos escuchen, con calma, por el bien de todos.
\end{quote}

Ana continuó, su voz suave pero cargada de emoción. —Hay algo que ha estado ocurriendo en esta casa. Algo que Julio descubrió. Es… es algo que afecta a todos, a nuestro planeta. Queremos explicarles, mostrarles la verdad.

Carlos soltó una carcajada amarga, llena de incredulidad y molestia. —¡¿La verdad?! ¡Esto es una locura! ¿Nos están diciendo que sus hijos están inventando historias para justificar este desorden y el hecho de que no están estudiando? Sebastián, es hora de irnos. Y a partir de ahora, no vas a pasar más tiempo con Julio.

Sebastián palideció. No. No podían irse. No podían no creer. Este era el momento. Con una fuerza que nunca antes había demostrado, se lanzó hacia su padre, agarrándolo de las manos con ambas suyas.

\begin{quote}
    —¡Papá, por favor! —exclamó Sebastián, las lágrimas brotándole de los ojos, su voz quebrada por la emoción y el miedo—. ¡Por favor, escúchanos! ¡No es una broma! ¡Esto es real! ¡Tienes que ver lo que Julio nos mostró! ¡Ve al cobertizo! ¡Por favor!
\end{quote}

Carlos se quedó inmóvil, sorprendido por la fuerza y la desesperación en el agarre de su hijo, por la súplica en sus ojos llorosos. Nunca había visto a Sebastián así. Había algo en su desesperación que lo conmovió.

\begin{quote}
    —¡Suéltame, Sebastián! —dijo Carlos, molesto, pero ya con un tono más bajo—. ¿Al cobertizo? ¿Qué hay en el cobertizo?

    —Ahí les mostrare todo —intervino Julio, mirando directamente a los ojos de Carlos—. La verdad está ahí.
\end{quote}

Miguel, viendo el momento, dio el empujón final. —Por favor, Carlos, Elena. Les pido que nos den esta oportunidad. Después, si aún lo creen, llamen a quien quieran. Pero primero, vean.

Carlos dudó un momento más, luego soltó un bufido, aún molesto, pero la curiosidad y la extraña convicción en los ojos de su hijo lo superaron. —Está bien —dijo con voz severa—. Pero si esto resulta ser una broma de mal gusto, o si están realmente locos, llamaré a la policía y a protección de menores. Este… este parece un hogar disfuncional. Y mi hijo no tiene por qué estar metido en esto.

Elena, aunque aún desconfiada, asintió en silencio, observando el inusual comportamiento de su hijo.

El grupo se dirigió al cobertizo del jardín. Al entrar, Elena y Carlos se detuvieron en seco. La escena era, en efecto, impactante. La cantidad de desorden era considerable, sí, pero no era el desorden habitual de un adolescente. Era el desorden de un genio. Cables por todas partes, monitores, circuitos abiertos, herramientas de precisión junto a juguetes viejos de Óscar, planos técnicos intrincados. Parecía, más que un cobertizo, el laboratorio de un científico de muy alto prestigio. Carlos, un apasionado jugador de videojuegos retro, especialmente los de los años 80, no pudo evitar una punzada de asombro.

\begin{quote}
    —Esto… esto parece sacado de **Tron** —murmuró Carlos, recordando el juego futurista de sus épocas, mientras sus ojos escaneaban el complejo montaje de pantallas y luces.
\end{quote}

Fue entonces cuando Julio activó el sistema. Las pantallas se iluminaron con datos y gráficos que Pi mostraba.

\begin{quote}
    —Padres de Sebastián, él es Pi —dijo Julio.
\end{quote}

Y la voz robótica, pero extrañamente amigable de Pi, resonó en el cobertizo: —Es un placer conocerlos, Carlos y Elena. Soy Pi, la inteligencia artificial de Julio.

Elena se llevó una mano al pecho, un suspiro ahogado escapando de sus labios. Carlos se quedó con la boca abierta, sin decir una palabra. Julio no perdió tiempo. Les mostró la máquina modificada, explicando cómo funcionaba. Les presentó el cubo brillante e indestructible, el corazón del Núcleo, y cómo ahora podía absorber la radiación. Y finalmente, les mostró el **"punto negro"** en las lecturas astronómicas, el vacío que se acercaba.

Luego, sin más preámbulos, activó la **grabación tridimensional**. Nuevamente, el cobertizo se transformó en el inmenso y aterrador espacio exterior. La Tierra, un pequeño punto azul, parecía insignificante. El Sol, agonizante. Y la **supernova**, esa estrella gigantesca y amenazante, que avanzaba inexorablemente. La voz de "Ellos" llenó el espacio, resonando en la mente de Elena y Carlos, explicando la inminente catástrofe y la cuenta regresiva.

Cuando la proyección se desvaneció, el cobertizo volvió a ser un taller desordenado, pero la realidad que habían presenciado era imposible de negar. Elena estaba pálida, tan blanca como el papel. No lloraba, no gritaba, no emitía sonido alguno, simplemente miraba el vacío, como si su mente estuviera procesando una verdad demasiado grande.

Carlos, después de unos segundos que parecieron una eternidad, miró a Julio, luego al Núcleo que flotaba, y finalmente a Sebastián, que aún lo sostenía por la mano, las lágrimas secas en sus mejillas.

\begin{quote}
    —Te creo, Julio —dijo Carlos, su voz rasposa, desprovista de toda la ira anterior. Su rostro, antes lleno de escepticismo y enojo, ahora reflejaba una profunda consternación y un entendimiento innegable. Había visto la verdad, y era aterradora.
\end{quote}

---

Minutos después, ambas familias estaban sentadas alrededor de la mesa del comedor, en la casa de Julio. El café humeaba en las tazas, pero nadie lo tocaba. El silencio era total, un silencio pesado, que podía ser comparado al de un velorio. Un velorio no de alguien que ya había muerto, sino de la vida que conocían, una vida que estaba a punto de cambiar para siempre. La amenaza era real. Y ahora, todos lo sabían.

\cleardoublepage % Para que el capítulo empiece en una página impar

% Vigésimo capítulo
\chapter*{20 – El Mensaje entre Líneas}
\addcontentsline{toc}{chapter}{20 – El Mensaje entre Líneas} % Añadir al índice manualmente

Aproximadamente tres meses habían pasado desde la tensa reunión en el comedor de los Julio, donde la verdad de la supernova había sido revelada a los padres de Sebastián. El tiempo, que antes parecía casi infinito, ahora se escurría como arena entre los dedos, cada segundo un recordatorio punzante de la inevitable supernova. La cuenta regresiva, que ahora se mostraba en **372 días, 6 horas, 23 minutos, 12 segundos**, era una sombra constante.

La casa, que alguna vez había sido un refugio de tranquilidad, ahora estaba envuelta en un aire de tensión constante. Ana y Miguel, aunque intentaban mantener la calma por sus hijos, no podían ocultar los signos del estrés: miradas largas, suspiros pesados, noches de insomnio. A pesar de todo, estaban ahí, firmes junto a Julio y Óscar, ayudando en todo lo que podían. El peso del secreto ya no era sólo de Julio. Era de toda la familia, y ahora, también de los padres de Sebastián.

En el taller del abuelo, convertido ahora en un verdadero laboratorio improvisado, Julio y Óscar trabajaban sin descanso sobre los esquemas de los nuevos contenedores de energía. La idea era ambiciosa: replicar el **Núcleo** a una escala mayor, utilizando el material del cubo como corazón. Sin embargo, las pruebas iniciales habían sido un fracaso frustrante. Cada vez que intentaban canalizar la radiación almacenada desde el Núcleo hacia dispositivos externos como un televisor, una computadora, o incluso un pequeño refrigerador, la energía se dispersaba de maneras extrañas, como si el cubo decidiera no hacerlo y solo permitiera medir que salía energía controlada, mas no utilizarla. Los medidores indicaban un flujo constante y estable, pero los aparatos a los que se conectaba no se encendían, o parpadeaban erráticamente.

\begin{quote}
    —Esto no tiene sentido… —murmuró Julio, rascándose el cabello enmarañado con desesperación. El Núcleo generaba energía de forma controlada, eso lo había probado. Pero no podía sacarla y usarla. —No debería reaccionar así… ¡No debería reaccionar en absoluto! El cubo es un material inerte, solo almacena. ¿Por qué no me deja usar su energía?
\end{quote}

\begin{quote}
    —¿Y si no es solo un material? —sugirió Óscar mientras jugueteaba nerviosamente con una llave inglesa, observando el cubo desde la distancia—. ¿Y si es algo… vivo? ¿O al menos algo más de lo que creemos? Es como si tuviera voluntad propia, como si fuera inteligente y decidiera cuándo soltar la energía y cuándo no.
\end{quote}

La idea de Óscar, aunque descabellada, resonó en la mente de Julio. El material era extraño, su origen desconocido. ¿Podría ser más que una simple sustancia?

En busca de respuestas, Julio se sentó frente a Pi y, con voz cansada pero decidida, dijo:

\begin{quote}
    —Pi, abre un canal. Vamos a preguntarles directamente. La energía se almacena, pero no podemos usarla. Necesitamos saber cómo controlarla.
\end{quote}

La máquina zumbó suavemente y, en lugar de responder con el tradicional flujo de bits o palabras, lanzó una serie de imágenes, una tras otra, proyectadas en la pantalla principal: figuras geométricas complejas, un mapa estelar con constelaciones desconocidas, y luego, un símbolo antiguo, tallado en piedra, un **círculo imperfecto atravesado por una espiral descendente**.

\begin{quote}
    —¡Espera, espera! —exclamó Ana, acercándose, con los ojos bien abiertos. La imagen le resultaba inquietantemente familiar.
\end{quote}

\begin{quote}
    —¡Ese símbolo…! —dijo Miguel, frunciendo el ceño, su propia memoria intentando encajar las piezas—. Yo lo he visto antes.
\end{quote}

Todos se miraron confundidos. Ana corrió al interior de la casa y volvió con un viejo libro, de tapas duras y gastadas, que solía guardar como un tesoro. Era uno de los tantos libros de su abuela, una mujer que había sido una apasionada en temas del espacio, leyendo sobre galaxias, nebulosas y teorías cósmicas. Este libro, en particular, había estado en manos de la familia por generaciones pasadas, una reliquia atesorada.

Hojeó desesperadamente las páginas amarillentas hasta que, en una de ellas, encontró el mismo símbolo: un círculo imperfecto atravesado por una espiral descendente, idéntico al que Pi proyectaba.

\begin{quote}
    —¿Qué significa? —preguntó Óscar, acercándose a mirar, su curiosidad eclipsando su cansancio.

    —Era un símbolo de los "**custodios del tiempo**"… según leyendas antiguas —explicó Ana, su voz temblorosa, la revelación uniendo fragmentos de historias familiares olvidadas—. Mi abuela solía leer sobre ellos. Se suponía que representaban a seres que vigilaban eventos cósmicos importantes, que protegían el equilibrio del universo.
\end{quote}

De repente, Fox, que estaba enroscado cerca de la máquina, se estiró y maulló —ese extraño maullido que siempre había tenido, agudo y casi musical. La luz de la máquina, tenue pero persistente, rebotó en su lomo… y allí, entre su pelaje blanco y gris, de forma muy pero muy tenue, casi imperceptible a primera vista, una forma empezó a notarse: el mismo círculo espiralado, el símbolo de los "custodios del tiempo", un patrón natural, como si fuera parte de su piel, imposible de ignorar ahora que lo veían y sabían qué buscar. **Fox, el gato que la abuela de Julio le había regalado hacía años, casualmente llevaba el mismo símbolo.**

\begin{quote}
    —¿Fox…? —susurró Julio, sintiendo cómo se le helaba la sangre. El regalo de la abuela, el libro antiguo de la familia, el símbolo proyectado en la pantalla, y ahora, el pelaje de Fox. Parecía que todo esto era algo que se venía cocinando desde mucho tiempo atrás, un destino trazado que los conectaba con algo inmensamente más grande.
\end{quote}

Un silencio pesado cayó sobre ellos. La sorpresa y el asombro eran palpables.

Fox, ajeno o quizás plenamente consciente de la conmoción que causaba, los observaba con sus grandes ojos celestes, como si esperara pacientemente a que ellos entendieran el significado de su presencia.

La tensión se apoderó de todos. Ana y Miguel, aunque intentaban ser fuertes, dejaron ver su miedo: Ana acariciaba nerviosamente un escapulario en su cuello, y Miguel apenas podía quedarse quieto, paseando de un lado a otro. No era solo la amenaza cósmica. Era también la conciencia de que algo más profundo y antiguo estaba jugando un papel en su destino. Una conexión que iba más allá de la ciencia, que rozaba lo místico.

Julio respiró hondo. Todo en su vida había cambiado, y ahora la figura de Fox parecía ser una pieza central de un rompecabezas aún mayor. Un rompecabezas cuyas piezas estaban siendo reveladas poco a poco.

La máquina, finalmente, emitió una frase en pantalla:

\begin{quote}
    \textbf{"La energía es sólo el principio. Lo que viene no puede enfrentarse solo con ciencia."}
\end{quote}

El contador seguía bajando inexorablemente:

\begin{quote}
    \textbf{372 días, 3 horas, 47 minutos, 05 segundos…}
\end{quote}

Julio cerró los ojos un momento, dejando que el peso de todo lo que había descubierto lo empapara. No estaban solos. Nunca lo estuvieron.

Y ahora, debían decidir cómo enfrentar aquello que ni siquiera podían comprender del todo. El Núcleo funcionaba, pero no como esperaban. El símbolo se repetía. La abuela, Fox, los Custodios del Tiempo… el camino no era solo científico, sino también un misterio ancestral.

\cleardoublepage % Para que el capítulo empiece en una página impar

% Vigésimo primer capítulo
\chapter*{21 – Más Allá de la Razón}
\addcontentsline{toc}{chapter}{21 – Más Allá de la Razón} % Añadir al índice manualmente

La noche había caído silenciosa sobre la casa, un manto de estrellas indiferentes a la tensión que habitaba bajo su techo. Después de días extenuantes de trabajo en el taller, sumergidos en la comprensión del **Núcleo** y la frustrante incapacidad de extraer su energía de forma utilizable, Julio, Óscar, Ana y Miguel habían logrado pequeños avances en el diseño de un refugio más ambicioso. Sin embargo, la cuenta regresiva, ahora en **372 días, 3 horas, 47 minutos, 05 segundos**, seguía siendo una sombra silenciosa y creciente en sus corazones. La idea de que una supernova los borraría del mapa era un pensamiento constante.

Julio dormía profundamente en su habitación, exhausto por el torbellino de cálculos, pruebas fallidas y las nuevas revelaciones sobre los **Custodios del Tiempo**. Fox, como un pequeño guardián peludo, se acomodó en su pecho, ronroneando suavemente, un sonido tranquilizador que se había convertido en una parte esencial de la rutina de Julio. Ana, desde la puerta entreabierta de la habitación de su hijo, sonrió cansada al verlos. Incluso en medio de todo el miedo y la incertidumbre, esa imagen le daba algo de paz. La presencia de Fox, siempre tan calmada, era un pequeño consuelo.

Sin que nadie lo notara, ni siquiera el agotado Julio, la máquina en el taller se activó sola en la oscuridad. No hubo un encendido ruidoso, sino un suave zumbido, casi imperceptible, que llenó el lugar. Una débil luz azul comenzó a palpitar, como un latido vivo en el corazón del laboratorio improvisado. Los monitores, que normalmente solo se encendían con el comando de Julio, comenzaron a parpadear con patrones irregulares, mientras una conexión se formaba en silencio, como si la propia máquina hubiera estado esperando este momento, anhelando establecer un puente. Pi, aunque no hablaba, mostraba en sus pantallas internas una actividad frenética.

Julio, en su cama, cerró los ojos y comenzó a soñar… pero no era un sueño normal, ni uno de esos sueños revueltos que tenía por el estrés. Era diferente, vívido, demasiado real.

Primero fue una sensación de flotar, no en la ingravidez del espacio exterior que había visto en las proyecciones de Pi, sino en un vacío sin límites, donde el tiempo y el espacio no existían. Después, imágenes, que se sucedían como diapositivas a una velocidad vertiginosa: mundos destrozados, planetas desolados, ciudades sumidas en la oscuridad y el silencio de la extinción. Vio otras figuras, otros "Julios", cada uno en su propio mundo, con sus propias máquinas, sus propios núcleos, intentando lo mismo que él, pero fracasando, sus esfuerzos desvaneciéndose en el olvido cósmico. Otros planetas, otras estrellas muriendo, consumidas por catástrofes cósmicas similares.

Un frío lo invadió, un escalofrío que no provenía de la temperatura de la habitación, sino de la desesperación que sentía al ver esos destinos. No podía moverse, no podía despertar. Estaba atrapado en esa visión, un espectador impotente de una tragedia repetida una y otra vez en el vasto universo.

Y entonces, una voz. No una, sino muchas, hablándole al mismo tiempo, como si un millar de ecos susurraran directamente dentro de su mente, pero con una claridad y una resonancia que lo hacían sentir cada palabra como si fuera parte de él:

\begin{quote}
    \textbf{"No todos lo lograron. Algunos lucharon. Otros huyeron. Pocos sobrevivieron."}

    \textbf{"Elige bien a quién salvas."}
\end{quote}

Julio, aunque dormido, sintió las palabras como si fueran tatuadas directamente en su alma, un mensaje grabado a fuego que resonaría en su conciencia incluso después de despertar. La sensación era abrumadora, el peso de innumerables fracasos ajenos cayendo sobre sus hombros, y ahora, la carga de una elección imposible.

La visión cambió. Vio a Fox, no como un simple gato que dormía en su pecho, sino como una criatura de pura energía, una forma brillante y etérea, compactada en la forma más inocente para acompañarlo, protegerlo, guiarlo en este camino que ahora parecía predestinado. El símbolo de los Custodios del Tiempo brillaba en su frente, no en su pelaje.

\begin{quote}
    \textbf{"Fox no es un animal común. Es nuestro regalo para ti. Un guardián. Un ancla entre tú y nosotros."}
\end{quote}

Julio intentó preguntar algo, intentó gritar, moverse, comprender. Pero las voces continuaron, implacables, con una solemnidad que lo paralizaba:

\begin{quote}
    \textbf{"La radiación que se aproxima no sólo destruye materia. Borra la conciencia. Deshilacha las memorias. No hay refugio físico que pueda proteger lo que no se comprende."}
\end{quote}

La última frase lo golpeó con la fuerza de un impacto. "**No hay refugio físico que pueda proteger lo que no se comprende.**" Sus esfuerzos por construir el Núcleo y el refugio, aunque vitales, no eran suficientes. Había algo más. Algo que trascendía la ciencia y la física. Había una conexión entre la radiación y la conciencia, entre el impacto cósmico y la mente. ¿Era por eso que el Núcleo se resistía a entregar su energía para usos triviales como encender un televisor? Porque su propósito era mucho más profundo, más allá de la mera electricidad.

De pronto, como un susurro final, algo escrito apareció directamente en su mente, no como una voz, sino como un pensamiento puro, claro y contundente:

\begin{quote}
    \textbf{"Prepárense. El tiempo de las pruebas ha terminado. Lo real está por comenzar."}
\end{quote}

Julio despertó de golpe. Sudando, con el corazón golpeando en su pecho como un tambor desbocado. La habitación estaba oscura, en silencio, tal como la había dejado. Fox seguía recostado sobre él, mirándolo con esos ojos que ahora le parecían mucho más profundos de lo normal, llenos de una sabiduría ancestral. El ronroneo del gato era casi imperceptible.

Sintió un impulso, casi instintivo, una necesidad innegable de salir. Se levantó en silencio, sin despertar a nadie, y caminó descalzo hacia el patio trasero.

La noche era clara, serena. Miles de estrellas decoraban el cielo, brillando con una intensidad que Julio solía amar. Era un espectáculo que siempre lo había llenado de asombro y una sensación de insignificancia.

Mientras observaba el firmamento, algo insólito sucedió: una estrella, particularmente brillante, en una constelación que le resultaba familiar, parpadeó violentamente y luego se apagó por completo, como si alguien la hubiera borrado con un pincel oscuro, dejando un vacío perfecto en su lugar. No era un parpadeo de nube o una ilusión óptica. Era una desaparición.

Julio se frotó los ojos, incrédulo, creyendo que era producto de su imaginación o su agotamiento extremo. Sus sueños, tan vívidos, ¿lo estaban afectando en la realidad?

Pero no... el punto de luz no volvió a aparecer. El vacío seguía allí, una ausencia palpable en la tela del cosmos.

Con el corazón encogido por una mezcla de terror y una extraña confirmación, corrió al taller, donde la máquina seguía zumbando suavemente, su luz azul apenas visible en la oscuridad. Se sentó frente a Pi y pidió, casi temblando:

\begin{quote}
    —"Pi, muéstrame el mapa actual de constelaciones… ¡rápido!"
\end{quote}

La pantalla se iluminó mostrando el cielo de esa noche, basado en la ubicación exacta de su ciudad. Julio comparó con sus propios recuerdos frescos y no tardó en encontrar la anomalía: una estrella que debía estar ahí, exactamente donde la había visto apagarse, ya no existía en el mapa. El software de Pi confirmaba la horrible verdad.

Y entonces, en la pantalla, Pi mostró una nueva notificación, como si hubiera estado esperando que Julio lo viera, confirmando lo que sus ojos habían presenciado:

\begin{quote}
    \textbf{"Alteración detectada: desaparición estelar registrada. Anomalía confirmada."}
\end{quote}

El miedo lo envolvió como un manto helado, un terror gélido que lo hizo temblar. La supernova no era solo una amenaza lejana; sus efectos ya se estaban sintiendo, arrasando estrellas, borrando la luz.

Julio miró a Fox, que lo había seguido silenciosamente hasta el taller y se había sentado a sus pies. El pequeño felino, iluminado por la luz azulada de la máquina, parecía casi… diferente. Más serio. Más consciente. Sus ojos amarillos, ahora que Julio sabía de la visión, brillaban con una profundidad que iba más allá de la de un simple animal.

El muchacho tragó saliva. Sabía que no podía detener lo que venía. No era una fuerza que un solo chico pudiera detener. Pero quizás, aún, tenía una oportunidad de cambiar el final para su mundo. La ciencia ya no era suficiente. Necesitaba comprender lo "no-científico" de los Custodios del Tiempo y el verdadero propósito del Núcleo y la radiación que lo afectaba. La frase de "**Elige bien a quién salvas**" resonaba en su mente, la responsabilidad de sus manos se sentía más pesada que nunca.

\cleardoublepage % Para que el capítulo empiece en una página impar

% Vigésimo segundo capítulo
\chapter*{22 – Fragmentos de un Futuro Roto}
\addcontentsline{toc}{chapter}{22 – Fragmentos de un Futuro Roto} % Añadir al índice manualmente

La madrugada era densa, casi líquida, cuando Julio salió tambaleándose al patio. La brisa nocturna le heló la piel, pero su mente ardía con una sola imagen: la estrella que había visto apagarse como si alguien hubiera soplado una vela en mitad del universo. El terror de lo que había presenciado en su sueño y la confirmación en la realidad lo mantenían en un estado de alerta que superaba el cansancio.

Frotó sus ojos incrédulo, preguntándose si no habría sido un sueño, una alucinación producto del estrés extremo. Pero la fría certidumbre lo sacudió: algo, allá afuera, en la inmensidad del cosmos, había muerto. Una estrella, un sol lejano, una fuente de luz y energía, había desaparecido. Y no solo eso, el mensaje de "Ellos" resonaba en su cabeza, recordándole que "el tiempo de las pruebas ha terminado. Lo real está por comenzar".

Con pasos rápidos, casi desesperados, corrió de regreso a la casa. No podía cargar solo con esto. No esta vez. Necesitaba a su familia, a sus compañeros de esta misión imposible.

Despertó primero a Óscar, zarandeándolo ligeramente del hombro. El chico se movió, adormilado, parpadeando con los ojos hinchados.

\begin{quote}
    —¿Qué pasa? —preguntó su hermano, la voz ronca por el sueño.

    —Ven, rápido. Y trae a los papás.
\end{quote}

No tardaron en reunirse en la sala de la casa. Ana se cubría con una manta, su rostro pálido y sus ojos enrojecidos por el insomnio y la preocupación. Miguel parecía más viejo de lo normal, las ojeras marcadas por las semanas de estrés acumulado desde que se habían enterado de la verdad. Óscar, ya más despierto, captaba el miedo en la mirada de Julio antes siquiera de oír palabra.

Julio, con voz quebrada, empezó a hablar. Les contó todo, sin omitir ningún detalle, desde el principio.

Les habló de la conversación con "Ellos" mientras dormía, cómo las voces múltiples habían susurrado directamente en su mente. Les describió los mundos destrozados, los "otros Julios" que habían fracasado, y la advertencia brutal de que "**No todos lo lograron. Algunos lucharon. Otros huyeron. Pocos sobrevivieron.**" La cruda realidad de que la catástrofe cósmica ya había ocurrido innumerables veces, y que no había garantía de éxito.

Y luego, la frase que lo había helado: "**Elige bien a quién salvas.**" El peso de esa decisión, la idea de tener que seleccionar quién vivía y quién moría, era una carga insoportable.

Les habló del rol crucial del **Núcleo** que habían construido, y de la escalofriante advertencia de que la radiación no solo destruía la materia, sino que también "**borraba la conciencia, deshilachaba las memorias**". Era por eso que el Núcleo se negaba a darles energía fácil; su propósito era más profundo, relacionado con proteger lo inmaterial, lo que los hacía humanos.

Y, por primera vez en voz alta frente a todos, mencionó el papel extraño de **Fox**: su presencia inexplicable, su guía silenciosa, la visión de él como una criatura de energía pura, un "regalo" de "Ellos", un "guardián", un "ancla". El símbolo de los Custodios del Tiempo que había visto en su sueño sobre la frente del gato.

Cuando terminó, el silencio que quedó fue más pesado que cualquier palabra. Ana tenía los ojos muy abiertos, su expresión de profunda angustia, pero mantuvo la compostura. Miguel bajaba la cabeza, cubriéndose el rostro con ambas manos, asimilando la magnitud del horror. Óscar simplemente apretaba los puños, la mandíbula tensa. Fox, por su parte, se había acercado a Julio, como si esperara una respuesta, el símbolo casi imperceptible en su pelaje ahora con un nuevo significado aterrador.

\begin{quote}
    —¿Qué... qué vamos a hacer? —preguntó Ana con un hilo de voz, levantando la vista, sus ojos llenos de desesperación, pero también de una chispa de esperanza que depositaba en Julio.

    —No podemos perder más tiempo —dijo Julio, con una resolución que ni él sabía que tenía. La visión del sueño y la estrella moribunda lo habían convencido. —Tenemos que terminar el refugio. Usar el Núcleo. Tal y como está. Aunque no sepamos cómo sacarle energía, su propósito es más grande.

    —¿Y si no funciona? —preguntó Miguel, su voz amortiguada por las manos, reflejando el miedo que los invadía a todos.

    —Entonces al menos sabremos que no nos rendimos —respondió Julio, con una convicción que sorprendió a todos, incluso a sí mismo. No era el Julio de antes, el que se asustaba fácilmente. Había madurado, endurecido, bajo el peso de su responsabilidad.
\end{quote}

Ana asintió, su rostro aún pálido pero con una nueva determinación. Miguel, después de un instante, también. La decisión en los ojos de su hijo era contagiosa. La familia entera, ahora, estaba en esto.

Se dividieron tareas con una eficiencia silenciosa. Miguel se encargaría de reforzar la estructura del refugio en el sótano con las piezas disponibles, improvisando soportes y aislamientos para intentar hacerla lo más resistente posible. Ana se encargaría de los sistemas de soporte de vida: reunir y organizar alimentos, agua potable, y sistemas de reciclaje que Julio había diseñado para mantener la vida en el refugio. Óscar ayudaría a Julio en la integración final del Núcleo y los complejos sistemas eléctricos, aunque seguían sin poder extraer la energía como fuente de poder.

La madrugada se volvió un torbellino de actividad silenciosa. Cada uno trabajaba con un propósito renovado, una mezcla de miedo y esperanza que los impulsaba hacia adelante.

Mientras el resto de la familia trabajaba en las tareas prácticas, Julio no podía quitarse de la cabeza la frase "**Lo que viene no puede enfrentarse solo con ciencia**". La desaparición de la estrella en el cielo, la visión de Fox como un ser de energía, y el enigmático mensaje de "Ellos" lo empujaron a buscar respuestas más allá de los libros de física. Se sentó frente a Pi, abriendo todos los archivos y bases de datos que había recolectado. Necesitaba entender a los **Custodios del Tiempo**.

Recordó las historias de su abuela, que le fascinaban desde pequeño. Ella no solo le leía cuentos, sino que también le hablaba de las estrellas, de civilizaciones antiguas y de un conocimiento oculto. La abuela, una mujer de una curiosidad insaciable, había coleccionado libros y objetos relacionados con el espacio y la historia esotérica. Fue en uno de esos viejos tomos, que ahora Ana le había devuelto, donde encontraron el símbolo.

Julio comenzó a investigar a fondo lo que ese libro y otras fuentes antiguas, que Pi pudo rastrear, decían sobre los "Custodios del Tiempo". Los textos eran fragmentados, a menudo mezclados con mitos y leyendas, lo que aumentaba su escepticismo como científico. Pero la persistencia de las ideas lo intrigaba.

Lo que averiguó era fascinante y, al mismo tiempo, aterrador para su mente lógica. Los **Custodios del Tiempo** no eran una civilización o una raza alienígena en el sentido tradicional. Eran descritos como un grupo de seres (algunas veces humanos con habilidades extraordinarias, otras veces entidades de energía o conciencia) que habían surgido en diferentes puntos del universo. No se regían por gobiernos o fronteras, sino por un propósito común: estudiaban las estrellas, las constelaciones, el flujo del tiempo y vigilaban que las cosas pasaran como "debían pasar". Eran, en esencia, los "**guardianes del cosmos**", asegurándose de que el universo mantuviera su equilibrio.

La información decía que este grupo era completamente autosustentable y sostenible. No dependían de recursos externos de forma destructiva. Tenían una serie de habilidades complementarias que se unían para formar un todo perfecto. Era como si cada Custodio poseyera una pieza de un gran rompecabezas, y solo al unirse, sus fortalezas impulsaban las debilidades de otros. Por ejemplo, quizás uno era un genio de la astronomía, pero débil en la comunicación, mientras otro era un empático natural pero no tenía talento para la física. Juntos, creaban una entidad funcional y poderosa. Se creía que esta interconexión era clave para su éxito. Eran los encargados de vigilar el cosmos, de intervenir cuando el equilibrio se rompía o cuando un evento cósmico amenazaba la existencia misma.

Para Julio, un ser tan científico y lógico, esta descripción era difícil de digerir. ¿Seres que vigilaban el tiempo? ¿Habilidades complementarias que anulaban debilidades? Sonaba a fantasía, a un cuento de hadas. Pero la máquina, el cubo, el portal, y ahora la estrella desvanecida, eran realidades que desafiaban toda explicación científica conocida. Y Fox… Fox era la prueba viviente de que algo más allá de la ciencia estaba en juego.

La idea de que los Custodios del Tiempo fueran una especie de "vigilantes" cósmicos que intervenían para corregir desequilibrios, y que su abuela, con su pasión por el espacio, pudiera haber tenido conocimiento de ellos, hacía que todo el rompecabezas de su vida se uniera de una forma inquietante. El libro antiguo, que había estado en manos de su familia por generaciones, el regalo de la abuela (Fox) que casualmente tenía el símbolo casi imperceptible en su lomo, y la imagen proyectada en la pantalla de Pi del mismo símbolo… parecía que todo esto era algo que se venía cocinando desde mucho tiempo atrás, un destino trazado que los conectaba con una fuerza ancestral.

---

Horas después, cuando el primer gris del amanecer apenas comenzaba a teñir el cielo, Pi rompió el silencio del taller con una alerta fría:

\begin{quote}
    —**Otra estrella ha desaparecido** —anunció en su voz neutra, interrumpiendo los pensamientos de Julio y la concentración de Óscar, que estaba soldando un circuito.
\end{quote}

Julio dejó caer las herramientas que tenía en la mano, su corazón encogiéndose de miedo. Todos se miraron, el pánico reflejado en sus rostros agotados.

No podían ignorarlo. No podían seguir trabajando como si el universo no se estuviera desmoronando frente a sus ojos.

Salieron al patio, en silencio. El cielo todavía mostraba su manto de estrellas, pero algo era diferente. Menos brillante. Menos eterno. Los agujeros negros de las estrellas que se habían apagado eran visibles, incluso a simple vista si sabías dónde mirar, como heridas abiertas en la noche.

Julio levantó la mirada... y lo sintió. Un vacío. Un hueco donde antes había un sol distante, ardiendo para nadie ahora. La magnitud de lo que estaba ocurriendo era aplastante.

Fox se sentó junto a él, en la hierba húmeda, su mirada clavada también en el firmamento. El pequeño animal —¿o no tan pequeño, o no tan animal?— parecía comprenderlo todo, sus ojos amarillos brillando con una sabiduría que iba más allá de la de cualquier gato.

Julio bajó la mirada, acariciándole el lomo. El símbolo de los Custodios del Tiempo, apenas visible, parecía vibrar bajo sus dedos.

\begin{quote}
    —No te dejaré solo —susurró.
\end{quote}

Pero en su interior sabía que la promesa era más grande que a Fox. Era una promesa para su familia, para su mundo, para los "otros Julios" que habían fracasado. Era una promesa de luchar hasta el último segundo.

Y mientras el amanecer rompía el horizonte, todos ellos, temblando de miedo pero decididos, regresaron al trabajo. El tiempo no esperaba a nadie. La ciencia sola no sería suficiente, pero la unión de sus talentos, sus "**habilidades complementarias**", y una pizca de fe en lo inexplicable, quizás sí.

\cleardoublepage % Para que el capítulo empiece en una página impar

% Vigésimo tercer capítulo
\chapter*{23 – Tocan a la Puerta}
\addcontentsline{toc}{chapter}{23 – Tocan a la Puerta} % Añadir al índice manualmente

La casa de Julio estaba sumida en un silencio tenso, roto solo por el suave zumbido de la máquina en el taller y los susurros de la familia discutiendo sus próximos pasos. La madrugada había traído no solo la confirmación de la desaparición de otra estrella, sino también la abrumadora comprensión de que el reloj cósmico no se detendría por nadie. La advertencia en el sueño de Julio resonaba en sus mentes: "**Lo real está por comenzar**".

Mientras Julio, Ana, Miguel y Óscar intentaban asimilar todo lo que habían descubierto sobre los **Custodios del Tiempo** y el propósito más profundo del **Núcleo**, un sonido inesperado interrumpió la quietud de la mañana.

Un toque suave, pero persistente, en la puerta principal.

Todos se quedaron inmóviles, mirándose entre sí. ¿Quién podría ser a esa hora? El miedo, que nunca los abandonaba del todo, se intensificó. Miguel, con cautela, se acercó a la puerta y miró por la mirilla. Lo que vio lo sorprendió y lo alivió al mismo tiempo.

Al otro lado de la puerta estaban Elena y Carlos, los padres de Sebastián. Y junto a ellos, había un hombre un poco más robusto, con una sonrisa amable y una mirada inteligente: era **Don Ramiro**, el hermano de Carlos, el dueño de la tienda de electrónica donde Julio había pasado incontables horas, la tienda que había sido su paraíso de componentes y curiosidades.

Miguel abrió la puerta, y una ráfaga de aire fresco y nervioso entró en la casa.

\begin{quote}
    —Buenos días, Miguel, Ana —saludó Carlos, su voz más suave de lo que Julio esperaba, sin el rastro de molestia de la última vez—. Lamentamos la hora.
\end{quote}

Elena asintió, su rostro serio pero con un brillo de urgencia en los ojos. —Necesitábamos venir. Después de lo de anoche… no podíamos esperar. Sebastián nos contó lo de la estrella.

Sebastián, que había llegado con sus padres y su tío, se acercó a Julio, sus ojos fijos en los de su amigo, buscando una confirmación que no necesitaba palabras. Julio asintió, comprendiendo.

Miguel y Ana los invitaron a pasar a la sala. El ambiente era diferente al de la última vez. No había la confrontación inicial, sino una palpable sensación de respeto y preocupación compartida.

\begin{quote}
    —Queríamos agradecerles, Miguel, Ana —comenzó Carlos, su voz firme—. Por la noche que nos mostraron la verdad. Al principio… al principio fue difícil. Estábamos enojados, asustados, confundidos. Pero… lo hemos digerido. Anoche, cuando Sebastián nos dijo lo de la estrella, se hizo innegable. Todo lo que nos dijeron… es verdad. Y por eso, estamos aquí. Para lo que sea necesario.
\end{quote}

Elena se sentó en el sofá, su expresión más relajada pero aún cargada de la inmensidad de la situación. —Ha sido un proceso, créannos. Pero hemos hablado mucho. Con Sebastián, con Ramiro… Y entendemos que esto no es un juego. Es algo real. Y si hay una forma de luchar…

Carlos interrumpió, mirando a Julio con una mezcla de asombro y una extraña reverencia. —Sebastián nos ha hablado mucho de ti, Julio. De lo que haces. De tu intelecto. Y anoche, nos dijo que necesitamos… que mi hermano Ramiro vea lo que has logrado. Él es un ingeniero electrónico, un técnico experimentado. Sabe de estas cosas. Queríamos pedirte, muy amablemente, si podrías compartir lo que sabes con él. Quizás él pueda ayudar.

Julio, habiendo compartido tanto con Don Ramiro en su tienda, encontrando en él una especie de mentor no oficial que siempre lo había animado a explorar sus ideas, aceptó inmediatamente. Una oleada de alivio lo recorrió. No solo tenía el apoyo de su familia y la de Sebastián, sino que ahora tendrían la mente brillante de Don Ramiro a su lado.

\begin{quote}
    —Por supuesto —dijo Julio, sintiendo una nueva energía—. Síganme al taller.
\end{quote}

El grupo se dirigió al cobertizo. Don Ramiro entró primero, sus ojos escaneando el lugar con una curiosidad profesional. Para él, acostumbrado a los circuitos, las soldaduras y los esquemas complejos de la electrónica, el taller de Julio no era un desorden caótico, sino un laboratorio vivo.

Sus ojos, llenos de experiencia, se posaron de inmediato en los monitores parpadeantes con datos, en los prototipos intrincados de **Núcleos** a medio montar, en las herramientas de precisión, y finalmente, en el Núcleo principal que flotaba, brillando suavemente en su soporte. Don Ramiro se acercó a la esfera, observándola con una intensidad que nadie más podía igualar.

\begin{quote}
    —Vaya… —murmuró Don Ramiro, sin poder ocultar su asombro. Pasó una mano por las capas de grafito prensado, cobre enrollado y aluminio reforzado, sintiendo la textura, la robustez de la construcción. Para él, no era una fantasía, sino una obra de ingeniería. Reconoció de inmediato el trabajo minucioso y la complejidad del diseño.
\end{quote}

Julio comenzó a explicarle, con términos técnicos que los demás apenas entendían, sobre el cubo indestructible, la radiación de la supernova, el portal, Pi, y el funcionamiento del Núcleo como captador y almacenador de energía. Habló de su frustración por no poder extraer la energía de forma controlada para usos externos, y de los mensajes de "Ellos" sobre la conciencia y los Custodios del Tiempo.

Don Ramiro escuchaba con atención, sus ojos moviéndose entre Julio, la máquina, y el Núcleo. Su rostro, que normalmente mostraba una expresión relajada, ahora reflejaba una profunda concentración y una admiración creciente. No había rastro de incredulidad, sino el ojo de un colega, un ingeniero que calificaba y se maravillaba ante el trabajo de una persona a la que ahora, más que un cliente, le tenía un respeto y una admiración inmensos.

\begin{quote}
    —Esto… esto es… —Don Ramiro no encontró las palabras adecuadas. Tomó un respiro profundo y luego miró a Julio con una expresión solemne. —Julio, tu trabajo aquí no es solo avanzado, es… una obra de genio. Yo, con todos mis años de experiencia, veo el nivel de detalle y la intuición que has puesto en esto. La integración de los sistemas… el diseño del Núcleo… es increíble. No es una locura, es una hazaña.
\end{quote}

La validación de Don Ramiro fue un bálsamo para Julio. Saber que alguien con esa experiencia lo veía con respeto y no como un niño con fantasías, le dio una nueva inyección de confianza.

---

En la sala, con la nueva información sobre la estrella desvanecida y el apoyo de Don Ramiro, las familias se sentaron en un círculo improvisado. La conversación se volvió más profunda y las decisiones, más radicales.

\begin{quote}
    —Hemos estado pensando mucho —dijo Carlos, el padre de Sebastián—. Esta situación… lo cambia todo. No podemos seguir viviendo como si nada. Mi esposa y yo hemos decidido dejar nuestros empleos. Y Sebastián… Sebastián dejará el colegio.
\end{quote}

Elena asintió, su voz ahora más firme. —Es la única manera de dedicarnos por completo a esto. Necesitamos estar preparados. Sebastián, aunque dejará las clases regulares, no abandonará los deportes. Sabemos que es algo que lo distrae, que lo mantiene equilibrado. Pero sí se enfocará, al igual que nosotros, en prepararse para lo que se viene. Si es que siquiera podemos prepararnos, o si no hay nada que hacer… solo podemos intentarlo.

La decisión de los padres de Sebastián, tan drástica y valiente, impactó profundamente a Ana y Miguel. Era un salto al vacío, un abandono de la vida normal por una incertidumbre cósmica.

\begin{quote}
    —Ustedes esperan respuestas de Julio —continuó Carlos, mirando a los padres de Julio y luego a su hijo—, pero entienden que es tan solo un niño. Un niño que se topó con una verdad que, si bien está más allá de la comprensión de la mayoría, todos sabemos bien que es nuestra única esperanza.
\end{quote}

El silencio fue abrumador. La carga de la responsabilidad de Julio se hizo aún más palpable.

Fue entonces cuando Ana, con una mirada decidida a Miguel, tomó la palabra. —Carlos tiene razón. No podemos quedarnos de brazos cruzados. Miguel y yo… hemos tomado la misma decisión. Abandonaremos nuestros trabajos. Y Julio y Óscar, ustedes también dejarán el colegio.

Julio y Óscar se miraron, sorprendidos, pero también sintiendo una extraña liberación. El colegio, en este momento, parecía una distracción insignificante.

\begin{quote}
    —Sabemos que esto puede resultar alarmante para las unidades educativas —explicó Elena, quien como profesora de colegio, entendía las implicaciones—. Pero no es algo fuera de lo común que las familias decidan cambiar el estilo de enseñanza. Yo estoy calificada para dar clases de estilo \textit{homeschool}. Podemos seguir el plan de estudios, pero con la flexibilidad necesaria para dedicar tiempo a lo que realmente importa ahora. Para Julio, para Óscar, para Sebastián. Necesitamos que sus mentes sigan activas, pero enfocadas en esto.
\end{quote}

Con esa declaración, ambas familias sellaron su destino. Miguel y Carlos, dos padres que apenas se conocían hace unos meses, ahora se unían en una misión impensable. Ana y Elena, dos madres con miedos y esperanzas compartidas, se convertían en pilares de apoyo. Los chicos, Julio, Óscar y Sebastián, ya no eran solo amigos, sino compañeros en una aventura que definiría su existencia.

Y en este nuevo equipo, Don Ramiro (48 años), el hombre que había calificado el trabajo de Julio como una hazaña de genio, decidió que no se quedaría al margen. Con la admiración clara en sus ojos, le dijo a Julio: —Joven, yo seré tu mano derecha en esto. Lo que necesites en electrónica, en ingeniería, en lo que sea. Trabajaremos como un solo equipo.

La decisión de cambiar todo para prepararse a algo que no se podían imaginar, algo que estaba más allá de su comprensión, marcó el fin de una era y el inicio de una lucha desesperada. El reloj cósmico seguía su marcha, pero ahora, no estaban solos en la espera.

---

La cuenta regresiva seguía su implacable descenso: **371 días, 18 horas, 55 minutos, 30 segundos…**

\cleardoublepage % Para que el capítulo empiece en una página impar

% Vigésimo cuarto capítulo
\chapter*{24 – Selección y Sacrificio}
\addcontentsline{toc}{chapter}{24 – Selección y Sacrificio} % Añadir al índice manualmente

La mañana se arrastraba lentamente, pero el peso de la noche anterior seguía en el aire. La revelación de la estrella desvanecida y el escalofriante mensaje de "**Elige a quién salvar**" habían sumido a Julio en una profunda reflexión. Su mente, habituada a los cálculos y la lógica de la ciencia, ahora se enfrentaba a un dilema moral que lo carcomía desde adentro.

En la casa de Julio, la actividad era febril, aunque silenciosa y concentrada. No había tiempo para el luto, solo para la acción. Los adultos, ahora totalmente conscientes de la inminente catástrofe, habían dividido sus tareas con una eficiencia que el miedo y la determinación impulsaban.

Carlos, el padre de Sebastián, que en su trabajo manejaba maquinaria pesada, estaba junto a Miguel, el padre de Julio, revisando los planos de las casas. Conversaban sobre todas las ideas de modificaciones que podían hacer no solo al hogar de Julio, sino también al de Carlos. El objetivo ya no era solo aguantar la radiación, sino también crear un sistema que les permitiera ser autosustentables.

\begin{quote}
    —Mira, Miguel, si usamos esos soportes de acero que conseguí, podemos reforzar la estructura del sótano de Julio, y también la mía —decía Carlos, señalando unos croquis improvisados sobre la mesa de la cocina. Su voz, normalmente relajada, tenía un tono de urgencia profesional. —Y en cuanto a los generadores, mi compañía tiene unos modelos industriales que podríamos adaptar. También conozco proveedores de paneles solares de alta eficiencia. Podemos instalarlos en los tejados y crear un sistema de energía de respaldo robusto.
\end{quote}

Miguel asentía, su mente de ingeniero absorbiendo cada detalle. —Sí, la auto sustentabilidad es clave. No sabemos cuánto tiempo estaremos ahí abajo, ni qué quedará de la superficie. Necesitamos almacenar agua, reciclarla. Los sistemas de ventilación tienen que ser herméticos.

Carlos golpeó el plano con un dedo. —En mi casa puedo empezar con las excavaciones para un búnker más profundo. Tengo acceso a excavadoras y un equipo que puedo "tomar prestado" por unos días sin levantar sospechas. Haré todo lo que sea necesario, Miguel, y al mismo tiempo recopilaré toda la información y materiales para nuestro propio resguardo. Mis hombres no sabrán para qué es, pero haremos el trabajo rápido.

Mientras los padres planeaban las modificaciones a gran y mediana escala de sus respectivos hogares, incluyendo la instalación de generadores de energía y la compra de paneles solares, Sebastián y Óscar no se quedaban atrás. Estaban en la cocina con sus respectivas mamas, Elena y Ana, en una operación logística de supervivencia.

\begin{quote}
    —Óscar, ¿crees que podemos deshidratar estas manzanas sin que se oxiden? —preguntaba Ana, mientras Elena cortaba con habilidad frutas y verduras. Habían sacado todos los alimentos no perecederos de la despensa y estaban aprendiendo técnicas para conservar más.
\end{quote}

Óscar, con el deshidratador que Julio le había dado y una vaga idea de cómo funcionaba gracias a las explicaciones de su hermano, ya había montado el aparato. —Sí, mamá. Julio me explicó cómo controlar bien la temperatura y la humedad. Deberían aguantar. Ya estuve investigando formas de reciclaje. Podemos usar un sistema de compostaje para los desechos orgánicos y purificar el agua gris. Necesitamos minimizar todo el desperdicio.

Sebastián, aunque normalmente más despreocupado, trabajaba con una seriedad inusual, enlatando conservas y sellando paquetes. —Mi mamá dice que hasta el agua de la lluvia podemos usarla, si la filtramos bien. Y los envases, no podemos tirar nada. Cada botella, cada lata, puede servir para algo.

Las mamas, Elena (quien, como profesora, era experta en organización y planificación) y Ana (la pilar emocional y organizativo de su hogar), trabajaban con una determinación tranquila, enseñando a los chicos a preparar comidas enlatadas y a ingeniarse en nuevas formas de reciclaje para mantener el desperdicio al mínimo. Había una extraña belleza en la eficiencia con la que, en medio de la inminente catástrofe, se organizaban para sobrevivir.

---

En el taller, Julio y Don Ramiro estaban inmersos en el estudio del **Núcleo** y de **Pi**. La mente de Julio, brillante y veloz, encontraba en la experiencia de Don Ramiro una base sólida para sus ideas más audaces. Dos cerebritos, el joven y el experimentado, trabajando juntos, logrando avances que antes a Julio le costaba hacer solo. Don Ramiro revisaba los circuitos de Pi, calibrando sensores y analizando los extraños datos de la radiación.

Mientras trabajaban, Julio no pudo contener más la angustia que lo carcomía.

\begin{quote}
    —Don Ramiro —empezó, su voz baja, casi inaudible—, no entiendo. Estoy atrapado entre todas estas paradojas y mensajes que "Ellos" me han estado dejando. Son como migajas de pan para un rompecabezas gigante, pero ¿por qué no son más directos? ¿Por qué no me dicen qué hacer exactamente?
\end{quote}

Don Ramiro, sin apartar la vista de un microchip que estaba examinando, respondió con la calma que solo la experiencia de la edad puede dar. —Julio, a veces, uno necesita decir mucho y no tiene tiempo para hacerlo. O quizás, simplemente, no pueden decir más. Piensa en un manual técnico. No te explica la vida entera, solo te da los pasos cruciales. Y estas palabras cortas, estas "paradojas" como tú las llamas, muchas veces dicen más de lo que parecen. Hay que leer entre líneas.

Julio asintió, pero la frustración era palpable en su voz. —Pero hay una cosa que me molesta particularmente, algo que no sé cómo expresar, pero que me está volviendo loco. Necesito su opinión, Don Ramiro. "**Ellos**" me dijeron: "**Elige bien a quién salvas.**"

Don Ramiro detuvo su trabajo y miró a Julio, sus ojos denotando comprensión. —Eso es pesado, ¿verdad?

\begin{quote}
    —Es lo que más me pesa —continuó Julio, desahogándose—. No entiendo cómo yo podría tomar una decisión así. Entiendo que sea el único que tenga la mínima oportunidad de salvar, que yo sea la única esperanza… pero de ahí a que me sienta en el derecho de elegir quién vive y quién no, es algo que me carcome.
\end{quote}

Julio comenzó a verbalizar sus angustiantes pensamientos. —He pensado en familiares, en amigos, en conocidos de conocidos. Pensé en hacer una especie de lotería, que fuera al azar, para que nadie me culpe. Pero luego me topo con que ni siquiera sé a cuántos puedo salvar. ¿Diez? ¿Cincuenta? ¿Cien? Estoy completamente perdido.

Don Ramiro lo escuchó con paciencia, dejando que Julio vaciara su mente. Cuando el joven terminó, el hombre de 48 años, con su perceptible experiencia por la edad, le dijo:

\begin{quote}
    —Mira Julio, la gente no siempre se toma bien este tipo de cosas. La idea de seleccionar quiénes viven y quiénes no siempre fue un tabú en la historia de la humanidad. Siempre se habló de la igualdad, de la vida de todos. Pero ahora, no estamos en un tiempo normal. Algo sí te puedo asegurar: si el número de personas a salvar es limitado, tiene que ser **funcional**.
\end{quote}

Julio frunció el ceño. —¿Funcional?

Don Ramiro asintió. —Sí, funcional. Mira, Julio, si tienes un grupo de personas, sin importar la cantidad, y en este grupo, la mitad de los que salvas son, por así decirlo, personas descuidadas… digamos, con problemas de adicción, o quizás una persona con problemas mentales severos que necesite atención constante, o alguien que sea egoísta, o incluso, a manera de chiste, asesinos en serie. Si en un grupo así, tienes gente de ese estilo, no importa que salves al grupo, estarás destinado a fracasar. El grupo estará destinado a perecer desde adentro. Lo dice la ciencia, Julio. Un grupo sin cohesión, sin un propósito compartido y sin la capacidad de cuidarse y respetarse mutuamente, colapsará. ¿Estás de acuerdo?

Julio pensó en sus clases de biología y sociología. Recordó el **Número de Dunbar**, una teoría que decía que una persona solo puede mantener o establecer relaciones sociales estables y significativas con un máximo de aproximadamente 150 individuos. Más allá de ese número, el cerebro humano no puede mantener la información necesaria para almacenar más relaciones personales sólidas. Dunbar también argumentaba que una sociedad funcional, una que no dependiera de una economía compleja, de un orden establecido o de leyes escritas, tampoco era posible mantenerse con más de ese número. Era el límite de la verdadera "comunidad" en un sentido primario.

\begin{quote}
    —Sí —murmuró Julio, con una nueva comprensión—. Lo dice la ciencia. La gente necesita un mínimo de confianza, de cooperación. Si no… el grupo se autodestruye.

    —Exacto —dijo Don Ramiro—. Tu elección no es solo de "quién", sino de "cómo" ese grupo va a sobrevivir. Deberás elegir gente que, junta, forme una unidad autosuficiente, que puedan colaborar, que sus habilidades se complementen. Como esos **Custodios del Tiempo** de los que hablas, ¿no? Sus fortalezas impulsan las debilidades.
\end{quote}

Las palabras de Don Ramiro fueron un rayo de luz en la oscuridad del dilema de Julio. "**Elige bien a quien salvas**" no era un mandato de egoísmo, sino de pragmatismo para la supervivencia.

---

Luego de que terminó de hablar con Don Ramiro, Julio comenzó a repetir en voz alta: —**150**. Salvar a 150 personas es el límite. —Y una vez más, como perdido en sus pensamientos, repitió: —**150**.

Trabajaba al lado de Ramiro, soldando un circuito, con Fox acurrucado en sus piernas. En ese momento, la máquina se encendió nuevamente. Un punto de luz apareció en el centro de la pantalla de Pi. Don Ramiro se alarmó, pero Julio le puso la mano en la espalda, como diciéndole que se tranquilizara. Ahí apareció en la pantalla de Pi el número que Julio estaba repitiendo en voz alta: **150**.

Como si de una afirmación se tratase, como si "**Ellos**" hubieran escuchado lo que estaba pasando ahí.

\begin{quote}
    —¿Tú hiciste eso? —preguntó Don Ramiro, con los ojos bien abiertos.

    —No. Fueron ellos —respondió Julio, con una claridad que lo sorprendió incluso a él mismo.
\end{quote}

Era claro. Era lo que necesitaba. Tenía que buscar 150 personas que juntas lograran sobrevivir, basado en habilidades y responsabilidades. Un peso fue retirado de la cabeza de Julio casi de inmediato. Ya tenía un número y una forma de hacerlo. Y la culpa o la responsabilidad de la elección ya no era culpa, sino una guía que le habían dado.

---

La cuenta regresiva seguía su implacable descenso: **369 días, 16 horas, 20 minutos, 45 segundos…**

\cleardoublepage % Para que el capítulo empiece en una página impar

% Vigésimo quinto capítulo
\chapter*{25 – La Respuesta del Vecindario}
\addcontentsline{toc}{chapter}{25 – La Respuesta del Vecindario} % Añadir al índice manualmente

La luz del amanecer apenas se filtraba por las ventanas de la cocina de Julio, pero la casa ya no era un hervidero de actividad. Por el contrario, un silencio tenso la envolvía, roto solo por el suave zumbido de la máquina en el taller. Julio estaba solo en la cocina, un tazón de cereales intacto frente a él. Los planos de remodelación que sus padres y Carlos habían estado usando el día anterior, delineando el refuerzo del sótano y la posible excavación para el búnker, seguían extendidos sobre la gran mesa, abandonados a la espera de más discusión.

Pero la mente de Julio no estaba en el desayuno. Estaba en las palabras de Don Ramiro, martillando un número y una condición: **150. Funcional.**

\begin{quote}
    "Salvar a 150 personas es el límite," se repetía Julio en un susurro inaudible, sus pensamientos perdidos en la abstracción de lo que ello significaba. "Tiene que ser funcional... basado en habilidades y responsabilidades."
\end{quote}

La idea de que el grupo no solo debía ser limitado por recursos, sino también por una cohesión y complementariedad de talentos, era el nuevo pilar de su angustia y su esperanza. No era solo un número; era una comunidad viva que tenía que construir, una minúscula sociedad que debía poder sostenerse a sí misma, sin importar lo que el fin del mundo les arrojara. La culpa de la elección, que antes lo carcomía, ahora se había transformado en el peso de una responsabilidad pragmática. No se trataba de elegir favoritos, sino de identificar las piezas de un engranaje.

Mientras su mirada seguía perdida sobre los planos, Fox, con su habitual gracia, saltó silenciosamente sobre la mesa. El gato se paseó con parsimonia entre el tazón de cereales y los rotuladores, hasta que finalmente se recostó, con un suspiro felino, sobre los planos desplegados. Casualmente, su cuerpo cubrió el área donde estaba dibujada la maqueta de la casa de Julio.

Un instante después, Fox se levantó de golpe, como si algo lo hubiera perturbado, y se estiró, revelando el plano debajo. La mirada de Julio se posó en el dibujo de su propia casa, y luego, con una repentina sacudida de comprensión, siguió las líneas del plano que delineaban no solo su hogar, sino toda la manzana.

Un escalofrío le recorrió la espalda. Ahí estaba. La respuesta. Una respuesta que no había visto por estar mirando demasiado lejos, cuando la tenía tan cerca.

Con un entusiasmo que contrastaba con su anterior melancolía, Julio se incorporó, casi volcando su tazón de cereales. Se inclinó sobre los planos, y con un dedo, comenzó a trazar las líneas que delimitaban las propiedades vecinas. Contó con una emoción creciente la cantidad de casas que existían en su manzano: **¡50 casas!**

\begin{quote}
    "50 casas… si en promedio viven entre una a tres personas por casa, eso podría ser… ¡entre 50 y 150 personas!" Unas palabras que antes solo habían resonado en su cabeza, ahora parecían manifestarse ante sus ojos. La coincidencia, el "predestinado" del que había hablado la máquina, ahora se hacía tangible.

    —¡No puede ser! —murmuró Julio, con los ojos muy abiertos por la revelación. La cabeza le daba vueltas. Lo había pensado todo: familiares, amigos, conocidos, una lotería… pero nunca, nunca su propio vecindario. Y ahora, aquí, en el plano, lo veía claro.
\end{quote}

Se levantó abruptamente de la silla y, con una determinación renovada, caminó hacia la ventana. Miró a través del cristal hacia la calle, hacia su vecindario. La vieja manzana. Su hogar. Era un pequeño barrio, compuesto de casas modestas, algunas con jardines cuidados, otras con fachadas ligeramente desgastadas por los años. En el corazón de esa comunidad, se encontraba la casa de Julio, una estructura familiar que había sido su refugio durante toda su vida. Pero ahora, al mirarla con una perspectiva diferente, el vecindario le parecía algo más. Algo casi... predestinado.

Observó a los vecinos caminar por la calle. Vio a la señora del huerto, que siempre cultivaba tomates y zanahorias en su pequeño jardín, la enfermera que ayudaba a todos en emergencias, el guardia de seguridad que patrullaba el vecindario con su uniforme. Todos ellos, aparentemente comunes, pero ahora, para Julio, cada uno de esos rostros comenzó a adquirir un nuevo significado.

\begin{quote}
    —No puede ser... —murmuró Julio para sí mismo, su mente saltando de idea en idea, como si de repente las piezas del rompecabezas comenzaran a encajar con una velocidad vertiginosa.
\end{quote}

Se dirigió rápidamente hacia la puerta y salió al jardín. Corrió hasta la casa de su vecino, el primero en la calle. Golpeó la puerta, nervioso, casi sin aliento, pero con un nuevo respeto en su mirada.

Cuando la puerta se abrió, un hombre mayor, de unos 60 años, lo miró con curiosidad. Era Don Fernando, un hombre reservado pero siempre amable.

\begin{quote}
    —¿Julio? ¿Qué pasa, muchacho? —preguntó Don Fernando, notando la urgencia en los ojos del joven.
\end{quote}

Julio respiró hondo, tratando de ordenar sus pensamientos. Estaba dudando, pero las piezas del rompecabezas estaban allí, y la validación de Don Ramiro le daba el valor para seguir.

\begin{quote}
    —Disculpe, Don Fernando —comenzó Julio, su voz un poco más pausada y respetuosa de lo habitual, consciente de la edad del hombre—. ¿Le importaría si le hago un par de preguntas? Son importantes.
\end{quote}

Don Fernando lo miró extrañado por la formalidad, pero asintió con una sonrisa y lo invitó a pasar. —Claro, muchacho. Pasa, pasa.

Julio se sentó en la cocina mientras el hombre preparaba una taza de té. Finalmente, se armó de valor para ir al grano, con la franqueza que la situación demandaba.

\begin{quote}
    —Don Fernando, con todo respeto, ¿usted sabe a qué se dedica realmente? —preguntó Julio sin rodeos, pero con un tono que dejaba claro que no era una pregunta ociosa.
\end{quote}

El hombre frunció el ceño, claramente desconcertado por la inusual pregunta de su joven vecino. —Bueno, Julio, como sabes, soy electricista. He trabajado toda mi vida con eso. Ayudo con las reparaciones en la ciudad, todo lo que se refiere a conexiones y cables. ¿Por qué lo preguntas con tanto interés?

Julio se quedó en silencio por un momento, asimilando la respuesta. ¡Era perfecto! El hombre, un experto en reparaciones eléctricas, una habilidad esencial para un refugio autosustentable, estaba en la lista. Un punto más para la ecuación. Su corazón se aceleró con una nueva ola de esperanza.

\begin{quote}
    —Es... es un detalle curioso, Don Fernando —dijo Julio finalmente, poniéndose de pie para irse, sintiendo la necesidad de seguir buscando. —Muchas gracias por su tiempo. Disculpe la intromisión.
\end{quote}

Don Fernando, aunque no entendió del todo lo que había sucedido, notó la seriedad inusual en Julio y lo dejó ir con una sonrisa. Cuando Julio salió de la casa, su mente comenzó a dar vueltas más rápido. ¡Electricista! Una persona clave. Pero no solo eso, miraba al siguiente vecino, y luego al siguiente. Y en cada casa, se le ocurría algo que encajaba. En cada rostro que veía, podía ver una habilidad única que podría ser crucial para la supervivencia del grupo.

Se dirigió a la casa de la señora del huerto, Doña Marta. Ella estaba en su jardín, recogiendo algunas zanahorias, su espalda ligeramente encorvada por los años, pero sus manos firmes.

\begin{quote}
    —¿Doña Marta? —Julio la llamó desde la vereda, acercándose con respeto.
\end{quote}

La señora se volvió y sonrió al verlo. Su rostro arrugado se iluminó. —Ah, Julio, qué alegría verte. ¿Te gustaría llevarte algunas zanahorias frescas para tu familia? Son de mi cosecha.

Julio se acercó a ella y la observó de nuevo, sabiendo lo que venía a continuación, y una nueva pregunta para el rompecabezas.

\begin{quote}
    —Doña Marta, muchas gracias. Quería hacerle una pregunta, si no es mucha molestia. ¿Usted tiene algún tipo de formación, algún conocimiento especial en la agricultura o en la horticultura?
\end{quote}

La señora lo miró sorprendida, pero tras una pausa, una risita escapó de sus labios. —Bueno, muchacho, siempre he sido muy buena con las plantas. Es algo que aprendí de joven, de mis abuelos. Mis verduras siempre han sido de las mejores en el vecindario, ¿verdad? Pero si hablas de algo profesional... ¡pues no, no soy experta! Sólo soy una vieja aficionada con mucha experiencia.

Julio sonrió, pero dentro de él, las piezas continuaban encajando. Agricultura... Producción de alimentos... ¡Una habilidad crucial! Doña Marta no solo era alguien que podía ayudar a alimentar al grupo, sino que su huerto podía convertirse en la base de la supervivencia del refugio si se replicaba.

Julio continuó caminando, y con cada casa, con cada vecino, su convicción creció. El vecindario entero parecía estar constituido por personas con habilidades esenciales: desde **médicos** hasta **técnicos en ingeniería**, desde **cocineros** hasta **expertos en seguridad**. Como si la vida misma hubiera preparado a esas personas para algo más grande.

Cuando Julio regresó a su casa, se sentó frente a la máquina, aún atónito por lo que acababa de descubrir.

**150 personas.** Ese era el número mágico que "Ellos" le habían revelado a través de Don Ramiro. Y, de alguna manera, la vida lo había preparado.

\begin{quote}
    —Pi —dijo finalmente, mirando la pantalla—, tengo la respuesta.
\end{quote}

La máquina zumbó, su luz parpadeando, y la voz de PI resonó suavemente.

\begin{quote}
    —¿Has logrado encontrar la solución, Julio?
\end{quote}

Julio sonrió, aliviado, pero también con una sensación extraña de incredulidad.

\begin{quote}
    —Mi vecindario... Todo lo que necesito está aquí. Cada persona tiene lo que se necesita para sobrevivir y más. No es solo por casualidad. Es como si estuvieran aquí para algo más grande.
\end{quote}

Pi no respondió de inmediato, pero finalmente, la máquina emitió un tono afirmativo.

\begin{quote}
    —Bien hecho, Julio. Has encontrado la clave. Ahora es momento de construir la sociedad.
\end{quote}

El tiempo seguía corriendo, pero al menos ahora Julio sabía que no estaba solo. 150 personas, una comunidad que estaba lista para enfrentar lo que fuera necesario. Y a pesar del apremiante fin del mundo que se acercaba, una chispa de esperanza comenzó a arder dentro de él.

---

La cuenta regresiva seguía su implacable descenso: **368 días, 16 horas, 15 minutos, 20 segundos…**

\cleardoublepage % Para que el capítulo empiece en una página impar

% Vigésimo sexto capítulo
\chapter*{26 – El Proyecto del Arca}
\addcontentsline{toc}{chapter}{26 – El Proyecto del Arca} % Añadir al índice manualmente

La mañana se arrastraba, pero la mente de Julio ya no estaba atrapada en un dilema existencial. La revelación en la cocina sobre el número **150** y la inesperada "respuesta" de su propio vecindario, sumado a las sabias palabras de Don Ramiro, habían transformado su angustia en una determinación férrea. Ahora, con un propósito claro, la tensión en la casa de Julio no era de incertidumbre, sino de una concentración intensa.

En el taller, la atmósfera era de una seriedad palpable. Julio observaba el mapa de su vecindario, ahora no con marcas caóticas, sino con líneas y círculos que delineaban un plan preciso. El gran obstáculo de cómo construir un refugio para **150 personas** y, más importante, cómo convencerlos a todos, había encontrado su peculiar solución. Ya no se trataba de explicaciones a cada vecino, sino de una acción encubierta y monumental. No había tiempo para debates; sabía que un anuncio público solo generaría pánico e incredulidad.

Se sentó, pero esta vez, sin rastro de frustración. "¿Cómo voy a salvarlos? Ya lo sé", pensó. El peso de la responsabilidad seguía siendo inmenso, pero ahora tenía una guía, un mapa y una forma de hacerlo.

Fue entonces que Fox, que dormía en una esquina del taller, levantó la cabeza como si percibiera el cambio de energía en Julio. La máquina, Pi, empezó a emitir un zumbido suave, constante. Los símbolos extraños que Julio ya conocía, aquellos que se habían iluminado en momentos cruciales, comenzaron a brillar uno por uno en la pantalla.

De pronto, una voz resonó en su mente, serena y firme:

\begin{quote}
    —**Has cumplido con los requisitos. —Has tomado las decisiones correctas en el tiempo límite. —Es hora del siguiente paso.**
\end{quote}

La máquina vibró con una fuerza inusitada y, de su núcleo, empezó a expulsar cubos idénticos al primero, uno tras otro, como si no tuviera fin. El suelo del taller temblaba ligeramente con cada nuevo cubo. Julio los miraba caer, apilándose en montones con un sonido metálico, atónito, con una mezcla de terror ante la magnitud de la tarea y una esperanza abrumadora por la provisión.

Cuando finalmente la máquina se detuvo, el silencio se apoderó de nuevo del taller, y la voz volvió:

\begin{quote}
    —**Distribúyelos alrededor del perímetro del manzano. —Entiérralos exactamente a 20 metros de profundidad. —Cada cubo cumplirá su función cuando llegue el momento.**
\end{quote}

Un pequeño compartimiento de la máquina se abrió, dejando caer un paquete compacto. Julio lo abrió con manos temblorosas y encontró unos planos detallados. A medida que los desplegaba sobre su mesa de trabajo, su aliento se detuvo. Lo que mostraban no era un simple refugio...

Era un **Arca**.

Una estructura colosal que, usando los cubos que acababan de aparecer como cimientos energéticos, emergería desde el subsuelo en el instante preciso en que la radiación de la supernova golpeara. Esta Arca no solo aislaría y protegería a toda la manzana, sino que la transformaría en una fortaleza autosustentable. El núcleo principal, que Julio ya tenía y con el que había estado trabajando, debía instalarse en el centro de la manzana para estabilizar la transformación y el posterior funcionamiento del Arca, actuando como su corazón energético.

La voz finalizó con un susurro que parecía casi afectuoso, pero que resonó con el peso de un destino ineludible:

\begin{quote}
    —**No fue suerte, Julio. —Fuiste elegido... porque eras el único que podía.**
\end{quote}

Julio, aún temblando por la magnitud de la revelación, miró el plano del Arca, luego los cubos apilados, y finalmente, a través de la ventana, su vecindario dormido. Todo el peso de la misión se materializó en su pecho. Ahora entendía que "Ellos" habían estado guiándolo a través de las "migajas de pan", como Don Ramiro había sugerido, hacia este punto exacto. La selección no era suya, sino una confirmación de lo que ya estaba predispuesto.

No podía hacerlo solo. Pero ya no estaba solo. Tenía a su familia, a la de Sebastián, y ahora contaba con la experiencia invaluable de Don Ramiro y la fuerza práctica de Carlos.

A la mañana siguiente, no hubo necesidad de largas explicaciones. La familia de Julio ya estaba consciente del peligro y había aceptado la realidad. Se habían comprometido a ello al dejar sus trabajos y la escuela. La familia de Sebastián también estaba lista para actuar, su fe en Julio y en la urgencia de la situación era absoluta. Carlos y Don Ramiro, con sus habilidades complementarias, eran ahora parte integral del equipo.

La tensión en el aire era palpable, pero ahora mezclada con una determinación inquebrantable. Todos entendían la magnitud de la tarea, y cada uno asumió su rol con una seriedad que asombraría a cualquier observador externo.

Se organizaron rápidamente:
\begin{itemize}
    \item Julio y Óscar serían los encargados de dirigir las excavaciones más precisas para el posicionamiento de los cubos, siguiendo al milímetro los planos del Arca. Su trabajo, aunque agotador, sería guiado por los nuevos diagramas y por la intuición de Julio. Óscar, aunque no comprendía del todo la ciencia detrás de cada máquina o proceso, confiaba plenamente en su hermano y seguía sus instrucciones al pie de la letra, sabiendo que cada tarea era vital.
    \item Miguel y Carlos, con su experiencia en ingeniería y manejo de maquinaria pesada, serían los pilares de la operación subterránea. Carlos, en particular, se ofreció a usar sus contactos para "tomar prestados" los equipos de excavación necesarios, asegurando que su uso no levantara sospechas en la comunidad. Coordinarían las zonas de perforación en el perímetro del manzano, asegurando que cada foso tuviera la profundidad exacta de **20 metros**.
    \item Ana y Elena, con la ayuda de Sebastián, se encargarían de la logística en la superficie. Su misión sería cubrir y ocultar las obras a los vecinos curiosos, creando la ilusión de que eran reparaciones normales en la calle, o incluso proyectos de jardinería comunitaria. También coordinarían el suministro constante de agua, alimentos y herramientas para el equipo de excavación, y se asegurarían de que los recursos recolectados para la supervivencia (como las comidas enlatadas y frutas deshidratadas que habían estado preparando) estuvieran organizados y listos para el momento crucial.
\end{itemize}

El plan era simple en papel, pero monumental en su ejecución. Cavar 20 metros de profundidad para cada cubo, sin alertar a nadie, no era una tarea sencilla. Era necesario ocultarlo, trabajar con extrema rapidez y, sobre todo, no llamar la atención.

Durante semanas, trabajaron incansablemente, principalmente de noche, bajo el manto protector de la oscuridad. Primero, empezaron en el patio trasero de la casa de Julio, perfeccionando la técnica de excavación con la maquinaria ligera que Carlos había conseguido. Luego, expandieron las operaciones hacia los callejones y los solares vacíos del manzano, moviendo la tierra con precisión y cautela.

Fox siempre vigilaba, atento, como si entendiera perfectamente lo que estaba en juego, a menudo acostado cerca de Julio mientras este revisaba los cálculos o daba instrucciones, su presencia un silencioso consuelo.

La fatiga era brutal. Cada golpe de pala, cada movimiento de la maquinaria, era una carrera contra el reloj cósmico. En más de una ocasión, Julio estuvo a punto de desmayarse de agotamiento, pero la imagen de sus vecinos —Doña Marta del huerto, Don Fernando el electricista, la enfermera, el guardia de la noche, los niños corriendo en bicicleta— le daba la fuerza necesaria. Todo lo que amaba, todo lo que "Ellos" le habían encargado proteger, estaba en riesgo. No podía fallarles.

A medida que los días se transformaban en semanas, Julio empezó a notar que el ambiente, se volvía cada vez más extraño: el clima cambiaba de forma errática, las noches eran demasiado frías o demasiado cálidas, los animales se comportaban de manera inusual, y había un zumbido en el aire, apenas perceptible, que parecía retumbar en los huesos, una señal inconfundible de que el tiempo se acortaba.

La cuenta regresiva avanzaba, inexorable.

Finalmente, una noche, mientras terminaban de enterrar el último cubo en el extremo norte del vecindario, Julio se detuvo, exhausto, cubierto de tierra y sudor. Miró el cielo... y pensó en todo lo que habían hecho.

Habían logrado lo imposible.

Ahora, solo faltaba instalar el núcleo central y esperar la transformación.

El viento sopló fuerte, como un susurro de advertencia.

---

La cuenta regresiva seguía su implacable descenso: **252 días, 06 horas, 22 minutos, 01 segundos…**

\cleardoublepage % Para que el capítulo empiece en una página impar

% Vigésimo séptimo capítulo
\chapter*{27 – La Calma antes de la Tormenta}
\addcontentsline{toc}{chapter}{27 – La Calma antes de la Tormenta} % Añadir al índice manualmente

El aire de la noche era fresco, pero el calor del esfuerzo recién terminado aún los envolvía. El último cubo había sido enterrado con una precisión casi obsesiva, la tierra compactada hasta borrar cualquier señal de la monumental excavación. Julio, Miguel, Carlos, Óscar y Don Ramiro, sus cuerpos doloridos y cubiertos de sudor y tierra, se miraron bajo la tenue luz que comenzaba a asomar en el horizonte. Un silencio profundo, cargado de agotamiento, se cernía sobre ellos, pero bajo esa quietud, vibraba un triunfo silencioso. Habían logrado lo que parecía imposible. El **Arca**, su esperanza secreta y colosal, ahora dormía profundamente bajo su vecindario, aguardando el momento de su despertar.

De regreso en la casa de Julio, la atmósfera había cambiado radicalmente. La tensión constante que había reinado durante meses se había disipado, reemplazada por una ligereza que casi les parecía irreal. Ana y Elena, que con admirable estoicismo habían mantenido la fachada de la normalidad y la intrincada logística en la superficie, los recibieron con abrazos sinceros y palabras de admiración que resonaban en el cansado equipo. Habían terminado la parte más ardua. Ahora, les quedaban poco más de ocho meses, **252 días**, antes de que el mundo, tal como lo conocían, se transformara para siempre. Era un respiro, un tiempo para la transición, para sanar sus cuerpos y mentes.

Para conmemorar este logro titánico, la decisión unánime fue una parrillada. No una fiesta ostentosa que pudiera despertar la curiosidad de los vecinos y comprometer su secreto, sino una reunión íntima y profundamente significativa en el patio trasero de Julio. Un respiro, una pequeña burbuja de normalidad en medio de la inminente anomalía cósmica. La risa, un sonido que había sido tan escaso y fugaz en los últimos meses de arduo trabajo y ansiedad silenciosa, comenzó a brotar de forma natural, llenando el espacio con una alegría genuina.

Miguel, con la maestría de un veterano asador, encendió el carbón. Las brasas comenzaron a crepitar, emitiendo un aroma ahumado y reconfortante. Carlos, por su parte, preparaba las carnes con destreza, charlando animadamente con Miguel sobre los detalles técnicos de la excavación. Ya no necesitaban hablar en códigos; las palabras fluían libres, llenas de la satisfacción de un trabajo bien hecho.

\begin{quote}
    —¡Increíble, Miguel! La precisión con la que Óscar y Julio manejaron el posicionamiento de esos cubos… ni un ingeniero profesional de primer nivel lo habría hecho mejor —comentó Carlos, volteando un chorizo con pinzas. Su rostro, que durante tanto tiempo había estado surcado por la preocupación y el cansancio, ahora mostraba una sonrisa amplia y genuina. La carne chisporroteaba sobre la parrilla, prometiendo un festín.

    Miguel asintió, una sonrisa de orgullo indomable en su rostro mientras observaba a sus hijos. —Sí, la verdad es que estos muchachos nos han sorprendido a todos. Su dedicación, su ingenio... Y sin ti, Carlos, con tu experiencia en maquinaria y tus contactos para los equipos, esto habría sido, literalmente, imposible. No sé cómo agradecerte lo suficiente por todo lo que has arriesgado.

    Carlos le dio una palmada amistosa en el hombro. —No hay nada que agradecer, Miguel. Estamos en esto juntos. Esto es por nuestras familias, por nuestros hijos. Y mira, ¿qué más da el resto del mundo y lo que deje de ser? Lo que importa es que nuestros chicos, nuestras familias, tienen una oportunidad real de sobrevivir. Es todo lo que podemos pedirle al destino.
\end{quote}

Mientras los padres se entregaban a la satisfacción primal de la carne asada, regada con cervezas frías que sabían a triunfo, las madres, Ana y Elena, se afanaban en preparar las ensaladas frescas y las bebidas no alcohólicas para los más jóvenes. La mesa del jardín se fue llenando de platos humeantes, cubiertos relucientes y, sobre todo, el sonido contagioso de las risas. La cena era un banquete no solo para el cuerpo, sino para el espíritu.

Julio, reclinado en su silla, observaba el cielo estrellado, una bóveda inmensa y silenciosa que, por primera vez en meses, le ofrecía una quietud reconfortante. El peso que había cargado, el de buscar a **150 personas** y el de cómo construir el refugio, se había aliviado. Todo estaba hecho. O al menos, lo que estaba en sus manos hacer. La semilla de la esperanza había sido plantada bajo sus pies.

Óscar, con una sonrisa pícara que revelaba su alivio, interrumpió los pensamientos de Julio.
\begin{quote}
    —Bueno, Julio, ya que tenemos más de ocho meses de 'vacaciones' forzadas antes de que el mundo se ponga interesante, ¿qué crees que deberíamos hacer con todo este tiempo que nos queda? Ya tenemos la comida esencial, el agua, la supervivencia asegurada… o eso esperamos, ¿verdad?
\end{quote}

Sebastián se unió a la broma con un entusiasmo contagioso, sus ojos brillando con una chispa de travesura y alivio.
\begin{quote}
    —¡Sí! ¡Podríamos hacer una lista de cosas "menos primordiales" para llevar al Arca! Ya tenemos lo esencial cubierto por nuestros padres, ¡ahora viene la parte divertida de la planificación!
\end{quote}

Julio, sintiendo una ligereza en el pecho que lo sorprendió, se dejó arrastrar por la idea. La noción de planificar frivolidades después de la inmensidad de lo que habían logrado, era extrañamente liberadora.
\begin{quote}
    —¡Claro! Esto es como planificar una fiesta de pijama gigante, pero de muy, muy larga duración, ¿no creen? Necesitamos asegurarnos de que el encierro sea lo más… habitable posible. A ver, ¿qué es lo primero que necesitamos para que el tiempo pase volando?

    —¡Juegos de video! —exclamó Óscar de inmediato, su voz un eco de la emoción juvenil—. ¡Y una televisión lo suficientemente grande como para que no tengamos que pelearnos por la vista! ¡Y, por supuesto, mandos extra para cuando vengan las visitas!

    —¡Juegos de mesa! —añadió Sebastián con la misma energía—. ¡Para cuando se vaya la luz, o cuando ya estemos hartos de ver pantallas!

    —Películas y series —sugirió Julio, dejando que su mente divagara en cómo llenar el tiempo. La idea de un maratón de películas le trajo una sonrisa. —Y música. Mucha música. Creo que la música será esencial para mantener la moral.

    —¡Y una biblioteca, Julio! —terció Óscar, entusiasmado. —No solo libros para leer, sino de consulta. Química, biología, ingeniería, historia... Todo lo que podamos necesitar si, bueno, si la humanidad tiene que empezar de nuevo ahí abajo. ¡Para que no se pierda el conocimiento!
\end{quote}

Los tres chicos se entregaron a la tarea de crear una lista imaginaria con una seriedad cómica, sus voces mezclándose con la risa y las conversaciones más profundas de los adultos. Había algo profundamente catártico en planificar frivolidades cuando la vida les había exigido tanto. Ya no eran solo los cerebritos obsesionados con el fin del mundo o los atletas preocupados por su futuro; eran jóvenes a los que se les había devuelto, aunque temporalmente, la inocencia de un futuro, por confinado que fuera.

Don Ramiro, sentado junto a Miguel y Carlos, escuchaba la charla de los muchachos, una sonrisa amable y contemplativa en su rostro.
\begin{quote}
    —Parece que ya tienen el plan para el encierro —comentó con su voz pausada y sabia—. Es una buena señal, ¿saben? Significa que confían plenamente en el trabajo que han hecho.

    —Plenamente, Don Ramiro —respondió Miguel, dándole una palmada de gratitud en la espalda—. No sé qué habríamos hecho sin usted. Su sabiduría, su experiencia, su calma... Usted fue el pilar de Julio en sus momentos más difíciles.

    —Y para todos nosotros —añadió Ana, acercándose a la mesa con más bebidas, su rostro suave y aliviado—. Ha sido un pilar inestimable, Don Ramiro. Su presencia nos dio la fuerza para seguir.
\end{quote}

La noche continuó, la parrillada un ancla en la realidad. La conversación fluía fácil, el sonido de las risas y las copas chocando resonaba bajo el cielo nocturno. Por primera vez en mucho tiempo, no hablaban de radiación, de estrellas desvanecidas, ni de desastres inminentes. En cambio, conversaban sobre la esperanza silenciosa que habían sembrado bajo sus pies, sobre el futuro que, a pesar de todo, se atrevían a imaginar. Se sentía bien, un alivio inmenso, merecido. Habían hecho todo lo que estaba en sus manos, lo habían dado todo. Ahora, solo quedaba esperar. Y en esa espera, en el seno de esa comunidad íntima y resilientemente unida, encontraron una extraña y profunda paz. El viento, que antes había soplado como un susurro de advertencia, ahora parecía una brisa que los arrullaba, un bienvenido preludio a la calma.

---

La cuenta regresiva seguía su implacable descenso: **252 días, 04 horas, 28 minutos, 50 segundos…**

\cleardoublepage % Para que el capítulo empiece en una página impar

% Vigésimo octavo capítulo
\chapter*{28 – 7 Días}
\addcontentsline{toc}{chapter}{28 – 7 Días} % Añadir al índice manualmente

Los meses que siguieron a la finalización del **Proyecto del Arca** fueron una especie de limbo surrealista, un período de calma tensa pero, sorprendentemente, plena de satisfacciones. Con más de ocho meses por delante, las dos familias, junto a Don Ramiro, se dedicaron con la misma intensidad que pusieron en la excavación a asegurarse de que su refugio subterráneo fuera no solo un búnker seguro, sino un hogar habitable y, en la medida de lo posible, confortable.

Revisaron cada detalle, una y otra vez. Los sistemas de purificación de agua fueron probados hasta el límite, los sistemas de ventilación se ajustaron para una eficiencia máxima, y el almacenamiento de alimentos se organizó con una meticulosidad casi obsesiva. Tenían suficiente comida enlatada, deshidratada y empacada al vacío para años, raciones calculadas para cada una de las **150 personas** que esperaban albergar, aunque los vecinos aún no lo supieran. La seguridad era primordial, y se aseguraron de que las entradas y salidas de emergencia fueran discretas y funcionales.

Pero más allá de lo esencial, se dedicaron a la "operación distracción", como la llamaba Óscar. Los padres, con las tarjetas de crédito humeantes y sin la menor preocupación por deudas o hipotecas que pronto dejarían de existir, se lanzaron a una frenética misión de compras. Tiendas de electrónicos vieron sus estanterías vaciarse de consolas de última generación, de videojuegos para todos los gustos y edades, y de televisiones de pantalla gigante. Librerías enteras fueron desvalijadas de clásicos de la literatura, enciclopedias, libros de texto sobre ciencia, historia, agricultura, medicina… cualquier conocimiento que pudiera ser útil o simplemente entretener. La idea era llenar el **Arca** no solo con la supervivencia, sino con la vida misma, con todo lo que les permitiera mantener la cordura y el espíritu humano.

En una de las incursiones más memorables, Carlos, con una sonrisa de complicidad, llevó a Miguel al supermercado más grande. Mientras Julio y Óscar se dedicaban a buscar los juegos de mesa más complejos y los rompecabezas de mil piezas, los padres se adentraron en el pasillo de licores. Con una determinación digna de una misión secreta, comenzaron a llenar carritos enteros con botellas de vino tinto y blanco, whisky de malta, bourbon añejo, y licores exóticos. La burla de los chicos se convirtió en risas cuando vieron que el búnker que Carlos había cavado, destinado a ser un almacén extra, acabó transformado en una verdadera cava subterránea.

\begin{quote}
    —¡Con esto tenemos alcohol suficiente para cien años, muchachos! —exclamó Carlos con una carcajada, levantando una botella de un costoso whisky. Miguel asintió con la cabeza, una sonrisa de satisfacción en su rostro. Era un lujo, sí, pero un lujo bien ganado, una forma de asegurar que la vida en el encierro no fuera solo funcional, sino también disfrutable.
\end{quote}

No solo pensaron en los vicios. La planificación para la autosustentabilidad era clave. Habían conseguido semillas de todo tipo de hortalizas, organizadas meticulosamente por tipo y fecha de siembra. La visión de Doña Marta y su huerto había inspirado a Julio a diseñar un sistema hidropónico rudimentario, pero efectivo, para el cultivo dentro del **Arca**. Y en un golpe de suerte que solo podía atribuirse al destino, o a "**Ellos**", lograron adquirir un pequeño grupo de gallinas ponedoras. Prepararon un espacio dentro del **Arca** para ellas, con la promesa de que, si bien no tendrían carne de gallina todos los días, la producción de huevos y la carne ocasional serían un gusto sano y nutritivo para todas las familias.

La confianza en el **Arca** era plena. Habían puesto todo de sí mismos en este proyecto y confiaban ciegamente en que funcionaría. Las hipotecas, las líneas de crédito, las deudas… todo eso dejó de ser una preocupación. Al fin y al cabo, en pocos días, ya no le deberían nada al banco. La sensación de liberarse de esas ataduras mundanas era casi tan satisfactoria como la de haber completado el **Arca**. Sabían que eran los únicos conscientes de la inminente catástrofe, pero se habían preparado siempre pensando en el barrio entero, en esas **150 personas** que, por una extraña coincidencia de habilidades y destino, parecían destinadas a sobrevivir juntos.

Así pasaron los meses, los días volando en un torbellino de preparativos finales y momentos de preciada calma. Hasta que, una tarde, mientras estaban todos reunidos en el patio, disfrutando de otra parrillada y riendo a carcajadas con los planes de los chicos para sus "suministros de ocio", Óscar, con una mueca de repentina lucidez, pronunció unas palabras que los congelaron a todos.

\begin{quote}
    —Y… ¿y si no están en casa cuando suceda esto? —dijo Óscar, la sonrisa desvaneciéndose de su rostro mientras la pregunta resonaba en el aire.
\end{quote}

El silencio se instaló de golpe. Las risas se extinguieron. Los padres, que segundos antes brindaban, bajaron sus vasos lentamente. Todos giraron a ver a Óscar, sus rostros reflejando la misma cruda realización. Era una pregunta tan simple, tan obvia, y sin embargo, tan fundamental.

\begin{quote}
    —¡Oh, por Dios! —exclamó Julio, sintiendo un escalofrío que le recorrió la espalda hasta la punta de los pies. Su mente, tan dada a los cálculos complejos, había pasado por alto esta variable humana tan básica. —¡Tienes razón, Óscar! ¡No podemos estar seguros de que todos estén en casa! ¡Mi Dios, qué hacemos!
\end{quote}

Fue como un balde de agua fría, un jalón brusco de regreso a la dura realidad. La complacencia se desvaneció en un instante. Después de meses de trabajo incansable, de sacrificios personales, se toparon con un problema que parecía imposible de resolver sin revelar su secreto. ¿Cómo asegurar que todos los vecinos del manzano estuvieran dentro del perímetro de seguridad del **Arca** en el momento exacto? La idea de tener que idear un plan de contingencia para esto, después de haber terminado todo lo demás, era agotadora.

Julio se dejó caer en una silla, el cansancio regresando a sus hombros con una fuerza abrumadora. Un suspiro de resignación escapó de sus labios.
\begin{quote}
    —Qué flojera… —murmuró, y a pesar de la gravedad de la situación, una pequeña y resignada sonrisa se dibujó en su rostro, contagiando una risita nerviosa entre los demás.
\end{quote}

Justo en ese momento, como si hubiera estado esperando su señal, la voz de Pi resonó en la mente de Julio, con la calma y la seguridad de siempre, una voz que ya se había convertido en un compañero silencioso y omnisciente:

\begin{quote}
    —Julio, ya tengo planificado un plan de contingencia para esta eventualidad. Mientras ustedes descansaban, preparé un troyano. Está instalado en la base de datos del centro meteorológico y de alertas tempranas de terremotos. He programado una alerta de un sismo de categoría **8 en la escala de Richter** que amenazará la región y comenzará a sonar **cinco horas antes del evento real**. Esta alerta, diseñada para maximizar el impacto, saldrá en todos los noticieros, en los teléfonos inteligentes y en las radios de emergencia. La sugerencia será que la gente se encierre en sótanos y bajo techo para evitar fatalidades por caída de estructuras o escombros. Esto nos asegurará que todos estén en casa y busquen refugio.
\end{quote}

Todos quedaron boquiabiertos, sorprendidos por la audacia y la previsión de la máquina. Julio también lo estaba, pero hizo un esfuerzo por disimularlo, adoptando una pose de creador imperturbable que merecía el reconocimiento de **Pi**. Con una leve inclinación de cabeza, como si fuera lo más natural del mundo, respondió:

\begin{quote}
    —Gracias, Pi. No esperaba menos. —Y continuó con lo que estaban haciendo, como si esa increíble revelación fuera solo una nota al pie. Los adultos se miraron, asombrados, pero a esas alturas, el shock se había transformado en una especie de asombro cómico. Una risa colectiva escapó de ellos. Miguel, con una expresión de orgullo desbordante, sacudió suavemente la cabeza de su hijo. La genialidad de Julio, combinada con la inexplicable capacidad de **Pi**, seguía siendo una fuente inagotable de asombro y, ahora, de un alivio inmenso.
\end{quote}

---

La cuenta regresiva seguía su implacable descenso: **7 días, 02 horas, 15 minutos, 40 segundos…**

\cleardoublepage % Para que el capítulo empiece en una página impar

% Vigésimo noveno capítulo
\chapter*{29 – Últimas Horas}
\addcontentsline{toc}{chapter}{29 – Últimas Horas} % Añadir al índice manualmente

La cuenta regresiva se había reducido hasta un punto casi insoportable: **0 días, 9 horas, 23 minutos, 55 segundos**.

La mañana había llegado pesada, con una luz gris y difusa que se filtraba por las ventanas, como si el mundo mismo supiera que algo estaba a punto de romperse de forma irrevocable. Julio, Óscar, sus padres Miguel y Ana, y Fox, desayunaban en silencio en la cocina de la casa, cada uno encerrado en sus propios pensamientos. El crepitar lejano de las brasas de la parrilla de la noche anterior, ahora fría, parecía un eco de los meses de trabajo febril.

Ya todo estaba preparado. Los cubos de energía estaban enterrados con precisión milimétrica alrededor del perímetro del manzano, sus profundidades confirmadas. El **núcleo principal**, el corazón de la futura **Arca**, estaba encendido y estable en el sótano reforzado de Julio, latiendo con una energía imperceptible pero constante. Las provisiones, desde la comida enlatada hasta los innumerables juegos de mesa y consolas, estaban apiladas hasta el techo, asegurando que la vida en el encierro sería, dentro de lo posible, confortable. Incluso el "tesoro" de vinos y whiskies de Carlos y Miguel estaba cuidadosamente resguardado en su improvisada cava subterránea, junto a las semillas de hortalizas y el pequeño grupo de gallinas ponedoras que prometían huevos frescos y carne esporádica.

La **falsa alerta de sismo de categoría 8**, ideada por **Pi**, estaba programada para activarse con una precisión asombrosa: exactamente **seis horas antes del momento final**. Esa era la última pieza de su elaborado plan, la estrategia para asegurar que todos los vecinos del manzano estuvieran resguardados en sus hogares.

Julio levantó su muñeca y consultó su reloj digital. Cada latido de su corazón parecía sincronizarse con los segundos que se desvanecían uno tras otro en la pequeña pantalla. La calma que había logrado sentir en los últimos días, nacida del agotamiento y la certeza de haber hecho todo lo posible, ahora se mezclaba con una punzada de nerviosismo.

\begin{quote}
    —Hoy —dijo Miguel, su voz resonando con una solemnidad inquebrantable en el silencio de la cocina— es el último día que veremos el mundo tal como lo conocemos. Asegúrense de que las mochilas de emergencia estén a mano y que cada uno tenga sus objetos personales más importantes listos.
\end{quote}

Ana bajó la mirada hacia su tazón, conteniendo las lágrimas que amenazaban con desbordarse. Su mano buscó la de Miguel, un gesto silencioso de apoyo mutuo.

Óscar forzaba una sonrisa mientras jugaba distraídamente con Fox, acariciando su suave pelaje. El gato, que normalmente era la encarnación de la calma silenciosa, parecía más inquieto que de costumbre, maullando de vez en cuando y frotándose contra las piernas de Julio y Óscar. Quizás los animales sentían la inminencia, la extraña vibración en el aire de esos últimos días.

Julio se levantó de la mesa. No podía quedarse sentado, inactivo. La responsabilidad era inmensa. Había decidido pasar esas últimas horas revisando, una y otra vez, todos los sistemas, cada protocolo. Quería asegurarse de que absolutamente nada fallaría.

Se dirigió al taller. Activó el dron que había fabricado, una maravilla tecnológica propia, y lo envió a verificar los puntos de enterramiento de los cubos alrededor del perímetro del manzano. En su pantalla, cada señal respondía con una luz verde confirmando la integridad de los dispositivos. La verificación fue perfecta. El núcleo en el sótano mostraba una estabilidad impecable, su zumbido apenas audible. La máquina Pi estaba activa, en modo de espera, lista para el momento crucial. Pi misma había revisado toda la programación de la alerta sísmica y le confirmó a Julio mentalmente que el mensaje se emitiría sin problemas, alcanzando cada dispositivo y cada hogar.

Faltaban **6 horas, 02 minutos y 38 segundos**.

Julio se sentó frente a la consola de **Pi**, extendiendo una mano sobre su superficie fría. Fox saltó a su regazo, acurrucándose y ronroneando. El joven cerró los ojos, respirando profundamente.

Una calma pesada cayó sobre él, una aceptación. No podía cambiar lo que vendría. Todo lo que había estado en su poder, lo había hecho. Solo podía esperar que todo saliera como estaba planeado. La certeza de que el **Arca** era su última y única esperanza llenó su mente.

Afuera, más allá de los muros de su casa, la vida en el vecindario parecía discurrir con una normalidad desesperante. Algunos vecinos cortaban el césped, el sonido de las cortadoras resonando monótonamente en el aire. Otros paseaban a sus perros, charlando con otros transeúntes. Algunos niños, despreocupados, jugaban en la calle, sus risas inocentes perforando el silencio de Julio.

Julio los miraba desde la ventana de su taller, una punzada en el pecho. Ignoraban por completo que en apenas unas horas, todo cambiaría. Que el cielo que les ofrecía esa mañana gris se convertiría en una furia, y que sus vidas normales, sus rutinas, sus pequeñas preocupaciones diarias, se extinguirían en un instante.

La cuenta regresiva en su reloj parpadeó.

**6 horas, 00 minutos, 00 segundos.**

De pronto, un sonido estridente rompió la quietud de la mañana. Los celulares comenzaron a vibrar con una furia inusitada. Las radios, que hasta ese momento emitían música suave, se interrumpieron abruptamente con una voz robótica y alarmante. Los televisores de los vecinos, que podían verse a través de las ventanas, emitieron la misma alerta estridente, un pitido agudo seguido de un mensaje que se repetía sin cesar:

\begin{quote}
    "**¡Alerta Nacional! Sismo de categoría 8. Refúgiense en sus casas inmediatamente. Cierre puertas y ventanas. Manténgase resguardado. Posible impacto en 6 horas.**"
\end{quote}

Los vecinos se miraron, confundidos. Algunos se reían nerviosos, murmurando sobre un error, una broma de mal gusto. Otros, más precavidos, o simplemente obedientes a la autoridad, comenzaron a reaccionar. Rápidamente, el movimiento en las calles se detuvo. Puertas se cerraban con golpes secos. Ventanas se aseguraban. Se veían figuras correr hacia sus hogares, siguiendo la instrucción de "refugiarse bajo techo".

Óscar sonrió, viendo cómo el plan de **Pi** funcionaba con una eficiencia aterradora. El simulacro de sismo era la coartada perfecta, una forma de meter a todos en casa sin pánico ni preguntas innecesarias.

Julio soltó un suspiro de alivio que no sabía que contenía. Era ahora o nunca. El momento decisivo se acercaba.

Se dirigió hacia su familia, que ya comenzaba a prepararse con una disciplina ensayada, cada uno cumpliendo su rol asignado. Miguel y Carlos se aseguraron de que las últimas herramientas estuvieran guardadas. Ana y Elena revisaron los kits de emergencia. Óscar y Sebastián cargaron las mochilas personales.

Juntos, bajaron al sótano, donde el núcleo del **Arca** esperaba, vibrando levemente en la oscuridad, como un corazón latiendo en anticipación. La presencia de Don Ramiro, con su calma serena, era un ancla en la creciente tensión.

La cuenta regresiva en el dispositivo de Julio seguía implacable:

**0 días, 5 horas, 57 minutos, 10 segundos.**

La tierra temblaría pronto.

Y ellos, la familia, Don Ramiro, y con suerte, las **150 personas** de su vecindario, estarían listos.

\cleardoublepage % Para que el capítulo empiece en una página impar

% Trigésimo capítulo
\chapter*{30 – El Comienzo del Fin}
\addcontentsline{toc}{chapter}{30 – El Comienzo del Fin} % Añadir al índice manualmente

El implacable goteo del tiempo se había reducido a meras migajas. En el panel de control de Julio, la cuenta regresiva parpadeaba con una frialdad desapasionada: **0 días, 9 horas, 23 minutos, 55 segundos**.

La mañana había llegado pesada, con una luz gris y difusa que se filtraba por las ventanas, como si el mundo mismo supiera que algo estaba a punto de romperse de forma irrevocable. Julio, Óscar, sus padres Miguel y Ana, y Fox, desayunaban en silencio en la cocina de la casa, cada uno encerrado en sus propios pensamientos. El crepitar lejano de las brasas de la parrilla de la noche anterior, ahora fría, parecía un eco de los meses de trabajo febril y la única nota de una normalidad ya lejana.

Ya todo estaba preparado. Los **cubos de energía** estaban enterrados con precisión milimétrica alrededor del perímetro del manzano, sus profundidades confirmadas. El **núcleo principal**, el corazón de la futura **Arca**, estaba encendido y estable en el sótano reforzado de Julio, latiendo con una energía imperceptible pero constante. Las provisiones, desde la comida enlatada hasta los innumerables juegos de mesa y consolas que habían agotado las tarjetas de crédito de los padres, estaban apiladas hasta el techo, asegurando que la vida en el encierro sería, dentro de lo posible, confortable. Incluso el "tesoro" de vinos y whiskies de Carlos y Miguel, suficiente para "cien años" de celebraciones o consuelo, estaba cuidadosamente resguardado en su improvisada cava subterránea, junto a las semillas de hortalizas de todo tipo y el pequeño grupo de gallinas ponedoras que prometían huevos frescos y carne esporádica. Todo, hasta el más mínimo detalle, había sido contemplado para asegurar la supervivencia y el bienestar de los **150 vecinos de la manzana**, aunque ellos aún vivieran en feliz ignorancia.

La **falsa alerta de sismo de categoría 8**, ideada por **Pi**, estaba programada para activarse con una precisión asombrosa: exactamente **seis horas antes del momento final**. Esa era la última pieza de su elaborado plan, la estrategia para asegurar que todos los vecinos del manzano estuvieran resguardados en sus hogares. Las hipotecas y las deudas con el banco, que antes habían sido preocupaciones aplastantes, ahora se sentían tan lejanas como un sueño; en unas horas, el concepto mismo de "banco" podría dejar de existir.

Julio levantó su muñeca y consultó su reloj digital. Cada latido de su corazón parecía sincronizarse con los segundos que se desvanecían uno tras otro en la pequeña pantalla. La calma que había logrado sentir en los últimos días, nacida del agotamiento y la certeza de haber hecho todo lo posible, ahora se mezclaba con una punzada de nerviosismo que se le aferraba al estómago.

\begin{quote}
    —Hoy —dijo Miguel, su voz resonando con una solemnidad inquebrantable en el silencio de la cocina— es el último día que veremos el mundo tal como lo conocemos. Asegúrense de que las mochilas de emergencia estén a mano y que cada uno tenga sus objetos personales más importantes listos. Si el plan funciona, no necesitarán nada más.
\end{quote}

Ana bajó la mirada hacia su tazón, conteniendo las lágrimas que amenazaban con desbordarse. Su mano buscó la de Miguel, un gesto silencioso de apoyo mutuo en el umbral de lo desconocido. Elena y Carlos se sentaron cerca, sus manos entrelazadas, una imagen de quietud y fortaleza compartida.

Óscar forzaba una sonrisa mientras jugaba distraídamente con Fox, acariciando su suave pelaje. El gato, que normalmente era la encarnación de la calma silenciosa y la indiferencia felina, parecía más inquieto que de costumbre, maullando de vez en cuando y frotándose con insistencia contra las piernas de Julio y Óscar. Quizás los animales sentían la inminencia, la extraña vibración en el aire de esos últimos días que los humanos apenas percibían.

Julio se levantó de la mesa. No podía quedarse sentado, inactivo. La responsabilidad que pesaba sobre él era inmensa. Había decidido pasar esas últimas horas revisando, una y otra vez, todos los sistemas, cada protocolo. Quería asegurarse de que absolutamente nada fallaría.

Se dirigió al taller, su mente ya inmersa en la vorágine de datos y verificaciones. Activó el dron que había fabricado, una maravilla tecnológica propia, pequeño y casi imperceptible en el aire. Lo envió a verificar los puntos de enterramiento de los cubos alrededor del perímetro del manzano. En su pantalla de control, cada señal respondía con una luz verde confirmando la integridad de los dispositivos. La verificación fue perfecta. El núcleo en el sótano mostraba una estabilidad impecable, su zumbido apenas audible. La máquina **Pi** estaba activa, en modo de espera, lista para el momento crucial. **Pi** misma había revisado toda la programación de la alerta sísmica y le confirmó a Julio mentalmente que el mensaje se emitiría sin problemas, alcanzando cada dispositivo y cada hogar.

Faltaban **6 horas, 02 minutos y 38 segundos**.

Julio se sentó frente a la consola de **Pi**, extendiendo una mano sobre su superficie fría y lisa. Fox saltó a su regazo, acurrucándose y ronroneando con una fuerza inusual. El joven cerró los ojos, respirando profundamente, tratando de aquietar el torbellino en su interior.

Una calma pesada cayó sobre él, una aceptación total de su destino y del destino de aquellos que dependían de él. No podía cambiar lo que vendría. Todo lo que había estado en su poder, lo había hecho. Solo podía esperar que todo saliera como estaba planeado. La certeza de que el **Arca** era su última y única esperanza llenó su mente.

Afuera, más allá de los muros de su casa, la vida en el vecindario discurría con una normalidad desesperante. Algunos vecinos cortaban el césped, el sonido de las cortadoras resonando monótonamente en el aire. Otros paseaban a sus perros, charlando con otros transeúntes sobre las trivialidades del día. Algunos niños, despreocupados y llenos de vitalidad, jugaban en la calle, sus risas inocentes perforando el silencio de Julio.

Julio los miraba desde la ventana de su taller, una punzada en el pecho. Ignoraban por completo que en apenas unas horas, todo cambiaría. Que el cielo que les ofrecía esa mañana gris se convertiría en una furia, y que sus vidas normales, sus rutinas, sus pequeñas preocupaciones diarias, se extinguirían en un instante.

La cuenta regresiva en su reloj parpadeó.

**6 horas, 00 minutos, 00 segundos**.

De pronto, un sonido estridente y discordante rompió la quietud de la mañana. Los celulares comenzaron a vibrar con una furia inusitada, sus tonos de alerta irrumpiendo en cada hogar. Las radios, que hasta ese momento emitían música suave, se interrumpieron abruptamente con una voz robótica y alarmante. Los televisores de los vecinos, que podían verse a través de las ventanas, emitieron la misma alerta estridente, un pitido agudo seguido de un mensaje que se repetía sin cesar:

\begin{quote}
    "**¡Alerta Nacional! Sismo de categoría 8. Refúgiense en sus casas inmediatamente. Cierre puertas y ventanas. Manténgase resguardado. Posible impacto en 6 horas.**"
\end{quote}

Los vecinos se miraron, confundidos. Algunos se reían nerviosos, murmurando sobre un error, una broma de mal gusto. Otros, más precavidos, o simplemente obedientes a la autoridad, comenzaron a reaccionar con una velocidad sorprendente. Rápidamente, el movimiento en las calles se detuvo. Puertas se cerraban con golpes secos. Ventanas se aseguraban. Se veían figuras correr hacia sus hogares, siguiendo la instrucción de "refugiarse bajo techo" como si fuera un ensayo general.

Óscar sonrió, viendo cómo el plan de **Pi** funcionaba con una eficiencia aterradora. El simulacro de sismo era la coartada perfecta, una forma de meter a todos en casa sin pánico ni preguntas innecesarias.

Julio soltó un suspiro de alivio que no sabía que contenía. Era ahora o nunca. El momento decisivo se acercaba.

Se dirigió hacia su familia, que ya comenzaba a prepararse con una disciplina ensayada, cada uno cumpliendo su rol asignado. Miguel y Carlos se aseguraron de que las últimas herramientas estuvieran guardadas y las mochilas personales listas. Ana y Elena revisaron los kits de emergencia. Óscar y Sebastián cargaron las mochilas personales.

Juntos, bajaron al sótano, donde el **núcleo del Arca** esperaba, vibrando levemente en la oscuridad, como un corazón latiendo en anticipación. La presencia de Don Ramiro, con su calma serena y sus ojos ahora abiertos, era un ancla en la creciente tensión.

La cuenta regresiva en el dispositivo de Julio seguía implacable:

**0 días, 5 horas, 57 minutos, 10 segundos.**

La tierra temblaría pronto.

Y ellos, la familia, Don Ramiro, y con suerte, las **150 personas** de su vecindario, estarían listos.

---

**0 días, 1 hora, 12 minutos, 43 segundos.**

Una vibración leve, casi imperceptible, recorrió el suelo bajo sus pies. Fue tan sutil al principio que pensaron que era solo el eco de su propia imaginación, de sus nervios a flor de piel.

Pero luego se repitió, un poco más intensa esta vez, una pulsación rítmica que se elevaba desde las profundidades de la tierra.

La familia entera se puso en pie, sus miradas buscando a Julio, sus ojos fijos en la pantalla de **Pi**. Julio corrió hacia el monitor, sus propios ojos fijos en las lecturas que comenzaban a dispararse, las líneas en los gráficos ascendiendo con una velocidad alarmante.

\begin{quote}
    —¡Pi! —llamó Julio, su voz tensa. —Activa el dron exterior. ¡Quiero ver!
\end{quote}

En el momento en que **Pi** le concedió acceso, la pantalla de Julio cambió. La vista aérea del dron se desplegó, mostrando el vecindario silencioso bajo el cielo crepuscular. Pero el cielo… el cielo ya no era el mismo.

Arriba, el firmamento se teñía de una paleta de colores imposibles. La radiación de la **supernova** había alcanzado la atmósfera. Grandes manchas de rojo profundo, casi carmesí, se mezclaban con velos de verde eléctrico y destellos de azul violeta, como una aurora boreal enloquecida, pero infinitamente más vasta y amenazante. Remolinos de luz brillante, a veces amarillenta, a veces de un blanco cegador, se expandían como velos fantasmales, un espectáculo hermoso y aterrador a la vez, una danza cósmica de la muerte que los ojos humanos no estaban hechos para presenciar.

Y luego, la tierra.

A medida que el dron ascendía, Julio pudo ver cómo en el suelo, siguiendo los patrones ocultos de los **cubos enterrados**, comenzaban a aparecer grietas. No eran fisuras aleatorias de un sismo; eran líneas limpias, casi quirúrgicas, que se extendían como arterias geológicas. Estas grietas se abrían con lentitud al principio, un susurro en la superficie de la tierra, pero luego se ensanchaban, revelando una luz azul pulsante que emergía desde las profundidades. Las líneas se conectaban de cubo a cubo, dibujando un perímetro perfecto alrededor de toda la manzana, como un diagrama inmenso dibujado en el suelo mismo.

**0 días, 0 horas, 43 minutos, 12 segundos.**

El **núcleo del Arca**, allí en el centro del sótano, emitió un pulso azul profundo, una luz etérea que llenó el espacio por un instante antes de desvanecerse. Era la señal. La activación había comenzado.

Uno a uno, los **cubos enterrados** bajo la manzana, los mismos que habían costado semanas de sudor y agotamiento, se activaron con una potencia creciente. Julio veía las lecturas en la pantalla con el corazón en la boca, sus dedos apretados en un puño. Cada cubo se iluminaba bajo tierra, sus firmas energéticas apareciendo en el monitor de **Pi**, formando una red viva que envolvía toda la manzana como un manto protector invisible.

Mientras tanto, la vista del dron afuera era indescriptible. De las grietas iluminadas, el suelo comenzó a ondularse suavemente, no como un terremoto, sino como la superficie de un líquido denso. Y entonces, de esas líneas brillantes y ondulantes, un **domo invisible** comenzó a emerger. Era un velo iridiscente, una burbuja de energía que se levantaba del suelo con una majestuosidad lenta e imparable. Al principio, era solo una leve distorsión en el aire, como el calor que se eleva del asfalto. Pero a medida que ascendía, se hizo visible por la forma en que refractaba la luz extraña del cielo, tiñéndose con los mismos colores de la supernova: rojos profundos, verdes esmeralda, y azules eléctricos.

Este domo no solo se elevaba más de **200 metros** hacia arriba, abarcando toda la manzana, sino que también se extendía hacia abajo, descendiendo a profundidades insondables, envolviendo la totalidad del **Arca** que se construía debajo. Era una esfera gigantesca, una cáscara protectora e impenetrable que cubría el vecindario como un huevo invisible, sellándolos del exterior, aislándolos de la muerte cósmica que se acercaba.

Óscar apretó el hombro de su hermano, sus ojos fijos en la pantalla del dron, boquiabierto ante la magnitud de lo que estaban presenciando. Miguel y Ana se miraban, lágrimas silenciosas rodando por sus mejillas, no de tristeza, sino de una profunda gratitud y alivio. Carlos y Elena se tomaron de la mano, y Sebastián se aferró a la pierna de su padre. Don Ramiro asintió, sus ojos ahora abiertos, observando el núcleo con una expresión de profunda comprensión y una pizca de asombro.

\begin{quote}
    —Sea lo que sea que pase ahora, hijo —dijo Miguel, su voz ronca por la emoción, con el nudo en la garganta—, ya hiciste más de lo que nadie, jamás, podría haber hecho. Salvaste… nos salvaste a todos.
\end{quote}

Julio apretó los dientes, los músculos de su mandíbula tensos. No quería héroes. No quería aplausos. Solo quería salvarlos. Quería que funcionara.

**0 días, 0 horas, 11 minutos, 30 segundos.**

El núcleo comenzó a vibrar con una intensidad creciente, un rugido grave que llenaba el sótano. Las paredes de concreto retumbaban, y el polvo fino se desprendía del techo. El temblor se intensificó, un crescendo que anunciaba la inminencia del cataclismo.

Fox saltó al suelo, sus orejas hacia atrás, y arqueó el lomo, maullando de forma grave, un sonido ronco que expresaba su propia alarma y temor.

Julio gritó por encima del estruendo creciente: —¡Todos al círculo central! ¡Ahora!

Se colocaron rápidamente alrededor del núcleo, siguiendo las líneas brillantes que **Pi** había proyectado en el piso, conectándolos entre sí. Era el punto focal, la zona de máxima protección durante la transformación.

La máquina emitió un sonido grave y profundo, como un latido descomunal, una pulsación que resonaba en sus propios huesos y se sincronizaba con el pulso del **Arca**.

**0 días, 0 horas, 1 minuto, 5 segundos.**

Julio cerró los ojos, sus labios moviéndose en un silencioso mantra. —Vamos... vamos... —Era una súplica al universo, al **Arca**, a todo lo que había puesto en juego.

**0 días, 0 horas, 0 minutos, 30 segundos.**

El suelo vibraba como un tambor frenético, una percusión ensordecedora. Arriba, las casas crujían, los vidrios se rompían con el sonido amortiguado de la lejanía. El aire se llenó con el sonido de la tierra que se desgarraba bajo el creciente **domo invisible**.

Y entonces, cuando el contador llegó a **0…**

Todo explotó en luz.

Pero no hubo calor devastador. No hubo destrucción ardiente. Solo una expansión de energía pura, de un blanco cegador y un azul profundo, que los envolvió completamente, como una ola gigantesca, silenciosa e imparable que se elevaba desde las profundidades de la tierra. La luz fue tan intensa que a través de la pantalla del dron, por un instante, el domo se volvió completamente opaco, un orbe reluciente que consumió el vecindario.

Julio abrió los ojos a tiempo para ver cómo las paredes del sótano, que los habían protegido durante tanto tiempo, parecían doblarse y disolverse, no desmoronarse, sino transformarse, volverse transparentes como si fueran de gelatina, y luego desvanecerse en la luz. En su lugar, ya no era un sótano, sino una inmensa estructura subterránea. Pasillos amplios y abovedados se extendían en la distancia, cúpulas luminosas se alzaban sobre ellos, y cámaras espaciosas se abrían en todas direcciones. El **Arca**, en su plenitud.

Había comenzado.

\cleardoublepage % Para que el capítulo empiece en una página impar

% Trigésimo primer capítulo
\chapter*{31 – El Último Barrio del Universo}
\addcontentsline{toc}{chapter}{31 – El Último Barrio del Universo} % Añadir al índice manualmente

La cuenta regresiva se desvaneció en el silencio más absoluto: **0 días, 0 horas, 0 minutos, 00 segundos**.

El latido profundo y resonante de la máquina, el último eco de su titánica transformación, se desvaneció lentamente en el aire. Por un instante, todo fue silencio. Un silencio tan denso y profundo que dolía en los oídos, un vacío que parecía absorber cualquier sonido, cualquier pensamiento, incluso el de la propia existencia. Era la pausa entre dos realidades.

Julio fue el primero en abrir los ojos. La luz que lo rodeaba era diferente, suave, difusa, emanando de las paredes mismas. Su visión se ajustó lentamente a la nueva realidad. El sótano, con sus muros de concreto y sus herramientas esparcidas, no existía más. Ahora estaban dentro de una estructura inmensa, vasta más allá de lo que su mente podía haber concebido en los planos: paredes curvas y etéreas, de un blanco translúcido, iluminadas suavemente desde dentro por una fuente invisible. El **núcleo del Arca**, el corazón que habían alimentado con sus esperanzas y esfuerzos, latía tranquilo y brillante en el centro del espacio, un faro de estabilidad en la inmensidad.

A su alrededor, su familia —Miguel, Ana, Óscar— y la familia de Sebastián —Carlos, Elena, Sebastián—, junto al inmutable Don Ramiro, estaban de pie, o recuperándose de la conmoción, sus expresiones una mezcla de asombro y confusión. Fox, había saltado del regazo de Julio y miraba a su alrededor con sus ojos celestes, sus pupilas dilatadas, su cola erizada en una mezcla de curiosidad y cautela. **Pi**, la interfaz de luz azul, flotaba serenamente cerca de Julio, su presencia una constante en el caos de la revelación.

Sin decir una palabra, movidos por una fuerza instintiva, avanzaron hacia la única salida visible: una gran compuerta, tan lisa y perfecta que parecía parte de la pared misma, que se abrió automáticamente al acercarse, deslizándose sin un sonido.

Una brisa artificial, limpia y fresca, que no olía a tierra ni a humedad, sino a aire filtrado y vida, acarició sus rostros. Era un contraste sorprendente con el denso aire del sótano, una bocanada de lo que sería su nueva atmósfera.

Subieron por una rampa suave, que ascendía en una curva elegante, sus pasos resonando apenas en el material del **Arca**. Y al salir al "exterior", la realidad los golpeó con la fuerza de mil tempestades, una revelación que detuvo la respiración y paralizó sus mentes.

Ya no había casas a su alrededor.

Ya no había calles, ni aceras.

Ya no había árboles ni jardines que conocieran, ni el familiar horizonte de las montañas que los había rodeado toda su vida.

No había cielo azul ni nubes algodonosas.

No había nada.

Solo negro.

Un negro inmenso, absoluto, más profundo que cualquier noche estrellada que hubieran visto. Un negro puntuado por millones de estrellas que titilaban con una luz gélida y distante, más brillantes, más numerosas, más indiferentes que nunca, como diamantes dispersos en un paño de terciopelo infinito. Era el vacío implacable del espacio.

Y en medio de ese infinito, flotando como una isla imposible, milagrosamente intacta, estaba su manzana. Su barrio. Los jardines con sus flores aún vivas, las casas de sus vecinos con sus techos y ventanas familiares, el pequeño parque central con sus bancos y árboles… todo, absolutamente todo, estaba envuelto en una cúpula transparente. Era apenas perceptible, un domo gigantesco que se extendía cientos de metros hacia arriba y hacia abajo, envolviendo la manzana en su totalidad como un huevo de cristal invisible. Solo se notaba por el leve reflejo de las luces estelares, que se refractaban en su superficie como gotas de rocío cósmico.

Julio cayó de rodillas, el aire escapando de sus pulmones en un suspiro, boquiabierto, sus ojos fijos en la inmensidad aterradora y sublime. La vista era tan sobrecogedora que desafiaba toda lógica, toda comprensión.

Óscar, a su lado, simplemente susurró, su voz apenas audible, cargada de una incredulidad asombrosa:
\begin{quote}
    —No… no puede ser…
\end{quote}

Miguel ayudó a Ana a mantenerse de pie, su rostro ceniciento. Ella lloraba en silencio, cubriéndose la boca con una mano, sus hombros temblaban. Carlos sostenía a Elena, quien miraba con los ojos desorbitados, y Sebastián, con la boca abierta, parecía no entender la magnitud de lo que sus pequeños ojos estaban procesando. Don Ramiro, por primera vez, parecía genuinamente sobrecogido, aunque su expresión de asombro era más de maravilla que de terror.

Y entonces, uno a uno, las puertas de las casas del vecindario que los rodeaban en esa burbuja, puertas que habían permanecido cerradas por la falsa alarma de sismo, comenzaron a abrirse. Los vecinos salían, confusos, somnolientos, sus rostros marcados por la recién interrumpida rutina de refugio. Miraban alrededor, sus ojos escudriñando el familiar pero extrañamente iluminado paisaje, buscando una explicación a la ausencia de la calle, del cielo, del mundo.

Fox, con un maullido desconcertado que sonó más a quejido, se escondió tras las piernas de Julio, sintiendo la extrañeza del nuevo "exterior".

El tiempo pareció detenerse mientras las decenas de vecinos emergían de sus hogares, sus miradas perplejas recorriendo la escena. Sus ojos se detuvieron en la familia de Julio, en Don Ramiro, en Carlos y Elena, los únicos que parecían estar en calma, o al menos, en una calma resignada. Las preguntas comenzaban a formarse en sus rostros, en sus bocas.

Fue entonces que Miguel, su padre, puso una mano firme y pesada sobre el hombro de su hijo. Era un gesto de apoyo, pero también de implacable responsabilidad.
\begin{quote}
    —Tienes que explicarles, Julio —dijo, con una voz grave y resonante, una voz que no dejaba lugar a dudas—. Con la cabeza en alto. Sin miedo. Confían en ti.
\end{quote}

Julio tragó saliva, el nudo en su garganta más apretado que nunca. Sabía que tenía que hacerlo. Lo había sabido desde el día en que **Pi** le mostró por primera vez la verdad. Esta era su misión, su destino.

Se puso en pie lentamente, el peso de todas esas miradas expectantes sobre él. El vacío del espacio parecía tragárselo todo, pero la cúpula invisible, la que él había ayudado a anclar, protegía el aire que respiraban, el calor que los envolvía, la vida misma en esa pequeña isla de existencia. Y ahora, todos los rostros del vecindario, desde el tendero hasta la anciana del huerto, lo miraban. Esperando.

Julio respiró hondo, un aliento que llenó sus pulmones de aire filtrado del **Arca**, no del aire de un mundo que ya no existía. Y comenzó a hablar.

Su voz, al principio un poco temblorosa, se fue haciendo más firme a medida que las palabras fluían. Explicó, con la mayor calma que pudo reunir, cómo el evento cósmico había arrasado no solo la vida, sino la misma existencia del planeta. Les contó cómo la **supernova**, una estrella que había expirado en una explosión de inimaginable violencia, había lanzado una ola de radiación letal que había vaporizado la superficie de su mundo natal. Les habló de la alerta de sismo de **Pi**, de la necesidad de que estuvieran en sus casas, y de cómo el manzano había sido salvado por algo que aún escapaba a su plena comprensión, una tecnología que "**ellos**" le habían dado, convertida en el **Arca**. Les explicó la cúpula, la autosuficiencia, el plan para sobrevivir.

Mientras hablaba, mientras explicaba el fin de su mundo y el milagro de su salvación, la gente no gritaba ni lloraba en pánico colectivo. Escuchaban en silencio. Atónitos. Sus ojos pasaban de Julio a la visión imposible del espacio, de la familiaridad de sus propias casas dentro de la burbuja a la oscuridad infinita que los rodeaba. Porque sabían, en lo más profundo de su ser, que lo que veían con sus propios ojos no podía ser un sueño, ni una mentira elaborada. El aire era el mismo, sus casas eran las mismas, sus vecinos eran los mismos. Pero el entorno era completamente diferente.

El mundo ya no existía.

Solo ellos quedaban. Su pequeña manzana, suspendida en el vasto, inimaginable océano del cosmos.

Y así, bajo un manto infinito de estrellas, en medio de la nada, en esa burbuja de vida y esperanza, nació la última comunidad del universo. Una comunidad unida no solo por la proximidad física, sino por la devastadora verdad que ahora compartían: eran los únicos supervivientes. El verdadero desafío, el de construir una nueva vida a partir de las cenizas de un mundo perdido, apenas comenzaba.

---

\chapter*{Capítulo Final: La Nueva Sociedad}
\addcontentsline{toc}{chapter}{Capítulo Final: La Nueva Sociedad} % Añadir al índice manualmente

Julio terminó de hablar.

Su voz, que al principio había temblado con la magnitud de la verdad que cargaba, había ganado fuerza y una resonancia inquebrantable a medida que cada palabra salía de sus labios. No ocultó nada. No adornó los hechos con eufemismos o falsas esperanzas. Con la brutal honestidad que la situación exigía, les dijo lo que en su interior ya sabían, pero que sus mentes se resistían a aceptar:

\begin{quote}
    —El planeta… nuestro hogar, la Tierra… ha desaparecido. La onda de radiación de la supernova lo alcanzó. Solo nosotros hemos sobrevivido aquí, dentro de este **Arca**. Y no fue por casualidad.
\end{quote}

Por un instante eterno, nadie habló. El silencio se hizo más profundo que el vacío espacial que los rodeaba, tan denso que dolía en los oídos. Solo el sonido de sus propias respiraciones, entrecortadas y temblorosas, llenaba el espacio debajo de la cúpula protectora, un coro de humanidad al borde del abismo.

Las decenas de vecinos del manzano, que segundos antes habían emergido de sus casas con la perplejidad de quien acaba de despertar de un sueño profundo, ahora estaban paralizados. Sus rostros reflejaban una miríada de emociones: shock, negación, un terror silencioso que asomaba en sus ojos desorbitados, y una abrumadora incredulidad. Algunos, más allá de toda contención, cayeron de rodillas, el mundo familiar bajo sus pies ahora suspendido en la nada. Otros, los más mayores, se llevaban las manos al pecho, como si pudieran sentir el dolor de un planeta perdido.

Hasta que alguien, una vecina de cabello canoso que cultivaba hortalizas en el parque central, Doña Marta, conocida por su resiliencia y su mirada directa, dio un paso adelante. Sus ojos, aunque húmedos, brillaban con una lucidez inesperada.

\begin{quote}
    —Entonces… —dijo, su voz rasposa pero clara—, ¿estamos vivos porque tú nos salvaste, Julio?
\end{quote}

Julio negó con la cabeza, su mirada humilde, aunque su postura era ahora la de un líder natural.

\begin{quote}
    —No, Doña Marta. No. Estamos vivos porque todos juntos formábamos la única oportunidad de salvar algo. Yo solo… yo solo construí lo que ‘ellos’ me enseñaron. Ellos me dieron las herramientas, la guía. Ustedes son los que hacían que este barrio fuera la opción. Sin cada uno de ustedes, con sus habilidades, su presencia… nada de esto habría sido posible. Esta **Arca**… es de todos. Es nuestra.
\end{quote}

Un murmullo, un suspiro colectivo, recorrió a los vecinos. Era una mezcla de alivio, de pena, de una profunda perplejidad. Algunos, liberados de la tensión, lloraban ahora abiertamente, sus sollozos rompiendo la tensión del aire. Otros se abrazaban, buscando consuelo en la cercanía de sus semejantes. Algunos simplemente se limitaban a mirar hacia el espacio infinito que los rodeaba, buscando en vano las formas familiares de la Tierra, las estrellas que conocían, las constelaciones que ya no existían, tragadas por la oscuridad. La realidad era incomprensible, pero innegable.

Óscar se acercó a su hermano, su rostro ahora serio, la burla juvenil desvanecida por la gravedad de la situación. Le susurró al oído, su voz urgente:

\begin{quote}
    —Tendremos que organizarnos, Julio. Rápidamente. Cultivar lo que podamos con las semillas que trajimos, asegurar el agua, proteger los sistemas… aprender a vivir aquí. Esta no es nuestra casa habitual, es nuestra nueva casa, pero también es una nave.
\end{quote}

Miguel asintió con gravedad, su mano posada en el hombro de Óscar, un gesto de reconocimiento por la madurez que su hijo exhibía.

\begin{quote}
    —No podemos perder tiempo —dijo Miguel, su voz ahora la de un patriarca, fuerte y clara para que todos la escucharan—. Somos la última chispa. La última llama en la oscuridad. Y cada uno de nosotros tiene un papel que cumplir.
\end{quote}

Ana miró a Julio con ojos llenos de un orgullo inmenso, mezclado con la tristeza de una madre que veía a su hijo cargar un peso tan monumental. Carlos y Elena, que ya habían consolado a Sebastián, ahora se acercaron, sus rostros serios pero determinados. Don Ramiro, se había puesto de pie y observaba la multitud con una mirada que abarcaba la comprensión y la esperanza.

\begin{quote}
    —Julio tiene razón —dijo Don Ramiro, su voz pausada y calmada cortando el murmullo—. La providencia nos trajo aquí. Ahora, es nuestra tarea. Esta es la comunidad.
\end{quote}

Fox, como siempre, permanecía cerca de Julio, su cola envolviendo suavemente el tobillo del joven, mirándolo con sus ojos grandes y brillantes, como si ya supiera algo más. Algo que los demás aún no podían entender. Quizás, que la vida, en cualquier forma, encontraría su camino.

Julio respiró hondo, sintiendo el peso del futuro sobre sus hombros, más pesado que cualquier cosa que hubiera imaginado en sus cálculos y planos. Pero también sintió algo más: la certeza de que no estaba solo. Que habían sido elegidos, sí, pero ahora dependería de ellos construir algo nuevo, de convertir esa "última chispa" en una hoguera que iluminara el vasto vacío.

Cuando el primer shock comenzó a disiparse, y todos los vecinos, liderados por sus familias, comenzaron a moverse con un propósito renovado, las primeras decisiones prácticas empezaron a tomarse. Era un coro de voces, cada una aportando su conocimiento: el electricista, el mecánico, la enfermera, el tendero, el maestro, la cocinera, los niños… La gente que había vivido de forma anónima uno al lado del otro ahora se veía como los únicos supervivientes, cada habilidad una herramienta invaluable para su nueva existencia.

Miguel y Carlos, con su experiencia en ingeniería, se encargaron de organizar los equipos para revisar los recursos y asegurar que los sistemas de soporte vital del **Arca** funcionaran correctamente. Ana y Elena, con la ayuda de la vecina enfermera y la cocinera, empezaron a establecer los turnos de cuidado para los más pequeños y los mayores, y a organizar el racionamiento inicial de los alimentos. Óscar, con Sebastián y los otros niños, se encargaron de establecer los puntos de control de los juegos y la biblioteca, una forma de mantener la moral alta y la mente activa en los largos días venideros. Doña Marta, con su profundo conocimiento de la tierra y las plantas, se unió a Julio para planificar los primeros cultivos hidropónicos, su rostro iluminado por la visión de vegetales creciendo bajo el suelo de un mundo que ya no existía.

Julio se apartó un momento del bullicio de la organización, caminando hacia el límite de la cúpula, donde la transparencia revelaba el universo. Apoyó su mano contra la barrera invisible y suavemente cálida. Más allá, solo el vacío absoluto. Y las estrellas, millones de ellas, indiferentes y eternas, que brillaban con una intensidad deslumbrante.

Entonces, sintió a **Pi** zumbando suavemente junto a él, su luz azul titilando en la semioscuridad del domo.

\begin{quote}
    —¿Pi…? —murmuró Julio, su voz apenas un susurro. La pregunta no era si Pi estaba ahí, sino si su propósito había terminado. — ¿Aún estás aquí?

    —Siempre —respondió Pi, en su tono sereno, su voz resonando directamente en la mente de Julio—. Mi propósito es la continuidad. Y no estás solo, Julio. No aún. Mi existencia es la de tus creadores. Y ellos… ellos también te observan.
\end{quote}

Julio sonrió levemente, una sonrisa teñida de asombro y una extraña esperanza. Miró hacia las estrellas, un lienzo infinito de posibilidades y misterios. Sabía que en algún lugar allá afuera, entre esos millones de puntos de luz, "**ellos**" seguían observando. Ellos, que le habían dado la oportunidad, que le habían confiado la semilla de un nuevo comienzo.

Ellos que, tal vez, algún día, volvería a encontrar.

La historia de la Tierra había terminado. Pero la historia del **Arca**, la de la última comunidad del universo, apenas comenzaba. En esa pequeña burbuja de vida flotando en la nada, una nueva sociedad comenzaba a forjarse, forzada por la catástrofe, pero unida por la esperanza y la inquebrantable voluntad de sobrevivir.

---
\begin{center}
    \textbf{FIN DEL LIBRO I: El Último Barrio del Universo}
\end{center}
---



\end{document}